% Supplement to "An algebra of mechanism composition".
% Sixty protocol profiles. Generated; see formal/v3/emit-tex.mjs.
\documentclass[11pt]{amsart}
\usepackage[margin=1.15in]{geometry}
\usepackage{amsmath,amssymb,amsthm,mathtools}
\usepackage[hidelinks]{hyperref}
\theoremstyle{remark}
\newtheorem{measurement}{Measurement}[section]
\newtheorem{remark}[measurement]{Remark}
\newcommand{\El}{\mathcal{E}}
\title{Sixty protocol profiles}
\date{}
\begin{document}
\maketitle

This supplement records, for each of the sixty protocols, the construction
exhibited for it, its canonical form, the verdict of the admissibility test,
the obligations it discharges and the obligations no element of the
vocabulary names. Every figure is generated from the same verification that
produced the measurements in the article.

\subsection{Curve}\label{sub:cat:dex:curve}

A depositor supplies any subset of a pool's coins, balanced or not, and receives a fungible claim on the whole reserve, exiting either balanced or into a single coin and paying an imbalance fee calibrated so that depositing one coin and withdrawing another costs about what the corresponding swap would cost. Stableswap pools quote off a hybrid of the constant-sum and constant-product invariants whose leverage term degrades toward constant product as the pool goes out of balance, concentrating liquidity at a fixed parity. Cryptoswap pools apply the same equation under a different leverage term around a moving internal price which the pool re-centres itself, consulting no external oracle and requiring no action from the provider. Settlement is atomic and per-pool, and a metapool composes by taking another pool's claim token as one of its own coins. A position ends on withdrawal --- the exchange carries no debt and so has no liquidation --- while a fixed share of every trading fee leaves the pool as an admin fee, is converted through another protocol's batch auction, and is distributed to lockers of the governance token.

\begin{measurement}\label{meas:cat:dex:curve}
Take $X = \{Ag,Em,Fd,Gp,Sh,St,Tg,Tp\}$, with canonical form
$\mathrm{ex}(X) = \{Ag,Em,Fd,Gp,Sh,St,Tg,Tp\}$, every element being primitive in it.
It is admissible: no requirement term is open, every element is warranted, and it arms no prohibition. It discharges 8 of the 22 recorded obligations; 14 are residue. No composite of at most three corpus protocols equals it; 3 corpus protocols are contained in it, and together they supply every element except \{Gp\}.
\end{measurement}

\begin{remark}
Stableswap and Cryptoswap are one equation under different leverage terms, and the vocabulary names the one and not the other, so the residue here is not a missing mechanism but a cut through the middle of a parameterised family. What the cut hides is the gate: the pool moves its own centre of concentration only when an exponential moving average of its own price has passed a step and the drawdown in a named profit functional of the pool state stays inside a bounded fraction of the profit that functional has accumulated --- a self-referential price, a monotone profit ledger and a ratio-bounded transition, none of them named, and a construction modelling the gate as repeg-if-profitable admits behaviour the protocol forbids. That safety property has a deadlock documented by the protocol itself, with a remedy list: a pool which has drifted too far earns too little to afford the move and strands its liquidity, and one published escape is a deliberately loss-making trade by a third party. The gate also cannot be pinned to an artefact, because both live implementations of one pool type publish the same version string and gate the repeg on mathematically different predicates --- the deployed one linear in accumulated profit, the repository's quadratic, the former the first-order expansion of the latter, agreeing while the profit ledger is young and diverging as profit compounds. A citation of that version without a bytecode hash names an ambiguous artefact, and that is a finding about decomposing deployed systems rather than a fact about this protocol.

The obligations no element discharges are these.
\begin{itemize}\setlength{\itemsep}{0pt}
  \item a Cryptoswap pool quotes off the IDENTICAL algebraic form with a different leverage term — K*D\textasciicircum{}(N-1)*Sum(xi) + Prod(xi) = K*D\textasciicircum{}N + (D/N)\textasciicircum{}N with K0 = Prod(xi)*N\textasciicircum{}N/D\textasciicircum{}N and K = A*K0*gamma\textasciicircum{}2/(gamma+1-K0)\textasciicircum{}2 — solved by Newton's method first for D then for xj at about 35k gas, with balances first transformed by a price vector so the invariant is always evaluated near equilibrium: liquidity is concentrated around a MOVING internal price\_scale rather than around parity
  \item the pool moves its own centre of concentration only when TWO conditions hold: an exponential moving average of the pool's own price has moved past the adjustment\_step, AND the loss of value of X\_cp = (Prod D/(N*pi))\textasciicircum{}(1/N) 'doesn't exceed half the profit we've made', tracked as a monotone accumulator (xcp\_profit) of that same functional
  \item the protocol runs several mutually incompatible implementations of that repeg gate in production at once, and they cannot be distinguished by the version string the contracts publish: the deployed and the repository implementations both report version 'v2.1.0' while testing mathematically different predicates
  \item the safety property in F3 has a documented deadlock: when the market price moves far from the last rebalance price the available liquidity falls, volume and fees fall with it, and 'without enough profit from fees, the pool cannot afford to rebalance and follow the price, leaving its liquidity stranded' — for which the protocol's own documented remedy includes wash-trading the pool with flash loans 'at a loss with no guarantee of a lasting fix'
  \item the fee is a function of pool imbalance, not a constant: 'Cryptoswap and all new Stableswap pools feature dynamic fees that adjust to increase returns for LPs when their liquidity is in high demand', governed by mid\_fee, out\_fee and fee\_gamma
  \item which asset the fee is taken in is fixed by pool type: stableswap takes it in the OUTPUT coin, cryptoswap takes it in the pool's own LP token
  \item the invariant's own parameter moves on a RAMP rather than a step: A is changed by ramp\_A with a minimum ramp time of 86,400 seconds, so the curve is time-varying by construction and a swap quoted at t and the same swap quoted at t+dt are priced by different invariants
  \item settlement is atomic and per-pool — each pool holds its own coins and there is no netting layer — and a metapool composes by making one of its coins another pool's LP TOKEN, tradeable through to the base coins but not depositable in them because of the contract size limit
  \item anyone may deploy a pool, but only the DAO may register the implementation it is deployed FROM: per-generation factories (CurveStableSwapFactoryNG, TwocryptoFactory, CurveTricryptoFactory, CurveL2TricryptoFactory) take the full parameter vector under asserted bounds (0 < adjustment\_step < 1e18+1, allowed\_extra\_profit < 1e18+1, fee\_max = 10*1e9, 86 < ma\_exp\_time < 872542), only set\_pool\_implementation is DAO-gated, and 'Pools created through the Factory are owned by the factory admin (DAO)'
  \item WHERE those emissions go is decided by a vote: veCRV holders vote a weight vector over gauges in the GaugeController, and that vector parameterises the emission function
  \item voting power is a non-transferable, non-reversible time-decaying lock: VotingEscrow.vy (Vyper 0.2.4, 0x5f3b5DfEb7B28CDbD7FAba78963EE202a494e2A2) with MAXTIME = 4*365*86400, a one-week minimum, unlock times rounded down to whole weeks, one lock per user, and power stored as a (bias, slope) piecewise-linear function with scheduled slope\_changes so decay needs no user action
  \item the lock also multiplies an LP's own emission rate, RIVALROUSLY: LiquidityGaugeV6.\_update\_liquidity\_limit sets working\_balance = min(l, l*40/100 + L*voting\_balance/voting\_total*60/100) and the reward rate is working\_balance/working\_supply, so an unboosted LP earns 40\% of their nominal share, a fully boosted LP earns 100\%, and one LP's boost lowers every other LP's yield through the shared denominator
  \item the protocol does not price its own fee stream: admin fees arriving in arbitrary tokens are converted by posting CoW Protocol CONDITIONAL ORDERS through CowSwapBurner (Ethereum 0xC0fC3dDfec95ca45A0D2393F518D3EA1ccF44f8b, Gnosis 0x566b9F24200A9B51b76792D4e81B569AF27eda83), and on chains without CoW the superseded route still runs — burn to MIM, bridge to mainnet, burn MIM to 3CRV
  \item the reference path for exercising that governance runs through another protocol: Curve's own voting library creates DAO votes through the Convex voter proxy 0x989AEB4D175E16225E39E87D0D97A3360524AD80
\end{itemize}
\end{remark}

\begin{remark}
$Gp$ is added: the corpus records no emergency authority here and there is one, a standing multisig whose powers are directional --- it may reduce a debt ceiling and never raise one, and may move fees without causing a liquidation --- so the witness names the body and leaves the directionality as residue. $Ve$ is dropped rather than carried, and the obligations show why: the lock is at once the vote, the weight vector that directs emissions, and a rivalrous multiplier on the holder's own emission rate through a shared denominator, and one contested symbol standing for all of that would assert an identity the gauge code refutes. $Fd$ is carried and means what it says here, distribution to a claim class, which makes this the only member of the category whose value return pays a claim class. The pool arithmetic is reconstructed from the published whitepapers rather than from the repositories, which are all-rights-reserved, released for information only, and grant no right of reproduction --- a constraint on the paper, not a gap in the record.
\end{remark}

\subsection{Fluid}\label{sub:cat:dex:fluid}

A lender deposits a single asset into a shared liquidity layer and receives a fungible vault share; a borrower opens a position represented by a transferable token, posting collateral and drawing debt against it. The two roles meet inside the exchange, because the collateral may itself be the exchange's inventory and so may the debt --- a swapper trades against borrowed principal, and the trading fees pay that principal down. Price is read off the constant-product curve contracted into a governance-set range about a centre which the pool shifts itself when a bound is breached, over a configured interval and within governance-set limits on how far the centre may travel; the deployed generation adds ticks and per-range positions in a singleton whose pool key names a controller. Settlement is atomic against the shared layer rather than against per-pool inventory. A position ends on repayment and withdrawal, or by liquidation --- and liquidation seizes no position: positions are bucketed into slots keyed by their ratio, a slot is closed whole and only far enough to restore health, and the closer is an ordinary trader taking a discount inside an aggregator route without knowing that is what they are.

\begin{measurement}\label{meas:cat:dex:fluid}
Take $X = \{Cl,Cp,Ct,Ex,Gp,Li,Pl,Sh,Tg,Tp,Up\}$, with canonical form
$\mathrm{ex}(X) = \{Cl,Cp,Ex,Gp,Li,Pl,Sh,Tg,Tp,Up\}$; the elements \{Ct\} are derived rather than chosen.
It is admissible: no requirement term is open, every element is warranted, and it arms no prohibition. It discharges 11 of the 22 recorded obligations; 11 are residue. No composite of at most three corpus protocols equals it; 1 corpus protocol is contained in it, and together they supply every element except \{Cp,Tp\}.
It is not minimal: \{Cp\} can be removed with every obligation still covered and admissibility intact.
\end{measurement}

\begin{remark}
The obligations this construction cannot state are the ones the protocol was built to have. One unit of capital is at once lent into the shared layer, borrowable against, and deployed as quoting inventory, and the inverse is stranger still: the inventory a swapper trades against is somebody's liability and the trade amortises it, so the vocabulary's partition of pool pricing from credit is not a simplification here but a mis-sort. Capacity is not a cap but a bound that relaxes on a schedule, with a slice held back for a liquidation that has not happened yet --- a reserved capability rather than a limit. Governance may deploy new pricing engines and not merely new pools, and which pools may exist is itself a governance decision, so the space of admissible curves is an object under live control rather than a fixed backdrop. $Li$ names third parties paid to close positions, and here the unit closed is an equivalence class rather than a position and the payment is execution quality delivered to someone who never asked to be a liquidator, so the symbol is carried on a mechanism that matches it in neither the unit nor the incentive.

The obligations no element discharges are these.
\begin{itemize}\setlength{\itemsep}{0pt}
  \item the pool moves its own centre without an owner action and without a profit test: the Center Price is 'the price at which the Pool rebalances itself once the Upper or Lower bounds are breached by the pool price', the move takes place over a Shift Time rather than instantaneously, and the centre may never leave [Center Min Price, Center Max Price]
  \item one unit of capital occupies three roles simultaneously: Smart Collateral 'is a single range order' that is lent into the Liquidity Layer earning supply yield, is borrowable against, AND is deployed as AMM inventory quoting swaps
  \item Smart Debt inverts it: the BORROWED PRINCIPAL is the AMM's inventory, so a swapper trades against someone else's liability and the trading fees pay that liability down — 'the first time debt can be transformed into a productive asset by using it as liquidity for the DEX', with borrowers 'even having the possibility of getting paid to borrow'
  \item capacity is a bound that RELAXES ON A SCHEDULE rather than a fixed cap: borrowing is governed by Base Limit, Current Limit, Max Limit, Expand Percentage ('the rate at which Limits would increase or decrease over the given duration') and Expand Duration, withdrawal by the same four plus Withdrawable, and the docs state the ceilings adjust 'in real-time based on utilization rates'
  \item a portion of the pool is reserved for an operation that has not happened yet: the Withdrawal Gap is a 'safety non-withdrawable amount to guarantee liquidations'
  \item the unit liquidated is an EQUIVALENCE CLASS, not a position: each position is assigned to a slot keyed by its ratio ('Drawing inspiration from Uniswap v3's slot-based liquidity'), and when a slot crosses the threshold 'the protocol aggregates all vaults in that slot for liquidation in a single transaction. Afterward, the system recalibrates the remaining vaults, adjusting their debt ratio positions down to the next slot' — taking only the minimum needed to restore health
  \item settlement is atomic on-chain against the shared Liquidity Layer rather than against per-pool inventory, and DEX v2 additionally advertises 'Flash Accounting (inspired by Uniswap v4)' alongside on-chain dynamic fees and on-chain limit orders that 'earn lending APR while waiting to be filled'
  \item the extension point is a CONTROLLER named per pool: the DexKey carries a controller address alongside tickSpacing, and the feature is documented as Hooks
  \item which pools may exist is decided by governance: 'DEX v2 operates with pools that are permissioned and listed on the Fluid Money Market ... Each pool has been approved through Fluid governance', discoverable through getD3PermissionedDexes() and getD4PermissionedDexes(), with a stated future transition to permissionless pools discovered from LogInitialize(dexType, dexId, dexKey, sqrtPriceX96)
  \item governance can create new PRICING ENGINES, not merely new pools: 'At its core, Fluid DEX v2 runs on a singleton contract built atop the Fluid Liquidity Layer ... Governance can deploy infinite DEX types, each with its own logic and math', with four types at launch (v1 Smart Collateral, v1 Smart Debt, Smart Collateral Range Orders, Smart Debt Range Orders)
  \item value return runs through a governance-toggled Revenue Cut into an on-chain buyback: 'Revenue Cuts can be turned on through governance voting and could consist of a cut from trading fees and/or a fee from any Smart Vault', and the deployed contract index carries Periphery/Buyback/\{FluidBuyback, FluidBuybackProxy, FluidBuybackCore\}
\end{itemize}
\end{remark}

\begin{remark}
The corpus record decomposes DEX v1 while DEX v2 is deployed, and one of its markers is wrong as recorded: it holds that there are no per-user ranges here to name, while the deployed singleton carries tick spacing in the pool key and a current tick in pool state and ships range orders on the collateral side and on the debt side both, so $Cl$ is carried rather than refused. $Tp$ is added, which the corpus omits --- the exchange publishes its own price series, and publishes maxima and minima over the window as well, which is more than the symbol names. The hook residue arrives here under its third naming convention in this category, a controller address sitting in the pool key beside tick spacing, so the extension point is now shared across the category's largest members and named by none of them. No value-return element is carried: buyback contracts stand in the deployed contract index and the revenue cut is a governance toggle, but whether the channel is running was not established at a primary source, and an element asserted on an unverified switch would be an element asserted from memory.
\end{remark}

\subsection{PancakeSwap}\label{sub:cat:dex:pancakeswap}

A provider deposits a pair into a constant-product pair carrying a fungible share, a tick-bounded position carrying a non-fungible one, a stable pool, or a Liquidity Book pool whose claim is fungible per price bin. Price formation follows the pool type: constant product on the reserves, the same curve restricted to a tick range, a stable-hybrid invariant near parity, and, inside a bin, constant sum --- so a trade that stays within one bin moves no price at all. The current generation separates custody from pricing, an immutable vault holding every balance and counting unsettled deltas in transient storage while independent pool managers for the concentrated and the bin families register against it, which makes the pricing engine pluggable at the contract level rather than only at the hook level. A position ends by burning the share or the position token and collecting; a provider may instead stake it into a farm and draw emissions on an administered allocation. There is no debt and so no liquidation, and the protocol's share of every fee is taken as a split declared at the protocol level with a token burn as one of its legs.

\begin{measurement}\label{meas:cat:dex:pancakeswap}
Take $X = \{Cl,Cp,Em,Fl,Gp,Ix,Sh,St\}$, with canonical form
$\mathrm{ex}(X) = \{Cl,Cp,Em,Fl,Gp,Ix,Sh,St\}$, every element being primitive in it.
It is admissible: no requirement term is open, every element is warranted, and it arms no prohibition. It discharges 9 of the 22 recorded obligations; 13 are residue. No composite of at most three corpus protocols equals it; no corpus protocol is contained in it.
It is not minimal: \{Cl\} can be removed with every obligation still covered and admissibility intact.
\end{measurement}

\begin{remark}
The pricing group offers a symbol for the constant-product curve, one for its restriction to a range, one for the stable hybrid, one for a weighted-geometric invariant and one for an oracle-priced inventory curve, and this protocol still deploys a live, documented, audited pool type that none of them names. Neither the weighted-geometric symbol nor the oracle-priced one discharges anything anywhere in this category, so the group's apparent resolution is redundancy where the pools are alike and silence where they differ. The hook residue recurs here with the permissions carried in the pool's own key rather than mined into the hook's address, which means one hook contract may run under different permission sets in different pools --- the same mechanism as Uniswap's under a binding that changes both the deployment model and who chooses the capability. Other obligations here concern time rather than state: a whole mechanism was retired and the holders of a third party's derivative of it were settled at par, and the terminal bound on token supply was itself lowered by a vote. The vocabulary records which mechanisms a protocol has; it cannot record that one was retired, that a liability to somebody else's derivative was assumed, or that a supply cap is governable.

The obligations no element discharges are these.
\begin{itemize}\setlength{\itemsep}{0pt}
  \item an Infinity LBAMM ('Liquidity Book') pool prices instead on DISCRETE BINS: liquidity sits in a bin at a fixed price, the invariant WITHIN a bin is constant sum — 'LBAMM follows the constant sum formula (X + Y = K)' — so a trade inside one bin has zero price impact, and the bin step rather than a tick spacing is carried in the pool key
  \item a v3 or CLAMM liquidity provider holds a per-range NON-FUNGIBLE position, wrapped as an ERC-721 by the NonfungiblePositionManager
  \item settlement in Infinity is deferred and net against an immutable ledger: vault.lock() then lockAcquired(data), operations run against a pool manager which calls vault.accountAppBalanceDelta(...), and the caller reconciles with take / settle / mint / burn (plus clear() for dust and the three-step sync-transfer-settle), with unsettled deltas counted in transient storage by SettlementGuard through UNSETTLED\_DELTAS\_COUNT and CURRENCY\_DELTA
  \item custody and pricing are SPLIT ACROSS CONTRACTS: an immutable Vault holds all balances and the accounting, while pricing lives in independent singleton pool managers (CLPoolManager and BinPoolManager) that register against it, so two AMM families share one ledger and tokens are transferred to the Vault rather than to the pool manager
  \item a pool may name a hook, and WHICH callbacks fire is carried in the POOL'S OWN KEY: the first 16 bits of the bytes32 parameters field are the hook registration bitmap and the next field is tick spacing (CL) or bin step (Bin) — the docs decode 0x...0a00c2 as permission bits 0000 0000 1100 0010 (afterInitialize, beforeSwap, afterSwap) with tick spacing 10 — so one hook contract may be registered with DIFFERENT permission sets by different pools and no address mining is required
  \item the current generation publishes no protocol-level price series: Infinity ships no oracle in the core and a custom oracle is an advertised hook use case
  \item the protocol's share of the fee is set by a swappable controller: ProtocolFeeController is Ownable2Step with protocolFeeSplitRatio = 33 * 1e4 (33\% of the total fee in hundredths of a bip) and defaultProtocolFeeForDynamicFeePool = 300, rejects a dynamic-pool default greater than 0.4\%, and exposes owner-only setProtocolFeeSplitRatio, setDefaultProtocolFeeForDynamicFeePool, per-pool setProtocolFee and collectProtocolFee
  \item the swap fee is split THREE WAYS at the protocol level, per tier, with the burn as one of the three legs: EVM v3 tiers pay LP 67 / burn 15 / treasury 18 at 0.01\%, 66 / 15 / 19 at 0.05\%, and 68 / 23 / 9 at both 0.25\% and 1\%, while Solana pools pay a uniform 84 / 8 / 8
  \item THE VALUE-RETURN CHANNEL IS SUPPLY DESTRUCTION, funded from five named sources: spot trading 15-23\% of trading fees, perpetual trading 20\% of all profits, CAKE.PADs 100\% of fees, Prediction 3\% of each round and Lottery 20\% of CAKE played, targeting 'an annual deflation rate of at least \textasciitilde{}4\% per year and a total CAKE supply reduction of \textasciitilde{}20\% by 2030'
  \item the terminal bound on supply is itself governable: the CAKE hard cap was lowered from 450M to 400M by a proposal that passed 2026-01-16
  \item the emission is directed across pools by an admin-mutable allocation-point vector, and the docs do not name the authority that sets it, calling it only 'the Chefs'
  \item the core is not upgradeable — the Vault is stated immutable and the docs assert 'The non-upgradeable core ensures stability' — while ownership of the pool-manager owner and of the fee controller moves by 2-step transfer, and the identity of the owner and of the pausable role is not published
  \item the protocol RETIRED a whole mechanism and settled the holders of a third party's derivative of it: veCAKE, gauge voting, farm boosting and the 5\% revenue share ended on 2025-04-23 (final gauge vote Epoch 37), Epoch 38 ran 2025-04-25 to 05-06, accrual and revenue share stopped 2025-05-07, a 6-month redemption window for directly staked CAKE closed 2025-10-23, third-party locker protocols were handled by whitelisting delegator addresses for 1:1 redemption, and the freed revenue-share allocation was 'redirected to the CAKE burn mechanism, increasing the burn rate for these pools from 10\% to 15\%'
\end{itemize}
\end{remark}

\begin{remark}
The corpus assigns a mutable implementation proxy on the strength of a mutable admin; at the current generation there is no proxy, the vault is stated immutable and the documentation asserts a non-upgradeable core, so $Up$ is dropped and what is actually there --- a fee controller swappable wholesale by an owner call, behind two-step ownership --- is left as residue. $Gp$ is added, which the corpus does not record at all: CLPoolManagerOwner is a PausableRole and exposes a pause reachable by a role strictly weaker than the owner while unpause is owner-only, so the emergency authority is one-directional. $St$ is carried for the deployed stable pool type, but the witness refuses to stretch it over the bins, because constant sum is a degenerate limit of $St$ in this vocabulary and not an instance of it. The corpus's summary of this protocol --- Uniswap with an emission schedule and a mutable admin --- fails on architecture rather than on flavour, and the split singleton, the second AMM family and the pool-key hook binding are the demonstration.
\end{remark}

\subsection{Raydium}\label{sub:cat:dex:raydium}

A provider deposits both sides of a pair into a constant-product pool and receives a fungible token claim on the two vaults, or chooses a price range in the concentrated pool and receives a position token whose mint authority is the program itself. The constant-product programs carry no collect-fees instruction at all, so fees stay in the vaults, inflate the invariant and are realised only on exit, whereas a concentrated position accrues against fee-growth globals differenced from its own checkpoint and is claimed explicitly. Price comes from pool state alone --- no external feed is consulted on a swap, a deposit or a withdrawal --- and settlement is atomic within a single instruction, a separate routing program composing a bounded number of pools inside one transaction. A position ends by burning the claim token, or by decreasing liquidity to nothing, collecting, and burning the position token. Nothing here can be liquidated: a concentrated position whose price leaves its range stops earning and converts wholly into one side, which is idleness rather than seizure.

\begin{measurement}\label{meas:cat:dex:raydium}
Take $X = \{Cl,Cp,Em,Gp,Ix,Sh,Tp,Up\}$, with canonical form
$\mathrm{ex}(X) = \{Cl,Cp,Em,Gp,Ix,Sh,Tp,Up\}$, every element being primitive in it.
It is admissible: no requirement term is open, every element is warranted, and it arms no prohibition. It discharges 8 of the 22 recorded obligations; 14 are residue. No composite of at most three corpus protocols equals it; 1 corpus protocol is contained in it, and together they supply every element except \{Cl,Cp,Tp\}.
It is not minimal: \{Cl\} can be removed with every obligation still covered and admissibility intact.
\end{measurement}

\begin{remark}
The stable program and the launchpad each price against a curve the vocabulary cannot name, and one of them publishes no invariant at all: it is closed source, ships no interface description, and its curve is stated nowhere at a primary source. Reproducing this protocol therefore means reproducing a description rather than a codebase over a substantial part of it, and no registry attests that any deployed program matches any published source --- verifiability is a property of a deployment rather than of code, and the vocabulary has no way to record its absence. Other residues concern state rather than mechanism: tick-array accounts can never be closed and their rent is locked permanently, so the footprint can only grow, and the oldest program keeps the account layout of a mechanism deleted from its own bytecode, leaving an on-chain type provably wider than the on-chain semantics. The fee obligations fail differently again, since which leg of a trade the fee is taken from is fixed at pool creation and a pool's tier is a reference to an immutable configuration it can never change --- so a fee here is a constant of the pool rather than a rate the protocol sets.

The obligations no element discharges are these.
\begin{itemize}\setlength{\itemsep}{0pt}
  \item a Stable AMM program is deployed and live at 5quBtoiQqxF9Jv6KYKctB59NT3gtJD2Y65kdnB1Uev3h, and no invariant for it is stated at any primary source: the program is closed source and publishes no IDL
  \item a CLMM liquidity provider instead holds a POSITION NFT — a supply-1 mint whose mint authority is the program — against a chosen price range, claims fees explicitly, and closes by decreasing liquidity to zero, collecting, then burning the NFT
  \item which leg the fee is taken from is fixed at pool creation and can never change: CLMM's CollectFeeOn is one of FromInput, Token0Only or Token1Only, so a pool may take its fee out of swap OUTPUT rather than input, changing the quoting equation by direction
  \item a pool's fee tier is a reference to an immutable AmmConfig chosen at creation and can NEVER be changed for that pool, with 19 live CPMM configs spanning 0.005\%-4\% and 18 live CLMM configs spanning 0.01\%-4\%
  \item the trading fee is split at the protocol level: CLMM and CPMM pay 84\% to LPs, 12\% to buyback and 4\% to treasury, while AMM v4 pays 88/12 with no treasury slice — verified against every live config, all of which carry protocolFeeRate 120000 and fundFeeRate 40000 on a 1,000,000 denominator, whereas AMM v4 uses a 10,000 denominator on which 25 means 0.25\%
  \item CPMM additionally charges a CREATOR fee that is not a slice of the trade fee at all — live configs show creatorFeeRate from 500 to 14,950 — plus a flat createPoolFee of 150,000,000 lamports (0.15 SOL) on CPMM and AMM v4
  \item THE VALUE-RETURN CHANNEL IS ACCUMULATION: the 12\% buys RAY on the open market and HOLDS it in a protocol-controlled account, DdHDoz94o2WJmD9myRobHCwtx1bESpHTd4SSPe6VEZaz, which holds about 83.54M RAY across six token accounts, while RAY total supply stands at 554,997,631.209468 against a 555,000,000 maximum — on the order of 2,400 RAY has ever been destroyed
  \item settlement is atomic within one Solana instruction against per-pool vaults, and multi-hop is bounded: the AMM Routing program CPIs across at most two pools in a transaction
  \item the protocol's on-chain state footprint can only GROW: CLMM TickArrayState accounts (60 ticks each, lazily initialised) can never be closed and their rent is permanently locked even when no tick in them is initialised, and CPMM pools and their PDAs are never closed even at zero liquidity
  \item AMM v4 retains a FOSSILISED INTERFACE for a mechanism that no longer exists: the 2026-07-22 upgrade removed the OpenBook/Serum dependency, all related CPIs and seven instructions (Initialize(0), MonitorStep(2), MigrateToOpenBook(5), WithdrawSrm(8), PreInitialize(10), SimulateInfo(12), AdminCancelOrders(13)), yet AmmInfo/StateData remain byte-compatible with inline open\_orders / market / market\_program fields as inert references and v1 swaps still demand 17-18 accounts of which several are never validated, with SwapBaseInV2/SwapBaseOutV2 (tags 16/17, 8 accounts) as the recommended path
  \item LaunchLab issues a new token against a bonding curve with NO pre-seeded liquidity, with curve\_type fixed at Initialize: 0 = constant product on VIRTUAL reserves, (V\_q + dq)(V\_b + b\_rem - b\_out) = V\_q*V\_b, 1 = fixed price, 2 = linear price
  \item the launch GRADUATES into a pool price-continuously and permissionlessly: base\_supply\_graduation is typically 0.8 x base\_supply\_max with the remaining 20\% seeding the pool, so the pre-graduation invariant is literally the post-graduation CPMM invariant; graduation fires when quote\_vault.balance >= quote\_reserve\_target, and 'Graduate is permissionless' — anyone may call it — whereupon the program CPIs CPMM CreatePool, revokes the base mint authority and sets LaunchState.status to Graduated
  \item the graduated pool's LP is split three ways by PlatformConfig under strict equality, platform\_scale + creator\_scale + burn\_scale = 1,000,000 checked by MigrateNftInfo::check, where the leg named 'burn' routes to the Lock program with is\_burn = true and the docs concede 'Despite the name, the LP tokens / position NFT are LOCKED in a program-owned escrow, not destroyed', while Burn\&Earn mints a transferable Raydium Fee Key NFT carrying the perpetual right to collectFees
  \item no registry attests that any deployed program matches any published source: OtterSec reports is\_verified false for AMM v4, CPMM, CLMM and LaunchLab, CPMM's record carries an empty executable\_hash, and the docs' stated mitigation — expected hashes per deploy in the repositories' releases — does not exist, because all three program repos have zero tags and zero releases
\end{itemize}
\end{remark}

\begin{remark}
The corpus files this protocol's fee channel under the same buy-and-burn heading as the protocols that burn, and the grouping is invalid: they destroy their token, this one buys it on the open market and holds it in an account it controls, its supply standing essentially unchanged against its stated maximum. $Fd$ names distribution to a claim class and no claim class is paid, so it is dropped and accumulation is left as residue --- which obliges the category to treat distribution, destruction and accumulation separately where the record carried a single heading. $Tp$ is added, which the corpus omits: the constant-product and concentrated programs each maintain an on-chain price series in a ring account that downstream protocols consume, with the qualification that the oldest program, holding the largest share of the inventory, keeps none, so the symbol holds of the newer programs rather than of the protocol. $Up$ and $Gp$ rest on chain evidence rather than documentation --- every queried program shares one upgrade authority, none is immutable, and the pause surface belongs to that same admin.
\end{remark}

\subsection{Aave V3}\label{sub:cat:lnd:aave}

A supplier deposits a fungible token into one shared reserve and receives a transferable claim token whose nominal balance grows as the reserve's index accrues. A borrower marks supplied reserves as collateral and draws an open-ended loan from the same reserve, with no maturity and no schedule, recorded as a debt token accruing by that same index. The price of the loan is a memoryless two-slope function of the reserve's current utilisation, held in a mutable per-reserve struct that only the pool configurator may write, and bounded at write time on its level, on where its kink may sit and on its shape. Solvency is a single scalar computed on chain from the oracle value of each collateral against the oracle value of the debt. A borrower closes by repaying; anyone else closes by repaying part of the debt and seizing collateral at a bonus, and how much of a position one liquidation may take is itself a computed quantity, falling out of how far underwater the position is and how large it is.

\begin{measurement}\label{meas:cat:lnd:aave}
Take $X = \{Aw,Ct,Ex,Fd,Fl,Gp,Ix,Li,Pl,Rb,Tg,Up,Xm\}$, with canonical form
$\mathrm{ex}(X) = \{Aw,Ex,Fd,Fl,Gp,Ix,Li,Pl,Rb,Tg,Up,Xm\}$; the elements \{Ct\} are derived rather than chosen.
It is admissible: no requirement term is open, every element is warranted, and it arms no prohibition. It discharges 12 of the 23 recorded obligations; 11 are residue. No composite of at most three corpus protocols equals it; no corpus protocol is contained in it.
It is not minimal: \{Rb,Tg,Xm\} can be removed with every obligation still covered and admissibility intact.
\end{measurement}

\begin{remark}
The residue here is the entire surface a risk team manages. What the vocabulary keeps is that money is pooled, that an index accrues, that a test is run, and that someone is paid to close a failing position; what it loses is every quantity those mechanisms are actually set to --- the shape of the rate curve and the interval its parameters must lie in, the allowance one liquidation may take and the remainder it may not leave, the caps that bound an asset's whole population rather than any individual position, and the permitted drift of a price feed. The sharpest instance is the steward: a named address, fixed at deployment, may move rates, caps, collateral parameters, correlated-asset categories and oracle growth limits with no vote, and of that mandate the decomposition names only the agent, and names it only by stretching an address-keyed permission gate. The enumerated domains it may touch, the maximum step per update, the minimum interval before the same parameter of the same asset may move again, and the principal's power to empty the envelope while being unable to remove the agent are all residue --- and across every instance of this shape in the category the domain, the magnitude cap and the revocation are residue without exception. Here the parts are not merely present but factored into one reusable validator, and that validator bounds the rate through the same code that bounds the caps, which is why the price of credit and the delegated mandate are one gap stated twice rather than two gaps.

The obligations no element discharges are these.
\begin{itemize}\setlength{\itemsep}{0pt}
  \item a liquid-staking-token feed is wrapped in an adapter that rejects any published value implying growth faster than a governance-set maximum yearly ratio, live example 10.64 per cent for wstETH
  \item how much of a position one liquidation may take is itself a computed quantity: at most half the debt when the health factor is above 0.95 and both the collateral and the debt exceed 2000 units of base currency, otherwise all of it, and in no case may the call leave less than 1000 base units of either side standing
  \item borrowing power is raised above the ordinary per-asset limits when the collateral and the debt belong to one governance-declared category of assets expected to track each other
  \item each reserve carries a supply cap and a borrow cap that bound aggregate exposure to the asset independently of whether any individual position is solvent
  \item the borrow rate is a memoryless two-slope function of the reserve's current utilisation alone -- base rate, a first slope up to a kink, a steeper second slope beyond it -- and the supply rate is that borrow rate scaled by supply utilisation and by one minus the reserve factor
  \item whatever curve is written must satisfy hard bounds checked at write time: the sum of base rate and both slopes may not exceed 1000 per cent, the kink must lie between 1 and 99 per cent, and the first slope may never exceed the second
  \item inside the strategy contract a permitted caller moves the entire curve in one transaction: there is no delay before the change takes effect and no limit on how far it may move
  \item that address may act only within an enumerated set of parameter families -- caps, rates, the collateral side, e-mode categories, and three kinds of oracle price cap -- and any individual asset or e-mode may be fenced out of its reach entirely
  \item each parameter carries its own maximum change per update, expressed either as a percentage of the current value or as an absolute step, and an update outside that range reverts
  \item no parameter may be moved again until a minimum interval has passed since that same parameter was last moved on that same asset, each pair carrying its own timestamp
  \item the principal may zero the envelope or fence any asset at any time, but cannot remove the agent: the agent's address is fixed at deployment and this contract holds no way to change it
\end{itemize}
\end{remark}

\begin{remark}
The corpus records $Bs$, $Cd$, $Em$, $Im$ and $Sl$; the witness carries none of them and adds $Aw$. $Im$ goes because the reserve configuration in the current tree holds no debt-ceiling field at all --- the bitmap's own comment records the unoccupied bits left where one used to sit --- so the fenced market the corpus called forced is not merely forced but unsupported, and the same drop follows downstream at the protocol that forked this one. $Bs$ and $Sl$ go because the Safety Module contracts were not read and this lane carries no element on memory; $Cd$ and $Em$ go for the same reason; $Aw$ arrives for the address-keyed administrative gate through which every parameter this protocol has is written. That leaves the exit-liquidity term of the pooled-claim law open where the corpus closed it, and the choice is deliberate rather than an omission: what actually guarantees a supplier's exit here is that a withdrawal reverts when the reserve is short and the curve then pays borrowers to repay, which is an incentive and not one of the mechanisms the term offers.
\end{remark}

\subsection{JustLend V1}\label{sub:cat:lnd:justlend}

A supplier deposits into a per-asset market and receives a token representing a pro-rata claim on it, whose redemption value rises through a growing exchange rate rather than through a growing balance. A borrower enters markets explicitly and a controller contract computes account liquidity from each entered collateral's collateral factor, valued by a push oracle whose price a permitted party writes rather than a feed the protocol pulls. Interest accrues into a per-market borrow index advanced every block a transaction touches the market, and the rate is quoted per block, so what a loan costs over a year depends on how fast the chain produces blocks. A liquidator repays at most a fixed fraction of the outstanding borrow and seizes collateral worth the repayment times a multiplicative premium, that fraction being a governed constant which does not vary with how far underwater the position is. The risk parameters carry permitted intervals checked in code; the rate curve carries none, and may be rewritten whole in a single call by one owner address.

\begin{measurement}\label{meas:cat:lnd:justlend}
Take $X = \{Aw,Ct,Ex,Gp,Ix,Li,Pl,Sh,Tg,Up\}$, with canonical form
$\mathrm{ex}(X) = \{Aw,Ex,Gp,Ix,Li,Pl,Sh,Tg,Up\}$; the elements \{Ct\} are derived rather than chosen.
It is admissible: no requirement term is open, every element is warranted, and it arms no prohibition. It discharges 10 of the 18 recorded obligations; 8 are residue. No composite of at most three corpus protocols equals it; no corpus protocol is contained in it.
It is not minimal: \{Sh\} can be removed with every obligation still covered and admissibility intact.
\end{measurement}

\begin{remark}
The instructive residue here is an absence. Every other protocol in this category carries a bounded delegate over something, and this one carries none; its risk parameters are bounded in code while its rate parameters are bounded by nothing whatever, a line the protocol deliberately drew and the decomposition cannot record. An absent bound and a present bound yield the same element set, so the vocabulary reports that a delay exists and cannot report that the delay is the only protection there is. That the primitive is missing here also dates it: this is a fork of a design older than the bounded mandate, and the stewards, curators and automators elsewhere in the category all arrived later and in repositories separate from the cores they govern, so the gap in the vocabulary is one of vintage rather than of scope. The rest is the category's own --- a rate quoted per block, so that a property of consensus leaks into a credit parameter, and a liquidation allowance fixed as a constant where the largest protocol here computes one.

The obligations no element discharges are these.
\begin{itemize}\setlength{\itemsep}{0pt}
  \item the rate is quoted and accrued per block rather than per second or per year, with a hard-coded 10,512,000 blocks per year matching a three-second chain, so the realised annual rate depends on how fast the chain produces blocks
  \item one liquidation may repay at most a fixed fraction of the outstanding borrow, the fraction being a governed constant that does not depend on how far underwater the position is
  \item the risk parameters are bounded in code at write time: a collateral factor may not exceed ninety per cent, a close factor must lie between five and ninety per cent, and a liquidation incentive must lie between one and one-and-a-half
  \item the borrow rate is a memoryless two-slope jump-rate function of current utilisation, and the supply rate is utilisation times the borrow rate net of the reserve factor
  \item nothing bounds what that curve may become: there is no maximum rate, no minimum or maximum kink, no requirement that the jump multiplier exceed the ordinary one, and no limit on how far a single update may move any of them
  \item every risk and every rate parameter in the protocol is set directly by one administrative address, with no per-parameter envelope, no step limit, no cooldown and no delegate of any kind
  \item the parameter changes that were actually executed are themselves committed to the repository as deployable contracts
  \item no contract in this repository rents chain execution resources, so the resource-rental business the corpus attributes to this protocol is not part of the artefact decomposed here
\end{itemize}
\end{remark}

\begin{remark}
The corpus records $Em$; the witness drops it, the controller having been read with no distribution mechanism in it, and adds $Aw$ for the single administrative address that writes every parameter and $Tg$ for the delay contract, the on-chain governor and the governance token standing behind it. Nothing else is added, and that is the point of this witness: it is the control in its category. No steward, no curator, no keeper and no bounded operator appears anywhere in the tree, and the decomposition would be identical whether the rate curve carried a permitted interval or none --- which is the strongest argument in the category that what is missing is a field on an element rather than an element. The corpus files the liquidation allowance here as though it were peculiar to this protocol and its ancestors; it is not, and the witness carries that residue under a category-wide reading with the three distinct policies the category actually runs.
\end{remark}

\subsection{Maple}\label{sub:cat:lnd:maple}

A lender deposits a single asset into a pool contract implementing the standard vault interface and receives pool shares whose value rises as borrowers repay; who may deposit is decided by a permission manager, so the lender is a known party rather than any address. The money is lent by a named pool delegate as dated loan contracts to named borrowers, with a payment schedule, a grace period and a balloon at maturity, and a second module issues loans with no maturity at all. The price of credit is not an output of the protocol but a term of a bilateral contract: a fixed annualised rate written into the loan at origination, with a fee on late payment and a premium added to the rate while a payment is late, unbounded in the loan contract and unmoved until both parties agree a new one by hash commitment. Whether a loan has defaulted, and when, is a judgement the delegate makes and executes --- closing a bad loan is an authorised act by a named party rather than a race, and no third party is paid to perform it. A loss draws down the delegate's own posted capital before it reaches the pool, and a lender leaves through a queue rather than on demand.

\begin{measurement}\label{meas:cat:lnd:maple}
Take $X = \{Aw,Bs,Ct,Fd,Ft,Gp,Pl,Sh,Sl,Sv,Tg,Up,Wq\}$, with canonical form
$\mathrm{ex}(X) = \{Aw,Bs,Fd,Ft,Gp,Pl,Sh,Sl,Sv,Tg,Up,Wq\}$; the elements \{Ct\} are derived rather than chosen.
It is not admissible: it leaves the requirement term $L1\!:\!Ex{\mid}Tp{\mid}At$ open. It discharges 10 of the 23 recorded obligations; 13 are residue. No composite of at most three corpus protocols equals it; no corpus protocol is contained in it.
\end{measurement}

\begin{remark}
The residue at this protocol is the protocol. There is no curve, no oracle in the pool layer, no permissionless liquidation, and no permissionless action anywhere on the credit side: depositing is gated, originating is delegated, defaulting is authorised, and the only act left to a stranger is a queued exit. What cannot be named is therefore not a parameter but a form of authority --- discretion whose precondition is the agent's own stake, a domain fixed by registered kind rather than by address so that any future instance of a kind is admitted automatically, a cap set by the principal rather than by the agent and sized to the pool rather than to any exposure inside it, and powers that exist only until configuration completes and then move elsewhere or vanish; each of those is a field the mandate schema would need and has not got, and the last is not a field at all but a lifecycle. This is also where the cross-lane characterisation of the delegate as accountable only reputationally fails, and fails in code: the agent here is bonded, the sanction is financial and automatic rather than a matter of standing, and stake belongs in the schema as a slot whose default is empty rather than in the definition as a constant. And the largest part of the mechanism was never in the code to begin with --- the underwriting, the custody, the governing law and the recourse against a named borrower --- which is why an honest decomposition from source leaves the truth term open rather than fitting a symbol to something that is not there.

The obligations no element discharges are these.
\begin{itemize}\setlength{\itemsep}{0pt}
  \item the price of credit is a term of a bilateral contract, not an output of the protocol: a fixed annualised rate written into the loan at origination, alongside a closing rate, a late-fee rate and a premium that raises the interest rate once a payment is overdue past the grace period
  \item the rate is agreed by the delegate and the borrower, and a loan cannot be funded until the borrower has accepted the terms
  \item nothing in the loan contract bounds that rate: there is no maximum, no minimum and no permitted interval on any of the four rate fields
  \item the rate does not move until both parties agree a new one: a change of terms is proposed as a hash commitment and takes effect only when the counterparty accepts it
  \item closing a bad loan is an authorised act by a named party rather than a race: the delegate or a protocol admin triggers a default, no third party is paid to do it, and the process may span several transactions and be finished later by a second authorised call
  \item the delegate may act only through instances the protocol registry recognises by kind -- a strategy, a liquidator factory, a withdrawal manager, a permission manager -- rather than through an enumerated list of addresses
  \item the delegate cannot originate a loan while its own posted capital is below the protocol's floor, and cannot withdraw that capital below the same floor
  \item the pool has a maximum size it may accept, and after configuration it is the protocol's admins and not the delegate who set it
  \item nothing limits how fast the delegate may lend: within the cap it may originate at will, with no cooldown, no step limit and no throughput budget
  \item the protocol's admins may replace the delegate, the registry alone may deactivate the pool, and every function in the pool manager stops while the protocol is paused
  \item several powers exist only until the pool is configured and then move to a different party or vanish
  \item there is no permissionless action anywhere on the credit side of this protocol: depositing is gated, originating is delegated, defaulting is authorised, and the only act anyone may take unbidden is to leave, which is queued
  \item no price source exists in the pool layer at all: the pool manager holds no oracle, and where the secured book's collateral is valued is not established on chain
\end{itemize}
\end{remark}

\begin{remark}
The corpus records $At$, $Em$, $Ex$, $Ix$, $Li$ and $Uc$; the witness carries none of them and adds $Bs$, $Fd$ and $Sl$. $Ix$ goes because the pool layer holds no index object, the share price rising because the pool's assets rise; $Ex$ goes because the pool manager holds no oracle, which leaves the truth term of the credit law honestly open, valuation here being a human act performed off chain; $Li$ goes because it is refuted in both clauses of its own definition, default being triggered by the delegate or a protocol admin, possibly across more than one transaction, with no bounty paid to any caller; $Uc$ goes because a second lane measured this book overcollateralised, so carrying it would assert something false about the loans; $At$ and $Em$ go uncited. $Bs$ and $Sl$ arrive because the delegate's own posted capital is drawn down before the pool's, up to a governed fraction of it, with the waterfall in code and origination blocked outright while that capital sits below the protocol's floor --- and the same second lane, reading this protocol's documentation rather than its source, refutes $Bs$ and drops it, that documentation heading its section on first-loss capital as unused because most pools are managed by the sponsor. Both readings are recorded and neither is reconciled: the mechanism is present and enforced in code, its live balance was not read on chain by this lane, and a vocabulary that cannot distinguish a protocol with a backstop from one with a backstop contract deliberately left empty is the result rather than the confusion.
\end{remark}

\subsection{Morpho}\label{sub:cat:lnd:morpho}

Credit is organised as many independent markets inside one contract, each fixed at creation by its loan asset, its collateral asset, its oracle, its interest-rate model and its loan-to-value, and anyone may create one from a combination governance has enabled. A lender supplies to a chosen market and receives shares of that market alone; a borrower posts that market's collateral asset and is tested against its single fixed loan-to-value at its own oracle. The rate is not a function but a controller --- a static curve around a target rate that itself drifts with the market's distance from target utilisation, clamped in code --- and no party sets it, every parameter of the controller being a compile-time constant. A liquidator repays an unhealthy position and seizes collateral at an incentive the protocol derives from that market's loan-to-value rather than from a parameter of its own, and a debt that survives with no collateral behind it is written down against that market's own suppliers and cannot reach another. Most depositor money arrives through a vault rather than directly: the vault, not the depositor, is the lender of record, and a curator decides which markets that money may enter and at what size while an allocator moves it between them.

\begin{measurement}\label{meas:cat:lnd:morpho}
Take $X = \{Aw,Ct,Ex,Fd,Gp,Im,Ix,Li,Pl,Sh,Sl,Tg\}$, with canonical form
$\mathrm{ex}(X) = \{Aw,Ex,Fd,Gp,Im,Ix,Li,Pl,Sh,Sl,Tg\}$; the elements \{Ct\} are derived rather than chosen.
It is admissible: no requirement term is open, every element is warranted, and it arms no prohibition. It discharges 11 of the 24 recorded obligations; 13 are residue. No composite of at most three corpus protocols equals it; no corpus protocol is contained in it.
It is not minimal: \{Pl,Ix,Sl\} can be removed with every obligation still covered and admissibility intact.
\end{measurement}

\begin{remark}
Here the residue is what the protocol is for. Its design claim is that a market's parameters are chosen once and never edited, that governance may widen a permitted set and may never narrow it, and that nobody at all may move the price of credit; the vocabulary can record only the absence of an upgrade path and of a pause, and the absence of a symbol means that nothing was observed, never that nothing can exist. The curator's mandate fares no better. Its delay is named, and named only because this instance happens to implement a delay as a queued change --- lowering a cap takes effect at once while raising one must wait --- where the same slot elsewhere in the category is a debounce, a one-sided cooldown or a refilling budget and is named by nothing; the set of markets the money may enter, the size each may hold, and the owner's power to remove the agent without notice are all residue. And the vault is itself a depositor in markets it does not control, so what a shareholder holds is a claim on a lender rather than a claim on a loan, which is a containment between protocols that a carrier composing by union of element sets cannot state at all.

The obligations no element discharges are these.
\begin{itemize}\setlength{\itemsep}{0pt}
  \item anyone may create a market from any combination whose interest-rate model and loan-to-value are on governance's enabled lists, and nobody may edit it afterwards
  \item the protocol checks nothing whatever about that oracle -- not its source, not its decimals, not its liveness -- and the address can never be changed for the life of the market
  \item a borrow or a withdrawal is refused whenever it would leave the market's borrowed assets above its supplied assets, and that is the whole of a lender's exit guarantee -- there is no queue, no ordering, no buffer and no reserve
  \item the borrow rate is produced by a closed-loop controller: a static curve of steepness four around a target rate, and that target rate itself drifts continuously toward whatever level drives utilisation to ninety per cent
  \item no party may change any parameter of that controller: the steepness, the target, the adjustment speed and the bounds are compile-time constants in a contract with no setter and no admin
  \item the controller's output is clamped in code: the target rate may never leave the band from one tenth of one per cent to two hundred per cent a year, so the realised borrow rate lies between 0.025 and 800 per cent a year
  \item the only rate decision anyone ever makes is which interest-rate-model address a market is created with, and that address is then part of the market's identity and frozen for its life
  \item no owner function can alter an existing market's oracle, model or loan-to-value, there is no proxy and no pause anywhere in the singleton, and the permitted lists may only be widened -- an enabled model or loan-to-value can never be disabled
  \item the vault is itself a depositor in markets it does not control, so what a vault shareholder holds is a claim on a lender rather than a claim on a loan
  \item the money may enter only markets the curator has admitted, and a market may not be admitted at all unless its loan asset is the vault's own asset and it already exists on the singleton, with at most thirty markets in the queue
  \item no market may hold more of the vault's money than its own supply cap, checked on every allocation
  \item the owner may remove the curator or any allocator at any moment, immediately and without a delay
  \item an allocation must end where it began: the total withdrawn across markets in one call must equal the total supplied, so the allocator may permute the deposit and can never move value out of the vault
\end{itemize}
\end{remark}

\begin{remark}
The corpus records $Em$, $Fl$ and $Sv$; the witness drops all three and adds $Fd$. $Fl$ goes because the singleton was read and holds no flash-loan facility, $Em$ because nothing in this lane's evidence cites one, and $Sv$ because the corpus's own marker concedes the point: $Sv$ names servicing and determination discretion over a loan, and a curator never touches a loan. $Gp$ is kept but belongs to the vault and not to the core --- the singleton has no pause, and the guardian's powers reach only changes that have been queued and not yet taken effect --- and this is where the corpus record fails on its sharpest claim, since the singleton cannot be upgraded and is not admin-less, its owner holding powers over ownership, over the permitted model and loan-to-value lists, over a bounded per-market fee and over the fee recipient. Among the alternatives the mandate's terms leave open, no element is forced onto the curator: its discretion is over which markets exist for a depositor's money and at what size, exercised before anything goes wrong, and $Sv$ would name a mechanism this protocol does not have.
\end{remark}

\subsection{SparkLend}\label{sub:cat:lnd:spark}

A supplier deposits an asset into a single shared reserve and receives an interest-bearing claim token whose nominal balance grows with a per-reserve index; a borrower marks collateral, draws an open-ended loan with no maturity, and is monitored by one on-chain health factor. Closing is by repayment or by a third party repaying and seizing collateral at a bonus, permissionlessly, exactly as at the protocol this one forked. The price of credit is where the fork stops being a fork. The curve has the same two-slope kinked form, but its base rate and its rate at the kink are not stored numbers: they are a fixed spread over a savings rate read live, by direct call, out of another protocol's own accumulator. No party sets that rate in place --- the spreads are fixed in the strategy contract's constructor, the level is whatever the tracked protocol sets, and when that protocol moves its savings rate this one's kink moves on the next interaction with no transaction here at all.

\begin{measurement}\label{meas:cat:lnd:spark}
Take $X = \{Aw,Ct,Ex,Fl,Gp,Ix,Li,Pl,Rb,Tg,Up\}$, with canonical form
$\mathrm{ex}(X) = \{Aw,Ex,Fl,Gp,Ix,Li,Pl,Rb,Tg,Up\}$; the elements \{Ct\} are derived rather than chosen.
It is admissible: no requirement term is open, every element is warranted, and it arms no prohibition. It discharges 8 of the 24 recorded obligations; 16 are residue. No composite of at most three corpus protocols equals it; no corpus protocol is contained in it.
It is not minimal: \{Pl,Rb,Up,Tg\} can be removed with every obligation still covered and admissibility intact.
\end{measurement}

\begin{remark}
Two mandates sit in this residue and they fill the same slot by different mechanisms: one keeper maintains a headroom above observed usage under a cooldown that applies to raising a cap and not to lowering it, and one controller spends a budget that refills continuously with elapsed time. The second bounds throughput and the first bounds frequency, and a schema with a single rate-of-change field records them as the same guarantee. Worse, the admin may switch the second bound off entirely for any one key in a single transaction with no delay, and the shape as this run has stated it has a slot for revoking the agent and none for revoking the bound, which is the more dangerous of the two. Underneath both sits the liquidity line itself, which is not a pool, there being no depositor and no claim, and is none of the other mechanisms the vocabulary offers for moving value between protocols. Above everything else stands the price of credit, which here is not administered but read: a direct call into another protocol's accumulator, so that the rate on this book moves when a different governance moves, and the code of this protocol depends on one upstream while the price of its credit depends on another.

The obligations no element discharges are these.
\begin{itemize}\setlength{\itemsep}{0pt}
  \item the borrow rate has the same two-slope kinked form as Aave's, but its base rate and its rate at the kink are not stored numbers -- they are read live, on every call, as a fixed spread over another protocol's savings rate
  \item the quantity the rate tracks is fetched by a direct call into another protocol's own accumulator, not from any feed, medianizer or reporter
  \item no party sets the rate in place: the spreads are fixed in the strategy contract's constructor and can only be changed by deploying a new strategy contract and repointing the reserve at it
  \item nothing bounds the rate on SparkLend's own side: there is no maximum on the spread and no maximum on the level, which is whatever the tracked protocol sets
  \item when the tracked protocol changes its savings rate, SparkLend's kink rate moves on the next interaction with no transaction on SparkLend at all
  \item the lender of last resort is not a depositor: liquidity issued by another protocol is moved into SparkLend and into other venues by an operator acting on that other protocol's behalf
  \item the keeper may act only on assets that have been given a configuration, and an asset with none is inert -- the computation returns its cap unchanged
  \item each configured cap has a ceiling it may never exceed and a headroom it maintains above current usage, so the new cap is the lesser of current usage plus headroom and that ceiling
  \item a cap may be lowered whenever the keeper likes but may not be raised again until a cooldown has passed since it was last raised, and no cap may be touched twice in one block
  \item the admin may delete any asset's configuration and may remove the keeper's role outright
  \item every permitted action is identified by a key naming an action and a venue together, and a key with no budget cannot be used at all
  \item each key carries a ceiling that is the most the agent can move in any one action
  \item spent budget refills continuously with elapsed time at a per-key rate, so what bounds the agent is not how often it acts but how much it may move per unit of time
  \item the admin may also switch the bound off entirely for any single key in one transaction with no delay, after which the agent's budget is unlimited and the checks short-circuit
  \item an oracle written and maintained in this protocol's own repository is supplied as the price source for markets in a different lending protocol
  \item the protocol's code is a fork of one protocol and the price of its credit is a function of a second protocol's state, so it depends on two upstreams for two different things
\end{itemize}
\end{remark}

\begin{remark}
The corpus records $Em$ and $Im$; the witness drops both and adds $Aw$. $Im$ goes for the reason it goes upstream, the fork inheriting a reserve configuration that no longer carries a debt-ceiling field; $Em$ goes uncited; $Aw$ arrives for the named roles that move this protocol's caps and move another protocol's liquidity across venues. The corpus records this protocol and its upstream as identical, and that is refuted here on the one axis the corpus itself says matters most: same functional form, different setter, different bound, different latency, different party in control. The identity survives only because the vocabulary cannot see rates, which is the corpus's own thesis demonstrated rather than asserted --- and while the element set does stay close to a subset of the upstream's, the control machinery does not, since separate stop mechanisms, a cap automator, a health checker and a liquidity controller, each an independently audited repository, have no analogue upstream at all.
\end{remark}

\subsection{Ethena}\label{sub:cat:cdp:ethena}

A whitelisted counterparty delivers an accepted asset to the minting contract and receives the newly created unit atomically in return, priced so that the issuer takes no profit on either side; redemption is the same trade reversed. What backs the unit is not a vault of collateral but a book --- liquid stablecoins and short-duration real-world assets that hold value on their own, alongside volatile holdings each paired with a short futures position of approximately equal notional. The backing never sits on the venue that holds the hedge: it is held by named third-party settlement providers and only delegated to derivatives venues as margin. The peg is defended by primary-market arbitrage alone, whitelisted parties creating the unit when it trades above par and destroying it when it trades below. Nobody borrows here, so nothing is liquidated; a holder who is not whitelisted leaves by selling, or by a cooldown that parks the unit in a silo contract for a period the administrator may set.

\begin{measurement}\label{meas:cat:cdp:ethena}
Take $X = \{At,Aw,Bs,Dp,Em,Fz,Gp,Ix,Rd,Rf,Sh,Sr,Tr,Wq,Xf\}$, with canonical form
$\mathrm{ex}(X) = \{At,Aw,Bs,Dp,Em,Fz,Gp,Ix,Rd,Rf,Sh,Sr,Tr,Wq,Xf\}$, every element being primitive in it.
It is admissible: no requirement term is open, every element is warranted, and it arms no prohibition. It discharges 13 of the 22 recorded obligations; 9 are residue. No composite of at most three corpus protocols equals it; 1 corpus protocol is contained in it, and together they supply every element except \{Bs,Dp,Em,Fz,Ix,Rf,Sh,Sr,Tr,Wq\}.
It is not minimal: \{Sh,Ix,Bs,Em\} can be removed with every obligation still covered and admissibility intact.
\end{measurement}

\begin{remark}
The residue here is the mechanism \emph{itself}. What makes the unit a dollar --- the hedge, the price service, the order router, the position reconciliation, the custodial mirroring --- is a private program operated by a company, and the model sees none of it. Around that sit four objects with no symbol: a par right that suspends itself precisely in the price region where it would be exercised, a speed limit on creating the unit differentiated by asset class, a reserve whose composition is a continuous discretionary investment decision rather than an attested holding, and a standing administrative claim on everything that lands in the contract except the unit itself. This is also the one member of the category with no instrument for the price of credit at all: the payout is a bookkeeping decision executed as a transfer, and no rate anywhere in the system is set to defend the peg.

The obligations no element discharges are these.
\begin{itemize}\setlength{\itemsep}{0pt}
  \item the right to create or destroy the unit is suspended automatically whenever the unit's price diverges unfavourably from its stablecoin reference beyond a configured threshold
  \item creation and destruction are capped per block, separately for stablecoin and non-stablecoin assets, under a global cap adjustable by a development multisig
  \item the backing is a book of six strategy families — crypto basis, non-crypto basis, DeFi lending, institutional lending, tokenised real-world assets and liquid stablecoins — whose allocation across categories is changed continuously at the operator's discretion
  \item the protocol's revenue is the funding payment of perpetual markets it does not operate and does not govern, received as an ordinary external participant
  \item the backing never sits on the venue that holds the hedge: it is held by named third-party settlement providers in bankruptcy-remote structures and only delegated to derivatives venues as margin on a rolling four-hour cycle, so that on a venue's failure positions are treated as closed and exposure is limited to unsettled profit and loss between cycles
  \item the prices used to quote creation and redemption are computed by an off-chain service that ingests six centralised venues plus two oracle networks, normalises and validates them, and publishes the result; no contract reads a price feed to test solvency
  \item the level of the payout is decided off-chain: a foundation subsidiary calculates an annual percentage yield weekly as part of internal accounting and executes it as a transfer, and no rate anywhere in the system is set to defend the peg
  \item two multisigs hold the constructive authority — one may change the supported collateral set and the custodian set, the other may grant and revoke every operational role — and no delay applies to either
  \item the administrator may remove any token except the protocol's own unit from the contract to an address the protocol controls
\end{itemize}
\end{remark}

\begin{remark}
This member is a hedged synthetic dollar rather than a collateralised debt position, and the witness carries no credit element, no collateral test and no liquidation path, so the canonical form separates it from the rest of the category without being told the distinction. Three corpus elements are dropped, each on a primary source. The oracle element goes because no contract reads a price feed to test solvency --- the price service, the order router and the reconciliation service are a private program run by a company, and the on-chain divergence limit compares an order's own implied price against its reference rather than reading a feed; the timelock element goes because the published role matrix specifies no delay for any role; and the upgrade element goes because the version transition was a redeployment to a new address, which is migration rather than upgrade in place. An execution element the corpus does not carry is added, the mint being a signed maker quote settled atomically against its own request identifier, and the perpetual-funding element is deliberately not used: this issuer takes the funding of markets it neither operates nor governs, as an ordinary outside participant.
\end{remark}

\subsection{Liquity}\label{sub:cat:cdp:liquity}

A borrower deposits ether or a liquid-staking receipt into a trove and draws the unit against it. In the first version there is no interest at all, only a one-time origination fee whose level rises with recent redemption activity and decays back; in the second the borrower chooses their own annual rate and may change it whenever they like. The peg floor is a direct redemption right at face value, exercisable by anyone at any time --- taken from the riskiest positions first in the first version, and in the second from the lowest-rate positions first, with volume allocated across markets in proportion to how little each has insured itself. Liquidation runs without an auction: a stability pool burns the unit against the debt and takes the collateral, and if that pool is empty the debt and the collateral are redistributed pro rata to every remaining borrower, who never opted in; the second version inserts a just-in-time tier between the two. Below a system-wide ratio the first version changes regime, and in the second a market shuts down of its own accord and switches redemption to no fee with a collateral bonus, paying whoever unwinds it.

\begin{measurement}\label{meas:cat:cdp:liquity}
Take $X = \{Bs,Cd,Ct,Em,Ep,Ex,Fd,Li,Rd,Sl\}$, with canonical form
$\mathrm{ex}(X) = \{Bs,Cd,Em,Ep,Ex,Fd,Li,Rd,Sl\}$; the elements \{Ct\} are derived rather than chosen.
It is admissible: no requirement term is open, every element is warranted, and it arms no prohibition. It discharges 9 of the 22 recorded obligations; 13 are residue. No composite of at most three corpus protocols equals it; 1 corpus protocol is contained in it, and together they supply every element except \{Ep\}.
It is not minimal: \{Em,Fd\} can be removed with every obligation still covered and admissibility intact.
\end{measurement}

\begin{remark}
The price of credit is residue here for a reason no other member supplies: the borrower sets it, so it is not an instrument the protocol holds at all --- and because the same number fixes the borrower's place in the redemption queue, what is bought is priority rather than credit. Delegating that choice is the category's second instance of the five-part bounded mandate with the principal inverted: granted sideways by one user over their own position rather than downward by governance, and priced by a fee denominated in the very rate it governs. Around those sit a liability that is itself a transferable token, a position redeemed to dust that can no longer be redeemed but still exists and still accrues interest, a fee retained inside the redeemed borrower's collateral with no residual claimant at all, and three ordered liquidation tiers whose payer and penalty change at each step, against an element that presumes one path paid to one third party. Immutability is residue too, and it needs three values rather than two: no key and no governance in the first version, and in the second, contracts that cannot be changed alongside a governance system that can do exactly one thing. Recording that as the absence of an upgrade, a delay and a pause erases the distinction between them.

The obligations no element discharges are these.
\begin{itemize}\setlength{\itemsep}{0pt}
  \item in V1 the borrower pays a one-time fee at origination whose level rises with recent redemption activity and decays back, floored at half a per cent and capped at five, and pays no interest at all thereafter
  \item in V2 each borrower chooses their own annual interest rate and may change it whenever they like, and that same rate fixes their place in the redemption queue: redemptions consume the lowest-rate positions in a branch first, ties broken last-in-first-redeemed
  \item interest on each position is computed simply on its own recorded debt at its own rate and compounded only when the position is touched, while the branch keeps two running aggregates so that pending interest for the whole branch can be computed without iterating positions
  \item in V2 redemption volume is allocated across branches in proportion to each branch's debt minus that branch's own stability-pool deposits, so the least self-insured market absorbs the most redemption pressure
  \item redemption carries a fee equal to half a per cent plus a base rate that rises with redemption volume and decays on a six-hour half-life, and in V2 that fee is retained inside the redeemed borrower's collateral rather than paid to anyone
  \item in V2 liquidation is attempted in three ordered tiers with different payers and different penalties: offset against the stability pool at five per cent, then a just-in-time liquidation in which a liquidator supplies the stablecoin for a hundred and five per cent of the collateral value, then redistribution
  \item in V1, when the system-wide total collateral ratio falls below one hundred and fifty per cent the whole system changes regime: positions under 150\% become liquidatable, any transaction that would worsen the aggregate ratio is blocked, and the origination fee drops to zero to pull collateral in
  \item in V2 each branch degrades on its own: below a critical ratio it blocks withdrawals, new debt and premature rate changes, and below a lower shutdown ratio it closes permanently to new borrowing and switches redemption to zero fee with a two per cent collateral bonus
  \item a borrower may delegate the choice of their own interest rate to a registered manager who declares a minimum and maximum band, may do nothing else whatsoever, and reprices the whole delegated set in one operation
  \item the option to reprice is itself priced: opening or changing a rate costs seven days of the borrower's own chosen interest, changing it again within seven days costs a further fee, switching manager costs an upfront fee, and a manager's own annual fee can be lowered but never raised
  \item a position redeemed below the minimum debt is not closed: it enters a state in which it can no longer be redeemed against but still exists and still accrues interest
  \item the debt position is a transferable non-fungible token, so one address may hold many of them and any of them may be sold
  \item the contracts cannot be changed: V1 has no administrative key, no governance and no upgrade path, and V2's core contracts are immutable and non-upgradeable while a governance system exists that can do exactly one thing
\end{itemize}
\end{remark}

\begin{remark}
This witness is the mirror image of the largest member's: law $L7$ leaves the credit term open to a redemption right, a peg-swap or liquidation capacity, and here it is satisfied by the redemption right alone, there being no peg-swap anywhere in the protocol. The emergency-pause element is refused for branch shutdown and the refusal is the finding --- that element names a guardian and presumes a holder, and this has none, the shutdown being a pure function of a measured ratio, per market rather than global, and improving the redeemer's terms instead of freezing them; the delegated-scope element is refused for the rate delegate, which has no expiry, no nonce and no domain separation; and the vote-escrow element is refused because voting power accrues with time held rather than time locked, with no escrow and no chosen duration. Both an emission element and a fee-distribution element are carried, because the entity spans two versions in which the same incentive is funded first from token issuance and then from protocol revenue, and collapsing them would blur a real change of funding source. One control-plane fact has no field in the record at all: the second version's core contracts are source-available under a business licence naming a licensor and a future change date, and are not open source until that date, which is a counterexample to any claim that permissionless forkability characterises this category.
\end{remark}

\subsection{Lista}\label{sub:cat:cdp:lista}

A borrower deposits one of a set of accepted assets --- including the protocol's own liquid-staking receipt --- through a front contract into the vault, and borrows the unit against it. Nobody votes on the borrow rate: a market-operations module recomputes it on a fixed cadence, and on any user interaction, as an exponential function of the deviation of the unit's price from a dollar, parameterised per collateral and clipped at a ceiling. The peg is defended by that controller, by a peg-swap that is free to enter and taxed to exit under a daily ceiling, by direct minting of the unit into two third-party lending markets to seed liquidity there, and by liquidation. Anyone, including the borrower, may trigger a liquidation; the collateral is sold by an auction that opens just above the unit price and declines linearly, and the trigger takes a fixed reward plus a proportional one while the debt is grossed up by a penalty. Savings pay a fixed rate adjusted by hand, and a saver withdraws immediately when pool liquidity allows and within a bounded period otherwise.

\begin{measurement}\label{meas:cat:cdp:lista}
Take $X = \{As,Cd,Ct,Em,Ex,Gp,Ix,Li,Ps,Sl,Wq\}$, with canonical form
$\mathrm{ex}(X) = \{As,Cd,Em,Ex,Gp,Ix,Li,Ps,Sl,Wq\}$; the elements \{Ct\} are derived rather than chosen.
It is admissible: no requirement term is open, every element is warranted, and it arms no prohibition. It discharges 12 of the 20 recorded obligations; 8 are residue. No composite of at most three corpus protocols equals it; no corpus protocol is contained in it.
\end{measurement}

\begin{remark}
The rate residue here is a controller rather than a policy: a feedback law from peg error to borrow rate, recomputed on a cadence and clipped, with no one in the loop --- and the same protocol sets its savings rate by hand, so two instruments of different kinds run inside one balance sheet. The price source is an ordered chain of providers with a defined degradation order and a hard band that rejects a price outright instead of incorporating it; the oracle element names an aggregation and can state neither half of that. The peg-swap is free one way, taxed the other, capped per day, and documented as taxed on exit in order to push exit onto a third-party market maker, which is a materially different instrument from the reserve-backed par swap the symbol names. The control-plane residue is the sharpest: governance is an off-chain vote in which only the core team may propose, implementation is applied by hand afterwards, and the same party holds a standing right to overrule the outcome. In the vocabulary that whole arrangement decomposes to a pause and nothing else, so the absence of a timelock reads as \emph{no timelock found} rather than as \emph{no chain at all}, and those are opposite epistemic states.

The obligations no element discharges are these.
\begin{itemize}\setlength{\itemsep}{0pt}
  \item that price is taken from an ordered chain of providers — a main source, a pivot and a fallback, drawn from four independent networks and differing per asset — and is then checked against a hard per-asset band and rejected outright if it falls outside it
  \item the borrow rate is recomputed every fifteen minutes, and on any user interaction, as an exponential function of the deviation of the stablecoin's price from one dollar read from an exchange oracle, parameterised per collateral by a base rate and an amplification factor, and clipped at twenty per cent
  \item the seized collateral is sold by an auction that opens at the unit price times 1.02 and declines linearly toward zero over a fixed horizon, pauses when it has run past a time bound or fallen about forty per cent, and must then be restarted by someone who is paid for restarting it
  \item that swap is free in one direction and taxed two per cent in the other, is capped in size, is subject to a daily redemption ceiling of five hundred thousand units, and is documented as taxed on exit in order to push exit onto a third-party automated market maker
  \item the solvency of that minted supply then depends on two lending markets the protocol neither operates nor governs, and the rates its unit earns on those markets are themselves managed by its own controller
  \item the rate that pool pays is a fixed number adjusted by hand, inside the same protocol whose borrow rate is set by a closed-loop controller
  \item parameters are changed by an off-chain vote on a hosted platform in which only core team members may submit a proposal, voting runs three days, a proposal passes on more than half of the tokens cast, and implementation follows by hand one to two weeks later
  \item the protocol accepts its own liquid-staking receipt as one of the assets backing its own stablecoin, so that one of its products' liabilities is another of its products' collateral
\end{itemize}
\end{remark}

\begin{remark}
Two residues the corpus records for this protocol are refuted here with evidence, and the refutations are why the witness looks as it does. The corpus assigns it the same rate-as-policy gap as the largest member of the category; nobody votes on this rate, so what it actually arms is the closed-loop controller gap the corpus records one row later, and the witness carries no element for the rate because none exists --- quantity adjustment names quantity rather than price, the oracle supplies the input without naming the feedback law, and the index names accrual rather than level. The corpus also records a gas-compensation-only liquidation incentive, on the ground that a rebate rather than a discount changes who shows up; the documented incentive is a fixed reward plus a proportional one on top of a penalty on the debt, the same structure the largest member uses, so that residue is struck and the liquidation element is carried on ordinary terms. Two further corpus elements are dropped --- the timelock, because no delay is documented and no vote is on-chain, and the upgrade element, for want of a citation rather than because a proxy is believed absent --- and a quantity-adjustment element is added, flagged as approximate, for unbacked minting into venues the protocol neither operates nor governs.
\end{remark}

\subsection{Sky}\label{sub:cat:cdp:sky}

A borrower locks an approved collateral asset into a vault under a collateral type and draws the issuer's stablecoin against it, bounded below by a debt floor, above by a per-type and a global ceiling, and per unit of collateral by a price carrying its own safety margin. The debt is not a fixed sum: it accrues continuously through a per-second accumulator that one contract advances and the core ledger folds across every vault of that type at once. The peg is defended from four directions at once --- the borrow rate makes debt expensive, the savings rate makes the liability attractive to hold, a peg-swap module trades the unit against a fiat-backed stablecoin at par, and liquidation removes vaults whose collateral no longer covers them. An unsafe vault is confiscated by any keeper, its debt booked to the balance sheet, and its collateral sold by a descending auction that opens above the oracle price and is bought atomically with an optional callback, so the buyer needs no capital of their own. There is no direct redemption right against the collateral: the peg-swap is the only par exit, and a borrower otherwise ends the position by repaying and freeing what they locked.

\begin{measurement}\label{meas:cat:cdp:sky}
Take $X = \{Aw,Cd,Ct,Em,Ex,Fd,Gp,Ix,Li,Ps,Sh,Sl,Tg,Tr,Up,Xf\}$, with canonical form
$\mathrm{ex}(X) = \{Aw,Cd,Em,Ex,Fd,Gp,Ix,Li,Ps,Sh,Sl,Tg,Tr,Up,Xf\}$; the elements \{Ct\} are derived rather than chosen.
It is not admissible: it arms $X19*$. It discharges 15 of the 24 recorded obligations; 9 are residue. No composite of at most three corpus protocols equals it; 2 corpus protocols are contained in it, and together they supply every element except \{Aw,Fd,Tr,Xf\}.
\end{measurement}

\begin{remark}
The gap the corpus records for this category --- no instrument for the price of credit --- is here a bounded delegated mandate, and writing it out part by part shows why one symbol could not hold it. Only the agent slot is nameable, and only by stretching a permission gate that names who may send a transaction rather than who may move a parameter; the bounded quantity, the magnitude bounds and the revocation have no element at all. The step-and-cooldown slot is misnamed by the timelock element twice over, because the cooldown runs since an action rather than before one, and because the write takes effect immediately. The rest of the residue is the control plane: nothing in the vocabulary names a franchise, a voting rule, a quorum or a threshold, and there is no fixed quorum here to name --- the threshold is whatever the incumbent slate currently holds, which is a materially different governance risk and is invisible to the model. What the model does see is a delay and a proxy, and it cannot say who holds either.

The obligations no element discharges are these.
\begin{itemize}\setlength{\itemsep}{0pt}
  \item the level of the borrow rate, of the savings rate and of the savings-reserve rate is a written number: no utilisation curve, market, auction, target or measured quantity derives it anywhere in the system
  \item that agent may write only parameters drawn from an enumerated set of rate identifiers, each with its own configuration record, and cannot touch debt ceilings, collateral onboarding, liquidation parameters or anything else
  \item each value the agent writes must lie between a governance-set floor and ceiling for that identifier, and a write outside the band reverts
  \item a single update may move a rate by at most a governance-set step, and no update may occur within a cooldown of the previous one; the update then takes effect immediately
  \item settlement clears by a descending price: the auction opens above the oracle price, falls along a chosen curve, is bought atomically at the current price with an optional callback so the buyer needs no capital of their own, and is reset by anyone once it has run past a time bound or fallen below a price bound
  \item the balance sheet mechanically converts imbalance into an auction in both directions: surplus above a high-water mark is sold off in fixed lots, and a deficit beyond a fixed lot triggers an auction of newly issued governance tokens for the stablecoin
  \item the issuer maintains two of its own liabilities simultaneously, converts between them one-for-one and permissionlessly, and pays a different savings rate on each
  \item a sub-entity receives its own collateral type backed by an oracle that returns a constant price, purely so that a nominal collateral balance makes its debt ceiling reachable, and its operators then draw and repay the stablecoin into an off-protocol buffer with no vote and no delay, whitelisted per action and rate-limited
  \item the franchise is exercised by locking the governance token into an approval-voting contract and approving a slate; the slate holding the most approvals is promoted, and only the promoted address may schedule an action
\end{itemize}
\end{remark}

\begin{remark}
Law $L7$ leaves the credit term open to a redemption right, a peg-swap or liquidation capacity, and this witness satisfies it through the peg-swap alone: the issuer grants no direct claim on collateral, which is a stronger statement than the record makes. One corpus element is dropped and three candidates are refused, each at primary source. The attestation element is dropped because no on-chain attestation object was found behind the real-world collateral, so its law cannot be checked; the time-weighted price element is refused because a queued next price is a delay rather than an average, and the two have different manipulation profiles; the delegated-scope element is refused for the rate setter, because that law requires expiry, a nonce and domain separation over the principal's own assets, and a facilitator has none of the three while acting over the whole system's economics; and the vote-escrow element is refused because the approval-voting contract locks without time-weighting and without a chosen duration. Against those, a share element and a tranche element are added, the first correcting a treatment the corpus already applies to an identically shaped share token elsewhere in this category, the second forced, the junior class being documented rather than settled by a waterfall.
\end{remark}

\subsection{USDD}\label{sub:cat:cdp:usdd}

A borrower deposits collateral --- the host chain's native token, its staked variant, or a fiat-backed stablecoin --- into a vault and mints the unit against it, staying above a minimum ratio that varies with the asset's volatility, and the debt accrues a stability fee through a rate accumulator. The peg is defended by over-collateralisation, by a peg-stability module that swaps a fiat-backed stablecoin for the unit at par with no service fee, and by liquidation. An unsafe vault is seized automatically and its collateral sold by a descending auction whose price decays at a fixed rate per second; the keeper takes a fixed incentive plus a proportional one, the proceeds repay the debt, and any excess collateral returns to the owner. The architecture also provides for auctions selling governance tokens when the protocol runs a deficit. There is no direct redemption right against collateral: the peg-swap is the par exit, and a savings share pays an administered rate to holders who deposit the unit.

\begin{measurement}\label{meas:cat:cdp:usdd}
Take $X = \{Cd,Ct,Ex,Gp,Ix,Li,Ps,Sh,Sl\}$, with canonical form
$\mathrm{ex}(X) = \{Cd,Ex,Gp,Ix,Li,Ps,Sh,Sl\}$; the elements \{Ct\} are derived rather than chosen.
It is admissible: no requirement term is open, every element is warranted, and it arms no prohibition. It discharges 9 of the 17 recorded obligations; 8 are residue. No composite of at most three corpus protocols equals it; no corpus protocol is contained in it.
It is not minimal: \{Sh\} can be removed with every obligation still covered and admissibility intact.
\end{measurement}

\begin{remark}
What this protocol cannot say is mostly what its sources do not say. The savings rate is administered by a published reaction function whose inputs are in the main outside the protocol --- a moving average of competitor yields, a central bank's policy rate, a competitive premium --- and nothing in the vocabulary names a rate whose level is a function of external market and central-bank rates. Behind a nominally permissionless vault stands a discretionary allocator: two named teams choose the venues, the caps and the phasing, and only the execution is on-chain. The descending auction, the contract whose only purpose is to convert holders of the issuer's own predecessor liability, and the separate issuance in three domains are residue as well. The control-plane residue is of a different kind from the others in this category: elsewhere the vocabulary is too coarse to distinguish arrangements that are documented, and here the arrangement is not documented \emph{at all}, which the record states as UNKNOWN rather than as absence.

The obligations no element discharges are these.
\begin{itemize}\setlength{\itemsep}{0pt}
  \item the level of the borrow rate is a written number set by an authority the documentation does not name
  \item the advertised savings rate is administered by a published reaction function whose inputs are a seven-day moving average of competitor stablecoin yields, the United States Federal Reserve policy rate, a competitive market premium, and the protocol's own underlying yield
  \item the yield paid on that share is produced by deploying idle capital into third-party lending venues, where the venues, the caps and the phasing are chosen jointly by two named teams after review of market conditions and liquidity, and only the execution is on-chain
  \item the auction clears by a price that decays at a fixed rate per second and can be reset by anyone who is paid for resetting it
  \item the architecture provides for auctions that sell governance tokens when the protocol runs a deficit
  \item the issuer exists separately in three domains, with its own token contract and its own local collateral in each, and mints locally against local backing rather than bridging a single supply
  \item a core contract exists whose only job is to convert holders of the issuer's own predecessor liability into the current one
  \item the protocol states that its framework is community-driven and that critical changes are subject to community review and approval, and names no governable parameter, no timelock, no voting threshold and no administrative address
\end{itemize}
\end{remark}

\begin{remark}
This is the smallest construction of the five, and it is small because most of what the corpus recorded for it could not be confirmed at a primary source. Five corpus elements are dropped: the timelock, because no primary source names one and secondary reporting about a successor governance stack could not be tied to this issuer's own parameters; the upgrade element, because the proxy listed among the core contracts is the borrower's own proxy rather than an upgrade proxy; the attestation, because the reserve it named issued the predecessor liability and the current architecture and governance pages do not mention it; the cross-domain element, because nothing is bridged --- the issuer exists separately in three domains, with its own token and its own local collateral in each, so the conservation law the symbol carries does not apply and asserting it would assert a relation that does not hold; and the emission element, because the savings subsidy it stood for is not evidenced anywhere. A share element is added, correcting a treatment the corpus already applies to an identically shaped share token elsewhere in this same category. The control plane is not merely thin: no public repository for the CDP core could be located, the governance pages name no governable parameter, no delay, no threshold and no address, and whether the deployed bytecode matches any published source is UNKNOWN.
\end{remark}

\subsection{Babylon}\label{sub:cat:lsd:babylon}

A Bitcoin holder locks BTC on Bitcoin, in a Taproot output it constructs itself, and never bridges the coin or surrenders the key; the protocol holds nothing. No transferable claim is issued: the position is a delegation registered on a separate chain against a Bitcoin inclusion proof, naming a finality provider chosen at the moment of locking and never reassigned, and rewards accrue in that chain's own token through a per-stakeholder gauge, net of the provider's commission. Exit is a race between paths written into the output --- wait out the staking timelock, or broadcast an unbonding transaction the covenant committee has already pre-signed and then wait again, during which the coin has left the staking output and remains liable. What the stake secures is the finality of that chain, and no third party may define the fault: the only slashable event is a provider signing conflicting blocks at the same height, and the same satoshis cannot be committed elsewhere while the output stands. When a provider does equivocate, the signature scheme discloses its own key, and anyone may then complete and broadcast the pre-signed slashing transaction, which destroys a fraction of the delegation into an output tagged with the protocol's name.

\begin{measurement}\label{meas:cat:lsd:babylon}
Take $X = \{Aw,Bs,Em,Ep,Fd,Ix,Rs,Sl,Xm\}$, with canonical form
$\mathrm{ex}(X) = \{Aw,Bs,Em,Ep,Fd,Ix,Rs,Sl,Xm\}$, every element being primitive in it.
It is admissible: no requirement term is open, every element is warranted, and it arms no prohibition. It discharges 11 of the 22 recorded obligations; 11 are residue. No composite of at most three corpus protocols equals it; no corpus protocol is contained in it.
It is not minimal: \{Bs,Sl,Em,Fd\} can be removed with every obligation still covered and admissibility intact.
\end{measurement}

\begin{remark}
The vocabulary presumes that a protocol holds the stake it slashes, and here nothing is \emph{held}. The securing relationship is a disjunction of script conditions on a chain that does not know this protocol exists, together with a set of transactions signed before the events that would justify them and waiting only to be broadcast; a conditional obligation already executed is neither a queue nor an escrow nor an option. Enforcement is stranger still: nothing is assigned to anyone, a secret becomes public, and a stranger finishes the job. Where the candidate symbol for this category would compress allocation and deposit defence, neither arises --- the staker names its provider inside its own output, and never parts with the coin --- and the carrier writes an obligation that cannot arise exactly as it writes one it could not name. Its answer to who may author the fault is likewise negative, and a negative is the one answer the carrier has no way to record.

The obligations no element discharges are these.
\begin{itemize}\setlength{\itemsep}{0pt}
  \item A Bitcoin holder locks BTC on Bitcoin, in a Taproot output it constructs itself with the NUMS point as internal key so that key-path spending is disabled, and never surrenders custody and never bridges the coin; the protocol holds nothing.
  \item The output carries exactly three spending paths: a pure timelock returning the coin to the staker after at most 65,534 blocks, an unbonding path requiring the staker plus a threshold of lexicographically sorted covenant keys, and a slashing path requiring the staker plus the named finality provider plus that same covenant threshold.
  \item [Vl sub-mechanism 2 of 5: stake allocation and scheduling] The staker names its finality provider inside its own Bitcoin output at the moment of locking; there is no scheduler, no allocation authority and no reassignment.
  \item The unbonding and slashing transactions exist fully signed before the events that would justify them, and are merely waiting to be broadcast.
  \item At most 60 finality providers are active, selected as the top providers by delegated value that have BTC-timestamped public randomness for the height in question.
  \item A provider commits Merkle-tree public randomness for future heights, and the commitment takes effect only after it has been timestamped into Bitcoin; a block is final once providers holding more than two thirds of voting power have cast one-time-signature votes, in a second round on top of the chain's ordinary consensus.
  \item [Who may write a slashing condition over the collateral] No third party may. The only slashable fault is equivocation at a single height, defined by the chain's own finality module; there is no analogue of an external service authoring its own slashing predicate.
  \item [Whether the same collateral can be committed twice] It cannot without a new Bitcoin transaction: a staking output names its finality provider at creation and the covenant pre-signs against that exact output, so the same satoshis cannot be re-committed while locked.
  \item [Vl sub-mechanism 5 of 5, trigger: penalty attribution] A provider that signs two conflicting blocks at one height mathematically reveals its own private key through the extractable one-time signature scheme; the module records the equivocation evidence such that anyone can extract the key, after which anyone may complete and broadcast the pre-signed slashing transaction.
  \item A user who holds both an active BTC delegation to an active finality provider and a BABY delegation to a CometBFT validator earns an additional reward computed from a score over active satoshis and active BABY jointly.
  \item [Vl sub-mechanism 3 of 5: deposit front-running defence] There is none and none is needed: the staker constructs the output, holds its own key and never hands the coin to a third party, so there is no deposit for an operator to front-run.
\end{itemize}
\end{remark}

\begin{remark}
The withdrawal-queue element is dropped from the corpus construction, because there is no queue: no ordering, no shared exit liquidity and no finalisation step exists, and each delegation unbonds independently on its own Bitcoin timelock as a per-delegation state machine. Slashable capital and loss allocation are retained but on a different magnitude from the one the record assumes: the live mainnet rate destroys far less of a delegation than a first-loss framing implies, so both symbols carry almost no economic content here, and the sanction that bites is the provider's removal from the active set and the stake's consequent failure to earn. Cross-domain verification is claimed for reading Bitcoin, where a light client and an inclusion proof supply the chain binding and finality its warrant asks for, and deliberately not claimed for writing checkpoints into Bitcoin, because nothing is delivered, nothing is addressed, and the destination chain has no notion of a sender. Protocol-funded emission is retained here where it is dropped elsewhere in this category, because these rewards come from the chain's own minting and fee machinery rather than from a service that submits them.
\end{remark}

\subsection{Binance staked ETH}\label{sub:cat:lsd:wbeth}

Anyone may send ETH to the token contract and be minted WBETH at the posted rate, with no account, no KYC and no allowlist; through the exchange the same token is a product behind a login, and the rights behind the doors differ. The claim is not a share of anything: a holder's balance never moves, and the ETH entitlement rises because a stored scalar rises, written by a whitelisted caller rather than derived from any state the ledger can see. Exit begins on chain and unconditionally --- the WBETH is burned and a request is recorded in a separate unwrap contract --- and does not end there, because the request becomes claimable only once an operator marks it allocated, out of ether the contract is topped up with from an off-chain balance. A blacklister may strand a named holder's claim while leaving that holder's token transferable and its mint untouched. Nothing here is slashable: there is no validator, no key, no operator set and no loss rule inside the ledger, and a loss reaches the token only if someone chooses to write a smaller number.

\begin{measurement}\label{meas:cat:lsd:wbeth}
Take $X = \{Fz,Gp,Ix,Rd,Up,Wq\}$, with canonical form
$\mathrm{ex}(X) = \{Fz,Gp,Ix,Rd,Up,Wq\}$, every element being primitive in it.
It is admissible: no requirement term is open, every element is warranted, and it arms no prohibition. It discharges 7 of the 16 recorded obligations; 9 are residue. No composite of at most three corpus protocols equals it; no corpus protocol is contained in it.
It is not minimal: \{Rd,Wq\} can be removed with every obligation still covered and admissibility intact.
\end{measurement}

\begin{remark}
The residue here is not a list of mechanisms the vocabulary missed. It \emph{is} the protocol. What backs the token, whether it exists, who holds it and what becomes of it when a validator is penalised all sit outside the ledger, and the ledger's own machinery --- an asserted index, a governor on the assertion, a request that is enforced against a fulfilment that is not --- has no symbol either. The redemption path is the sharpest of these: the burn and the queue entry are code, the payment is a decision, and the queue element presumes an obligation to clear. The validator element, meanwhile, is inapplicable here rather than approximate, and the carrier writes that exactly as it writes a mechanism it could not name --- the position where the conflation costs most, because here the absence is the whole design.

The obligations no element discharges are these.
\begin{itemize}\setlength{\itemsep}{0pt}
  \item Only whitelisted callers may write the exchange rate, and each write is bounded by a per-caller allowance that refills on a per-caller interval; a write that is zero, unchanged, or exceeds the remaining allowance reverts.
  \item A registered request becomes claimable only after an operator marks it allocated and after a minimum lock of two days (MIN\_LOCK\_TIME, with a settable lockTime); the contract is under no obligation to be funded, and is topped up from an off-chain balance through rechargeFromRechargeAddress.
  \item The backing of the token is an off-chain balance sheet: no attestation object naming an attester, an assurance level, a frequency or a staleness bound is published for WBETH or BETH.
  \item [Vl, all five sub-mechanisms] The protocol performs no validator action on chain of any kind: there is no validator registry, no key registration or vetting, no withdrawal-credential proof, no allocation of stake across operators, no deposit front-running defence, no exit request or exit signalling, no slashing record and no operator set. The three writes that touch the staking side (moveToStakingAddress, supplyEth, setNewEthStaked) move value to a Binance-controlled address or assert a number, and the trail ends on chain.
  \item [What the collateral secures, who may write a slashing condition over it, whether it may be committed twice] Nothing on chain, no one, and unanswerable: there is no restaked collateral, no slashing condition, no party empowered to author one, and no on-chain record of the backing position against which a second commitment could be detected.
  \item The same ERC-20 is simultaneously a permissionless on-chain claim mintable by anyone at a posted rate and an exchange product gated by an account and by terms of service, with different rights behind each door; through the exchange, wrapping and unwrapping BETH are at zero fee.
  \item The conversion rate is updated once daily at 00:00 UTC, and staking and redemption through the exchange are suspended each day from 23:45 to 00:15 UTC while it is applied.
  \item The rate was initialised at 1:1 with BETH at launch on 2023-04-27 08:00 UTC, and the same token address serves Ethereum and BNB Smart Chain.
  \item Slashing exposure, reward computation and the reserve sit entirely inside the custodian; a loss reaches the token only if a whitelisted caller writes a smaller number, and may not be written at all.
\end{itemize}
\end{remark}

\begin{remark}
The witness is built from verified bytecode and nothing else --- there is no repository and no specification --- and on that evidence alone the corpus record is overturned rather than refined: the rate is not incremented by an admin key but written through a verified rate-limiting contract that bounds each caller's move per interval. The access element is dropped, because the payable deposit requires only a non-zero value and the login the record describes belongs to the exchange and not to the contract; the share element is dropped, because the rate is written independently of supply, and carrying it beside the index would count one claim twice under its dual descriptions; and a queue, a freeze and a pause are added, because the unwrap contract records requests against a minimum lock, the claim path is blacklistable and both contracts are pausable, none of which the record carries. Attestation is deliberately not claimed: the published reserve methodology names no attester, states no frequency and does not enumerate this asset, so its warrant is unmet, and the unmet warrant is the finding. Loss allocation is deliberately not claimed either, because a loss here is not a rule the ledger executes but a smaller number a whitelisted caller may choose to write, at a size each write's refilling allowance caps.
\end{remark}

\subsection{EigenLayer}\label{sub:cat:lsd:eigenlayer}

A depositor sends a fungible token into a whitelisted strategy, or points a validator's withdrawal credentials at a pod it deploys itself and proves them on chain, and receives deposit shares --- an internal, non-transferable accounting claim rather than a token --- which it then delegates to an operator, transferring no custody and no transaction authority but only slashing liability. Nothing accrues to the share: services fund reward submissions, and earners claim against posted distribution roots. What the collateral secures is an arbitrary external service whose duties and faults the core protocol does not model; that service authors its own slashing predicate and calls for a proportion, and the protocol validates the arithmetic of a slash and never its ground. The same deposit backs many services at once and no unit of slashability backs more than one, because an operator's allocations cannot exceed a conserved total that only slashing reduces. Exit is a queue behind a long delay, with the beacon-chain exit queue sitting underneath it for native stake and modelled nowhere in these contracts.

\begin{measurement}\label{meas:cat:lsd:eigenlayer}
Take $X = \{Aw,Bs,Ep,Fd,Gp,Rs,Sh,Sl,Up,Wq,Xm\}$, with canonical form
$\mathrm{ex}(X) = \{Aw,Bs,Ep,Fd,Gp,Rs,Sh,Sl,Up,Wq,Xm\}$, every element being primitive in it.
It is admissible: no requirement term is open, every element is warranted, and it arms no prohibition. It discharges 12 of the 21 recorded obligations; 9 are residue. No composite of at most three corpus protocols equals it; no corpus protocol is contained in it.
It is not minimal: \{Bs,Sl,Wq,Fd,Up\} can be removed with every obligation still covered and admissibility intact.
\end{measurement}

\begin{remark}
The construction covers the accounting and misses the protocol. What is allocated here is not principal but \emph{exposure} over principal that does not move, and the conserved budget that makes the allocation exclusive is the central object of the design; the vocabulary has a symbol for slashable capital and none for a budget over that capital's slashability. Nor can it record who is permitted to define the fault --- the axis on which this protocol and Babylon sit at opposite ends, an open predicate at one end and no third-party predicate at all at the other, and both ends read as an empty array. The delays are the same story: the vocabulary names an epoch and a queue, and cannot say that a delay exists so that a counterparty who has already been left behind keeps its recourse. What is left is a loss recorded before it is effected and then cleared by any passer-by, and a remedy paid to a named recipient at an address fixed for the life of the set --- neither an allocation nor a pause, and neither of them nameable.

The obligations no element discharges are these.
\begin{itemize}\setlength{\itemsep}{0pt}
  \item [Vl sub-mechanism 1 of 5: validator key registration and vetting] A depositor may instead deploy an EigenPod at a CREATE2 address salted on its own staker address, deposit 32 ETH with the pod as withdrawal credentials, and then prove on chain that a named validator's withdrawal credentials point at that pod; validator state is thereafter maintained by periodic checkpoints.
  \item A staker delegates its entire position to an operator; the delegation transfers no custody, no assets and no transaction authority, only the right to be counted as that operator's weight and the corresponding slashing liability.
  \item [Vl sub-mechanism 2 of 5: stake allocation and scheduling] An operator allocates a magnitude, a proportion of its delegated stake in a given strategy, to a named operator set; allocations are held per (operator, strategy, operator set) with a current magnitude, a pending diff and an effect block, and deallocations are processed through a per-(operator, strategy) queue in order.
  \item [Whether the same collateral can be committed twice] The sum of an operator's allocated magnitudes across all operator sets can never exceed INITIAL\_TOTAL\_MAGNITUDE (1e18, one WAD), which is reduced only by slashing, so the same *deposit* backs many operator sets but each unit of *slashability* is exclusively allocated and no two operator sets can slash the same stake.
  \item [Who may write a slashing condition] Any AVS may create an operator set, define its slashable strategies and its membership roster, and call slashOperator with a proportion in WAD; the core protocol validates the arithmetic and never the ground for the slash. ELIP-006 states the consequence plainly: an attacker who gains access to AVS keys can drain the entirety of operator and staker allocated stake.
  \item Slashed shares are first marked burnable-or-redistributable, and only after SLASH\_RESOLUTION\_DELAY\_BLOCKS may anyone permissionlessly clear them, sending them either to the burn address 0x00000000000000000000000000000000000E16E4 or to the operator set's redistribution recipient.
  \item Native ETH and EIGEN are excluded from redistribution: native ETH because paying a recipient would require beacon exits whose queue can range from hours to weeks, and EIGEN because it requires a delay on token protocol outflow for its intersubjective machinery.
  \item For faults that are obvious to human observers but not provable in the EVM, the designed fallback is forking the EIGEN token: a challenger creates a fork in which the malicious stakers are slashed, and AVSs and users choose which fork to respect.
  \item [Vl sub-mechanism 3 of 5: deposit front-running defence] There is none, and none is needed: the pod is deployed at an address deterministic in the staker's own address, and the same party that supplies the 32 ETH also owns the withdrawal credentials it is deposited against, so the adversary Lido's DepositSecurityModule exists to defeat, an operator capturing a pool's ETH by pre-registering the pubkey against different credentials, has no counterpart here.
\end{itemize}
\end{remark}

\begin{remark}
The corpus record's central residue claim for this protocol --- that the same collateral is simultaneously bonded to mutually unaware slashing laws --- is refuted rather than refined, and by the protocol's own arithmetic: the governing improvement proposal and the allocation contract's documentation both state that the sum of an operator's magnitudes can never exceed its total, which is precisely the property that stops distinct operator sets slashing the same stake. Protocol-funded emission is dropped and a fee-and-distribution element added in its place, because these rewards are funded by the service that submits them and reach earners as a distribution of received value, and the record carries the emission symbol with no supporting marker. The governance-timelock element is dropped under cite-or-drop: no first-party source names one, which leaves the upgrade warrant met only by a scoped pause, and that is an absence of evidence rather than evidence of absence. Delegated execution authority is deliberately not claimed for delegation, because what moves is liability and not the right to act; and the index element is deliberately not claimed for the pod's scaling factors, because they are per-position and fall in one direction only, and flattening that asymmetry would erase the design.
\end{remark}

\subsection{ether.fi}\label{sub:cat:lsd:etherfi}

A user deposits ETH into the liquidity pool and receives eETH, a share of a ledger whose nominal balance moves with the pool's reported value; weETH is the non-rebasing wrapper, and it is the form every integration and every cross-chain deployment holds. The same principal is staked and restaked at once, because the node contract holding the validators' withdrawal credentials is itself a pod; which services that stake secures, and therefore who is empowered to author the conditions under which it is destroyed, is settled by the protocol in contracts outside this repository. Operators do not select themselves: they bid on chain for the right to be assigned a validator, and the winning bid is forwarded to the treasury. Exit is a choice among doors --- instant for a fee, metered by a refilling bucket and floored on pool value; a standard queue whose claim rate is frozen at an administrator's finalisation and thereafter invariant under rebases; and a priority queue for whitelisted addresses under a minimum delay. Nothing sits beneath the holder: operators post no bond, so a consensus-layer or restaking loss reaches eETH as a smaller rebase and stops there.

\begin{measurement}\label{meas:cat:lsd:etherfi}
Take $X = \{Aw,Ep,Ex,Fd,Fz,Gp,Ix,Rb,Rd,Rs,Sh,Sl,Tg,Up,Wq\}$, with canonical form
$\mathrm{ex}(X) = \{Aw,Ep,Ex,Fd,Fz,Gp,Ix,Rb,Rd,Rs,Sh,Sl,Tg,Up,Wq\}$, every element being primitive in it.
It is admissible: no requirement term is open, every element is warranted, and it arms no prohibition. It discharges 14 of the 25 recorded obligations; 11 are residue. No composite of at most three corpus protocols equals it; 1 corpus protocol is contained in it, and together they supply every element except \{Ep,Ex,Fd,Fz,Gp,Rb,Rs,Sl,Tg,Wq\}.
It is not minimal: \{Ix,Ex,Ep,Rd,Sl,Tg,Up,Gp\} can be removed with every obligation still covered and admissibility intact.
\end{measurement}

\begin{remark}
The protocol's own file layout is the argument: registration, the market for the slot, the key and predeposit path, exit and consolidation, and the restaking leg are separate contracts, and the candidate symbol for this category would collapse every one of them into itself. The market is the sharpest --- a priced, cleared, on-chain primary auction for the right to operate is not an emission, not a fee and not a permission gate, and nothing names it. Beneath it sit a refilling budget metering a state transition, which recurs across this category and has no symbol; a claim severed from the accounting index at a named block, which is the difference between a queue position and a fixed-price forward; and a menu of exits over one claim, where the routing is itself the price of leaving. The same principal is committed to consensus and to an external security market at once: within that market a conserved budget bounds the second commitment, and across markets nothing here bounds it or discloses the aggregate. The party that bears the loss neither chooses nor can observe on chain who may write the conditions under which its principal is destroyed, and that is the thinnest point in the carrier for this application.

The obligations no element discharges are these.
\begin{itemize}\setlength{\itemsep}{0pt}
  \item A single rebase may not raise the pool's total value by more than a hard cap expressed in basis points.
  \item Node operators do not self-select: they bid on chain for the right to be assigned a validator, with bids bounded below and above by minBidAmount and maxBidAmount and a distinct whitelistBidAmount for whitelisted bidders, and the winning bid's revenue is forwarded to the treasury.
  \item [Vl sub-mechanism 3 of 5: deposit front-running defence] A 1 ETH INITIAL\_DEPOSIT\_AMOUNT is deposited first, pinning the validator's withdrawal credentials to the EtherFiNode contract, before the remaining balance is committed; validators range from MIN\_VALIDATOR\_SIZE\_WEI of 32 ETH to MAX\_VALIDATOR\_SIZE\_WEI of 2048 ETH.
  \item [Vl sub-mechanism 4 of 5: exit signalling] Exit and consolidation requests are issued from EtherFiNodesManager against a validator pubkey hash, with FULL\_EXIT\_GWEI of 2,048,000,000,000, and each of unrestaking, exit requests and consolidation requests is metered by its own named rate-limit identifier.
  \item The choice among three exit doors, differently priced and carrying different rights, is itself a mechanism: instant is fee-bearing and bucket-limited, standard is free and rate-frozen at an admin's finalisation, priority is whitelisted and delay-bounded.
  \item A refilling token-bucket budget meters unrestaking, exit requests, consolidation requests and instant redemptions.
  \item [Who may write a slashing condition over the restaked collateral] Neither ether.fi's contracts nor the eETH holder: AVS selection and registration live outside this repository, in etherfi-avs-operator, and the README states restaking and AVS management are strictly controlled by the Protocol.
  \item [Whether the same collateral can be committed twice] Yes: the same principal is committed to Ethereum consensus and to EigenLayer operator sets at once. Within EigenLayer the second commitment is bounded, because an operator's magnitude allocations cannot exceed its total magnitude; across venues nothing in ether.fi bounds or discloses the aggregate exposure.
  \item A node contract may execute a bounded set of arbitrary calls into its own EigenPod or into a named external target, gated per caller, per function selector and per target address.
  \item Deposits were acquired through an off-chain loyalty-points programme whose balances were later converted into a token.
  \item Some validators are operated as distributed validators, one key split into keyshares across a cluster with a signing threshold, in the solo-staker programme.
\end{itemize}
\end{remark}

\begin{remark}
The tranche element is dropped: the tiered note structure the corpus record describes as a bond in tranche language now sits in the repository's archive directory, retained for storage-layout continuity, so a decomposition carrying it decomposes a dead mechanism. The bond element goes with it, because the whitepaper states that permissioned operators are required to bid and not to post collateral and that solo stakers post none either, and no other first-loss layer could be identified --- which leaves loss allocation doing the whole of the work, and the absence of any junior layer beneath the holder as the result. A permission element, a freeze element, a direct-redemption element and an index element are all added, because the record carries no access or freeze symbol for this protocol at all, the instant door is a genuine priced redemption right, and the wrapper is the adapter that keeps a rebasing claim admissible in balance-invariant ledgers. The record's multi-venue restaking framing is not inherited: within a single venue it is refuted by that venue's conserved allocation budget, and across venues the restaking contract imports the interfaces of one venue only, so the framing is unknown rather than established.
\end{remark}

\subsection{ApeX}\label{sub:cat:perp:apex}

A venue whose order book is run off-chain by an operator and settled on-chain against validity proofs. The trader posts margin to a bridge contract, trades against a book they cannot observe being constructed, and relies on the proof system for settlement integrity rather than on the sequencing being agreed. The venue never states where its matching engine runs; the stack it is built on does.

\begin{measurement}\label{meas:cat:perp:apex}
Take $X = \{Ad,Au,Ct,Ex,Fd,Li,Ob,Pf,Sh,Xf,Xm\}$, with canonical form
$\mathrm{ex}(X) = \{Ad,Au,Fd,Ob,Pf,Sh,Xf,Xm\}$; the elements \{Ct,Ex,Li\} are derived rather than chosen.
It is admissible: no requirement term is open, every element is warranted, and it arms no prohibition. It discharges 12 of the 22 recorded obligations; 10 are residue. No composite of at most three corpus protocols equals it; no corpus protocol is contained in it.
It is not minimal: \{Xf,Xm\} can be removed with every obligation still covered and admissibility intact.
\end{measurement}

\begin{remark}
The operator's discretion is the whole of this venue's risk and none of its decomposition. Two facts found on the second pass make the point sharper than any absence could. The published escape hatch does not work: on four of the five chains for which the venue publishes it, the exit call reverts, because the deployed verifier returns false unconditionally. And the exit deadline is counted in blocks rather than in time, so the same nominal window is a fortnight on one chain and the better part of a year on another. A venue with a documented escape hatch, a venue with a working one, and a venue with neither are three different objects; the vocabulary writes all three the same way.

The obligations no element discharges are these.
\begin{itemize}\setlength{\itemsep}{0pt}
  \item the venue does not state where its matching engine runs: the documented architecture is an Access / Application / Blockchain layering whose execution layer 'leverages the zkLink engine' and whose settlement layer 'provides proof verification', reached through an operator-hosted REST and WebSocket API
  \item no constraint on the operator's ordering or matching discretion is published: there is no stated fairness rule, no sequencing commitment and no named sequencer, and zkLink X's own sequencing is described as delegable to third parties without ApeX stating whether it delegates
  \item state transitions are finalised by validity proofs on a zkLink X app-rollup, batched at app-rollup and L2 levels before Ethereum settlement, with proofs first committed on a host chain that forwards sync hashes
  \item the data-availability layer under the venue is not stated: zkLink X offers Ethereum, Celestia, EigenDA, Polygon Avail or a zkLink-organised committee, and ApeX's selection is documented nowhere
  \item a censored trader exits in two steps — submit a forced exit on zkLink L2, wait for on-chain confirmation, then claim on L1 by interacting directly with the ApeX zkLink contract on the chosen chain — with no deadline or exit-window duration stated anywhere
  \item one account balance spans six chains: the chain of deposit is a coordinate on a single account object rather than a separate account per domain
  \item eight non-USDT collaterals are valued as Amount x Index Price x Collateral Ratio under tiered haircuts and per-asset supply caps (USDC 99\% on the first 5,000,000 then 97.5\%, capped at 20,000,000; USDT 100\% and uncapped), while a second balance — Perpetual Balance — is computed without the haircut
  \item no insurance fund exists: the string does not occur anywhere in the venue's 383 KB documentation corpus, and no capital is documented as absorbing a bankrupt position before auto-deleveraging
  \item the venue renders a third party's order books — Polymarket's — inside its own interface, margined from the user's Omni account and settled in Polymarket USD
  \item no upgrade authority, timelock, guardian or pause power over the contracts custodying user funds is documented, and the contracts in question are zkLink's rather than the venue's
\end{itemize}
\end{remark}

\begin{remark}
The witness drops the staked backstop the corpus assigns, because no insurance fund is documented anywhere in the venue's own materials, and an element carried on an undocumented mechanism is a forced fit. The remaining terms are discharged by elements the venue plainly has.
\end{remark}

\subsection{Aster}\label{sub:cat:perp:aster}

A venue that has moved its order books into the consensus of its own proof-of-stake-authority chain, with a closed validator set and a node that is not open source. The corpus records it as an operator-run book, which its current architecture contradicts.

\begin{measurement}\label{meas:cat:perp:aster}
Take $X = \{Ad,Bs,Ct,Em,Ex,Li,Ob,Pf,Sb,Sd,Sh,Sv,Tp,Vl,Xf,Xm\}$, with canonical form
$\mathrm{ex}(X) = \{Ad,Bs,Em,Ob,Pf,Sb,Sd,Sh,Sv,Tp,Vl,Xf,Xm\}$; the elements \{Ct,Ex,Li\} are derived rather than chosen.
It is admissible: no requirement term is open, every element is warranted, and it arms no prohibition. It discharges 15 of the 26 recorded obligations; 11 are residue. No composite of at most three corpus protocols equals it; no corpus protocol is contained in it.
It is not minimal: \{Tp,Vl,Em,Xf,Xm\} can be removed with every obligation still covered and admissibility intact.
\end{measurement}

\begin{remark}
A closed validator set with a closed-source node is the sharpest case in the category: the trust model is entirely a question about who the validators are and what they may do, and the carrier has no party sort in which to ask it. The second pass sharpened it further. The published validator set is not the set that matters --- withdrawals clear on a signature threshold far smaller than it, over a balance that is not small --- and there is neither an escape hatch nor slashing, both confirmed as negatives rather than assumed. The element set records a book, a funding transfer and a margin test, and none of that.

The obligations no element discharges are these.
\begin{itemize}\setlength{\itemsep}{0pt}
  \item all order matching and trading logic are processed and stored on-chain, inside Aster Chain — the venue's own Layer 1, launched 16 March 2026, running Proof-of-Staked-Authority consensus at 50 ms blocks with zero gas
  \item external validator participation is not supported in Phase 1, and core chain contracts and RPC infrastructure are not open-sourced, as stated policy, in order to protect proprietary matching engine algorithms
  \item no validity or fraud proof covers execution; safety rests on a permissioned PoSA validator set, while zero-knowledge verification is applied at the order layer — every hidden order requires ZK verification whether or not Account Privacy is enabled
  \item data storage is hybrid: key trading and account events are recorded on chain while certain operations remain off chain for privacy protection and system efficiency
  \item there is no forced-inclusion path, censorship deadline or escape hatch: withdrawals are batched and paid on 2-of-3 validator confirmation, typically in about 30 seconds, with a fee covering validator operating costs
  \item margin may be posted in yield-bearing instruments — asBNB, USDF, asUSDF, asBTC and asCAKE — each of which is a claim on another protocol's strategy with its own minter contract, its own peg and its own auditor
  \item on Shield Mode and 1001x the trader's counterparty is the ALP pool at an oracle price, with no order book and no inventory-repricing curve; risk is managed by fees, caps and a liquidation loss rate rather than by a curve
  \item Shield Mode liquidation triggers on last price against a liquidation price derived from margin, leverage, entry, funding and a liquidation loss rate that varies by pair and leverage band (BTC/ETH 75\% at 1-200x, 70\% at 201-500x, 65\% at 501-1001x), and the 1001x venue charges no opening fee, permits no margin top-up, charges a dynamic PnL-based closing fee and caps net ROI at roughly 300\%
  \item a hidden order reveals neither size nor presence to the public book while sharing the book's liquidity, always carries ZK verification of order data whether or not Account Privacy is on, and produces hidden positions excluded from all public data and portfolio analytics; a fresh one-time stealth address is generated per trade so that activity cannot be correlated
  \item listing is decided by validator vote: any active validator holding at least 20M staked ASTER may propose, and passage requires a double majority of more than 50\% of stake and more than 50\% of validators, formalised as AOS-1 for spot listing since 25 June 2026
  \item no timelock, guardian, pause authority or upgrade path for the chain or its contracts is documented
\end{itemize}
\end{remark}

\begin{remark}
This construction departs furthest from the corpus. The corpus decomposition leaves the loss-absorption term of the perpetual-funding requirement open and is rejected for it; the witness supplies the elements the venue's current design actually carries, and the rejection disappears. The departure is a correction of the record, not a repair of the protocol.
\end{remark}

\subsection{edgeX}\label{sub:cat:perp:edgex}

A venue in the middle of a migration, and the direction is the opposite of the usual one. Its validity-proof stack still custodies almost all of the value while the newer chain holds a rounding error, and the newer chain's settlement rests on interactive fraud proofs with a small data-availability committee. On the evidence the current regime is weaker than the one it replaces.

\begin{measurement}\label{meas:cat:perp:edgex}
Take $X = \{Ct,Ex,Li,Ob,Pf,Up,Xf\}$, with canonical form
$\mathrm{ex}(X) = \{Ob,Pf,Up,Xf\}$; the elements \{Ct,Ex,Li\} are derived rather than chosen.
It is not admissible: it leaves the requirement term $L4\!:\!Ad{\mid}Sl{\mid}Bs$ open. It discharges 9 of the 18 recorded obligations; 9 are residue. No composite of at most three corpus protocols equals it; no corpus protocol is contained in it.
\end{measurement}

\begin{remark}
The rejection is a result about the tables and not a criticism of the venue: the requirement demands a named loss-absorption mechanism, and this venue documents none. An undocumented mechanism and an absent one are different things, and the vocabulary cannot tell them apart --- the same defect the options category exhibits from the other side. Nor can it record that a venue has moved from validity proofs to fraud proofs, which is the single most consequential fact about it.

The obligations no element discharges are these.
\begin{itemize}\setlength{\itemsep}{0pt}
  \item matching runs in an operator-run engine and only state roots are anchored to Ethereum, with the operator as sole block proposer
  \item the venue has occupied two different loci while keeping the product, the account model and the API shape: V1 was 'powered by the StarkEx Layer 2 engine' as a validium, and V2 is a modular rollup that has 'temporarily adopted the Arbitrum technology stack as a transitional framework', with V3 announced as a sovereign AppChain
  \item no constraint on the operator's ordering, matching or inclusion discretion is published for either generation: there is no fairness rule, no sequencing commitment and no forced-ordering property, and the Cairo program that executed V1's state transitions is marked 'Code unknown'
  \item V1 batched trades through the StarkEx service to the SHARP shared prover, whose STARK proofs were verified on Ethereum before the StarkPerpetual contract would accept a state root; for V2 the documentation never states whether there are validity proofs, fault proofs, or neither
  \item V1 is a validium: proof construction relies fully on data that is not published on chain, held by a data-availability committee with a threshold of 2 of 6 whose members are not publicly known and who have no on-chain assets at risk of being slashed in a data-withholding attack; V2's data-availability regime is not stated
  \item on V1 a user may force the sequencer to include a trade or a withdrawal by submitting the request through L1, and if it is censored or down for 7 days anyone may freeze the system and withdraw by Merkle proof — but a forced TRADE requires the user to find a counterparty by out-of-system means; V2 asserts trustless permissionless withdrawals at any time even under operator sequence failure while specifying no mechanism, no deadline and no contract
  \item a deposit from a non-Ethereum chain is credited by transferring an equivalent amount from the venue's own asset pool on that chain, with rebalancing latency of up to 2 hours for single deposits over \$300K
  \item a user may onboard by email, in which case an MPC wallet is created for them and a fresh address is funded over Arbitrum
  \item what absorbs a bankrupt position is not documented: the strings 'insurance fund', 'auto-deleveraging', 'bankruptcy price', 'socialised loss' and any liquidation-fee schedule are absent from the venue's entire 414 KB documentation corpus
\end{itemize}
\end{remark}

\begin{remark}
The witness is deliberately thin. The corpus carries a staked backstop, a cross-domain verification element, a share claim and emissions; none is supported at a primary source for the current system, so all four are dropped. What remains is what can be evidenced, and the construction is rejected because of it.
\end{remark}

\subsection{Hyperliquid}\label{sub:cat:perp:hyperliquid}

The largest venue in the category matches orders in a book whose state is part of its own chain's consensus, so the sequence of trades is agreed by the validator set rather than asserted by an operator. Margin is posted on the same chain, funding is a long/short transfer, and a position that fails its margin test is closed by a liquidator and, past that, by auto-deleveraging against profitable accounts. Its backstop is a depositor-funded vault rather than staked capital, which the construction records as loss allocation over a pro-rata claim rather than as a bond.

\begin{measurement}\label{meas:cat:perp:hyperliquid}
Take $X = \{Ad,Ct,Ep,Ex,Gp,Li,Ob,Pf,Sh,Sl,Tp,Vl,Wq,Xf,Xm\}$, with canonical form
$\mathrm{ex}(X) = \{Ad,Ep,Gp,Ob,Pf,Sh,Sl,Tp,Vl,Wq,Xf,Xm\}$; the elements \{Ct,Ex,Li\} are derived rather than chosen.
It is admissible: no requirement term is open, every element is warranted, and it arms no prohibition. It discharges 15 of the 25 recorded obligations; 10 are residue. No composite of at most three corpus protocols equals it; no corpus protocol is contained in it.
It is not minimal: \{Tp,Sl,Ep,Wq\} can be removed with every obligation still covered and admissibility intact.
\end{measurement}

\begin{remark}
The residue is not incidental to this venue; it is what distinguishes it from the other four. That the book is in consensus, that no operator may reorder it, and that a user can exit without the operator's cooperation are three separate facts, and the vocabulary states none of them.

The obligations no element discharges are these.
\begin{itemize}\setlength{\itemsep}{0pt}
  \item the book, the margin state, the matching and the liquidations are the chain's own state machine — HyperCore is executed by the validator set under HyperBFT consensus, not by an off-chain operator
  \item within a block, actions are sorted into three classes — actions containing no GTC or IOC order, then cancels, then actions carrying at least one GTC or IOC — with proposed order preserved inside each class
  \item the correctness of every fill, funding payment and liquidation rests on two-thirds of staking power being honest; there is no validity proof, no fraud proof, and the node that computes the state is distributed as a signed pre-compiled binary with no published source
  \item the full order-book and account state is replicated to every validator and served from public nodes and the public API; nothing is posted to, or reconstructible from, any other chain
  \item a trader has no unilateral exit: there is no forced-inclusion transaction, no censorship deadline and no escape hatch, so leaving the venue requires the validator set to sign
  \item collateral is pooled across every perp and spot market by default, and a trader may instead ring-fence a single position in isolated margin, or partition capital into sub-accounts and vaults
  \item one pool, HLP, simultaneously makes markets through multiple strategies, performs backstop liquidations, supplies USDC in Earn and accrues a share of trading fees, and it is community-owned rather than protocol-owned
  \item the right to list a spot ticker is sold by a 31-hour Dutch auction decaying to 500 HYPE, and the right to deploy a perp DEX requires 500k HYPE staked and maintained for at least 183 days, slashable by stake-weighted validator vote up to 100\% for invalid state transitions and up to 50\% or 20\% for downtime and degradation
  \item a third-party front end can attach its own fee to order flow it originates, up to 0.1\% on perps and 1\% on spot, charged inside the exchange's fee logic, requiring an ApproveBuilderFee action from the user's main wallet, at most 10 active approvals and at least 100 USDC in the builder's perps account
  \item orders must be placed within 1\% of the oracle price, with HLP exempt so it can keep quoting, and open-interest caps scaled to liquidity block new positions when hit
\end{itemize}
\end{remark}

\begin{remark}
The construction is exhibited rather than derived: the corpus decomposition carries a staked backstop this venue does not have --- the vault is unbonded and unslashable depositor capital --- so the witness drops it and lets socialised loss over a pro-rata claim carry the obligation instead. Among the alternatives left open by the loss-absorption term, that is the one the code supports.
\end{remark}

\subsection{Lighter}\label{sub:cat:perp:lighter}

The only venue in the category whose deployed system can be checked against its published source: the circuits are open and the verifier is regenerable, which an independent reviewer has confirmed. Orders are matched by a sequencer whose output is proven, and a censorship deadline gives the trader a unilateral exit.

\begin{measurement}\label{meas:cat:perp:lighter}
Take $X = \{Ad,Aw,Ct,Em,Ex,Fd,Gp,Li,Ob,Pf,Rf,Sh,Sl,Tg,Tp,Up,Xf,Xm\}$, with canonical form
$\mathrm{ex}(X) = \{Ad,Aw,Em,Fd,Gp,Ob,Pf,Rf,Sh,Sl,Tg,Tp,Up,Xf,Xm\}$; the elements \{Ct,Ex,Li\} are derived rather than chosen.
It is admissible: no requirement term is open, every element is warranted, and it arms no prohibition. It discharges 17 of the 24 recorded obligations; 7 are residue. No composite of at most three corpus protocols equals it; 2 corpus protocols are contained in it, and together they supply every element except \{Ad,Ct,Em,Ex,Fd,Li,Ob,Pf,Rf,Sh,Sl,Tp\}.
It is not minimal: \{Pf,Tp,Sl,Sh,Em,Fd,Tg,Up\} can be removed with every obligation still covered and admissibility intact.
\end{measurement}

\begin{remark}
That this venue is the only one whose deployment can be checked against its source is the most important fact about it, and it is invisible: no element distinguishes a proven sequencer from an asserted one.

The obligations no element discharges are these.
\begin{itemize}\setlength{\itemsep}{0pt}
  \item price-time priority is enforced by the address space rather than by a rule: an ask at price p with nonce o occupies leaf index p*2\textasciicircum{}O + o and a bid occupies p*2\textasciicircum{}O + (2\textasciicircum{}O - 1 - o), so index order IS priority order, and every insertion proves non-crossing and every trade proves the supplied maker was the highest-priority crossing order by requiring crossingSize = 0 along the root-to-leaf path
  \item the book is maintained by an off-chain sequencer whose only power is ordering — since execution and matching are verified by the prover, every valid state transition is deterministically inferable from the transaction order and the oracle data
  \item batches are proven by Plonky2 circuits — TurboPlonk arithmetisation over Goldilocks with FRI — aggregated block to segment to batch, wrapped into a gnark PLONK proof over BN254 for on-chain verification against the Aztec Ignition setup, and the deployed ZkLighterVerifier can be regenerated from the published circuits
  \item data availability is defined by a sufficiency predicate rather than by a venue: an Account Delta Tree initialised empty each batch receives every update also applied to the Public Account Tree, is serialised, compressed and posted to Ethereum in a blob-carrying transaction with the market data, so replaying the blobs in order reconstructs the account state exactly — and two circuits bind the claim, a Blob Polynomial Equivalence Circuit evaluated at the point Ethereum's precompile uses and Account Delta Traversal Circuits at the same Fiat-Shamir point
  \item a censored trader can force inclusion by queueing a priority transaction directly through the Ethereum contracts — deposits, withdrawals, order creation, order cancellation and burning of pool shares are all forceable — and if the sequencer misses the 14-day deadline, anyone may trigger desert mode, after which commitments stop, the state freezes, and users withdraw against the frozen root using blob data alone, with open positions settled at the latest mark price
  \item standard accounts pay zero maker and taker fees at 300 ms taker and 200 ms maker latency, Premium accounts pay 0.4 to 2.8 bps scaling down with staked LIT in exchange for lower latency, and Plus accounts pay a flat 0.5 bps
  \item a second sovereign deployment of the same stack, Lighter on Robinhood, runs with USDG as base collateral and SPY and USO as margin assets, under its own verifier contract
\end{itemize}
\end{remark}

\begin{remark}
The witness adds elements the corpus omits --- a quote mechanism, a timelock, a loss-allocation rule and a permission gate --- each supported at the venue's own source. It also declines to carry a withdrawal queue: what the venue documents is a censorship deadline, which bounds the operator rather than rationing redemption, and the queue element names the second thing.
\end{remark}

\subsection{CIAN Yield Layer}\label{sub:cat:yld:cian}

A depositor brings the accounting asset, or one of a set of related staking tokens through a converter path, and receives a transferable pro-rata share of the vault. The yield is a recursive leveraged loop in an external lending market: supply the staking token, enter its correlated-asset mode, borrow the accounting asset against it, re-stake, repeat. The leverage, the venue, the swap route, the flash-loan source and the moment are all chosen by an off-chain agent holding a cloud-hosted key, which reaches the strategy through a contract whose sole privileged entry point forwards arbitrary batched calls. Exit is not the standard redemption path, which reverts outright: a depositor files a request naming a payout token and waits for a designated operator to settle it, at an amount that operator decides. The vault also refuses to report its own assets, and therefore refuses deposits and redemptions, if that agent has not lately written the price by which every share is valued.

\begin{measurement}\label{meas:cat:yld:cian}
Take $X = \{Aw,Ct,Ex,Fl,Gp,Ix,Sh,Up,Wq\}$, with canonical form
$\mathrm{ex}(X) = \{Aw,Ct,Ex,Fl,Gp,Ix,Sh,Up,Wq\}$, every element being primitive in it.
It is admissible: no requirement term is open, every element is warranted, and it arms no prohibition. It discharges 9 of the 19 recorded obligations; 10 are residue. No composite of at most three corpus protocols equals it; no corpus protocol is contained in it.
\end{measurement}

\begin{remark}
This is the construction on which the category's central claim turns, and it turns on what is absent rather than on what is present. Written out, the contract holds an envelope --- a ceiling on the ratio, a fraction cap per strategy, a deposit cap, a slippage floor, a registered set of destinations, a pause --- and a human supplies the point inside it. Nothing computes the target, schedules a rebalance, defines when a target has been missed, or compels the agent to act at all, and the one function that computes the amount needed to reach the ceiling is a helper the agent is free to ignore. So the strategy is not an element the vocabulary failed to name; it is not on the chain to be named, which is a difference in \emph{kind} rather than in coverage, and it predicts correctly that no further reading of this protocol's bytecode would ever produce it. The shape reaches past one vault: among this protocol's own destinations is a vault curated by another agent under another mandate, and the largest object in this category by capital is such a curator --- the one the requirement tables reject, whose distinguishing feature has no symbol at all.

The obligations no element discharges are these.
\begin{itemize}\setlength{\itemsep}{0pt}
  \item that index is written by an off-chain operator at its own discretion, with no oracle input, no attester, no staleness recourse beyond a revert, and no derivation from realised profit and loss
  \item the vault refuses to report its own assets — and therefore refuses deposits and redemptions — if the index has not been written within three days
  \item total deposits into the vault are capped at a fixed quantity of the accounting asset
  \item the vault's assets are deployed as a recursive leveraged loop in an external lending market: supply a staking token, enter its correlated-asset mode, borrow the accounting asset against it, re-stake, repeat
  \item the ceiling the vault enforces is a self-imposed inner bound below an immutable maximum, and the position is deliberately kept pressed against it
  \item that agent supplies the leverage amount, the swap route, the slippage floor and the flash-loan source as call arguments; nothing in the contracts computes them, schedules a rebalance, defines when a target has been missed, or compels the agent to act at all
  \item capital may be forwarded only to strategy contracts the owner has registered with the vault
  \item no single strategy may hold more than a set fraction of the vault's total assets, checked at the moment of transfer
  \item the amount a request settles for is decided by the entity that confirms it
  \item eight per cent of gains is taken by the operator as a performance fee, under a hard ceiling of fifteen per cent, with the documented exit and management fees currently set to zero
\end{itemize}
\end{remark}

\begin{remark}
The witness adds a permission gate and drops emissions. The gate is the crux of the decomposition: the entry points that lever and unlever the position are callable by one stored address and no other, and that address resolves to a contract whose only privileged function forwards batched calls from a cloud-hosted signing key, so the whole authority path from key to position is a single enumerated caller. Emissions go because there is nothing to emit --- the documentation still reads that a token is coming, the listing carries no symbol, and what is actually paid is an off-chain point balance. The threshold test and the external feed stay, and among the truth-source alternatives the feed is the one the code supports, but both are inherited rather than owned: the price comes from the lending market's own oracle, and the ceiling the vault checks is a self-imposed inner bound below that venue's immutable maximum, which the vocabulary cannot tell apart from merely being subject to a test.
\end{remark}

\subsection{Convex Finance}\label{sub:cat:yld:convex}

Deposits arrive by separate paths into separate contracts. A liquidity-pool token is staked into the corresponding upstream gauge and credited one-for-one into a per-pool reward contract, withdrawable at any time and without fee. The upstream's governance token, once deposited, joins a single permanent escrow and comes back as a transferable wrapper --- a claim with no redemption function, whose only exit is sale --- while the protocol's own token can be locked for a fixed minimum term to obtain the weight that decides how that escrow votes. None of the yield originates here: it is the upstream's emissions and the upstream's admin fees, taken at a pooled boost no small depositor could reach alone, and shared out after a levy. Nobody chooses where a depositor's capital sits --- the depositor names the pool and the protocol stakes it there --- and the discretion the protocol does hold is over votes cast inside protocols it does not own.

\begin{measurement}\label{meas:cat:yld:convex}
Take $X = \{Em,Ep,Fd,Gp,Ix,Rd,Sh,Sr\}$, with canonical form
$\mathrm{ex}(X) = \{Em,Ep,Fd,Gp,Ix,Rd,Sh,Sr\}$, every element being primitive in it.
It is admissible: no requirement term is open, every element is warranted, and it arms no prohibition. It discharges 8 of the 18 recorded obligations; 10 are residue. No composite of at most three corpus protocols equals it; no corpus protocol is contained in it.
It is not minimal: \{Sh,Rd\} can be removed with every obligation still covered and admissibility intact.
\end{measurement}

\begin{remark}
Strip the vote layer out of this construction and what remains is indistinguishable from any staking-rewards contract, which is the measure of how much of this protocol the vocabulary cannot see. The mechanism that makes the whole arrangement rational --- a reward rate on one balance multiplied by a different party's escrowed holding of a different token --- has no symbol, and neither does the escrow, nor the weight that decays with the time left on it, nor the one-way valve that turns a redeemable position into a claim with a market and no redemption. Nor does the acquisition and exercise of governance rights inside protocols the holder does not control, which is a containment relation between protocols in a model that composes them only as peers. One item of residue here is of a different kind, and it is a finding about the vocabulary rather than about this protocol: the construction must record that no deployed code can be replaced, and the only way it can record that is as the absence of a mutable proxy. The vocabulary can name mutability and cannot \emph{assert} immutability, so a protocol whose control plane is a commitment never to change is written down as a protocol that merely has not been given the means to.

The obligations no element discharges are these.
\begin{itemize}\setlength{\itemsep}{0pt}
  \item the reward rate on every pooled deposit is multiplied by a boost derived from the protocol's share of a different token's escrowed supply, held by the protocol rather than by the depositor, up to a hard multiple of two and a half
  \item the governance token of the upstream protocol, when deposited, is added to a single permanent escrow position from which no party can ever withdraw it
  \item that escrow is re-extended toward its maximum term by any permissionless deposit call whenever it would otherwise begin to decay
  \item in exchange the depositor receives a transferable token whose only exit is sale on a secondary market, since the position behind it has no redemption function
  \item voting weight over the escrowed position is non-transferable and is a function of the time remaining on a lock
  \item lock-holders vote in a proposal on this protocol whose weights are anchored to a historical epoch balance, and a separate executor contract later casts the aggregated result into the foreign protocol's gauge controller through a registered surrogate
  \item the protocol acquires and exercises governance rights inside four foreign protocols that it does not control
  \item half a per cent of that revenue is paid to whichever anonymous address calls the harvest function
  \item arbitrary third-party reward tokens attached to an upstream gauge are forwarded through per-pool stash contracts and virtual-balance children whose balances mirror the parent pool without holding its assets
  \item no deployed contract's code can be replaced: the protocol commits to immutability and offers shutdown-and-redeploy as its only migration path
\end{itemize}
\end{remark}

\begin{remark}
The witness adds streaming accrual, a direct redemption right and fee distribution, and drops the proxy and the timelock. The first three are read from the code: each harvest is paid out linearly over a rolling period rather than at once, the pooled receipt is withdrawable one-for-one at any time and without fee, and the levy on harvested revenue is split between wrapper-stakers, lock-holders and the treasury, with wrapper-stakers electing which asset they take their share in. The last two are read from their absence --- no deployed core contract sits behind a proxy, none carries a timelock, and the only migration path the protocol offers is to shut itself down and be redeployed --- and among the accounting alternatives the shares and the redemption right are mutually redundant, both kept because the receipt is a share of a pool and is simultaneously withdrawable at par, and dropping either loses one of those facts. Emissions stay, and must be stated exactly: the schedule is intact in code and has stepped down to the last cliff before it stops, so a construction that reproduces this protocol's economics from issuance reproduces a protocol that no longer exists.
\end{remark}

\subsection{Huma Finance V2}\label{sub:cat:yld:huma}

A depositor brings a stablecoin and receives a transferable token in one of two classes: a class paying an administered rate that the issuer states is subject to revision at its discretion, or a class paying no cash yield whatever and maximising an off-chain point balance instead. The holder switches between them instantly, for the cost of gas. The yield originates in transaction fees paid by businesses using an accelerated cross-border payment rail --- in practice short-duration advances to a single named counterparty --- and alongside this runs a permissioned product with the tranched architecture, where approved lenders take a senior or a junior claim, borrowers draw against a revolving line, a receivable-backed line or an outright purchase of the receivable, and one address alone sets the credit limit, sets the rate and declares default. Exit from the retail claim is first-come-first-served against a global daily cap that resets at midnight, with requests above the cap rejected outright rather than queued; exit from the tranched product is an epoch batch, senior filled first and shortfalls pro-rated into the next round. A holder who will not wait sells the claim into a third-party market maker.

\begin{measurement}\label{meas:cat:yld:huma}
Take $X = \{Aw,Bs,Ep,Fd,Ft,Gp,Sh,Sv,Tr,Uc,Wq\}$, with canonical form
$\mathrm{ex}(X) = \{Bs,Ep,Fd,Ft,Gp,Sh,Sv,Tr,Uc,Wq\}$; the elements \{Aw\} are derived rather than chosen.
It is not admissible: it leaves the requirement term $L3\!:\!At$ open; it arms $X11a*$. It discharges 12 of the 24 recorded obligations; 12 are residue. No composite of at most three corpus protocols equals it; no corpus protocol is contained in it.
\end{measurement}

\begin{remark}
The rejection is not a defect in the witness. The requirement tables ask for an attestation behind undercollateralised credit and this protocol has none, so the term stays open and the verdict states precisely what is true of it --- and the residue then says why that matters, because the financed obligation is an assertion by its own minter and the solvency of the pool rests on agreements in a legal system the code cannot reference. On the other side of the same protocol the category's mandate appears again, one step further than elsewhere: the address that determines an outcome on a claim also sets the \emph{price} of the credit and sizes the exposure, and of those three powers only the first is named. It is also the only agent in this category required to hold capital in the very pool whose credit it underwrites, and the vocabulary cannot record that either, so the one instance where the discretion is bonded reads the same as the instances where it is not. What is left is a retail claim whose holder elects between cash and points, and a redemption rule that is a rationing cap plus a promise made outside every contract.

The obligations no element discharges are these.
\begin{itemize}\setlength{\itemsep}{0pt}
  \item the default claim class pays an administered eight per cent annual rate, stated as subject to monthly adjustment at the issuer's discretion
  \item the holder may switch instantly, for gas only, to a second claim class that pays no cash yield at all
  \item in exchange the zero-yield class maximises an off-chain point balance which is the holder's entire return
  \item a depositor may optionally commit for three or six months to multiply that point rate; commitments can be lengthened but never shortened, and breaking one retroactively forfeits the bonus, including on transfer
  \item redemption is first-come-first-served against a global daily cap that resets at midnight UTC, and requests above the cap are rejected outright rather than queued, against a documented service level of one to seven days
  \item a holder may bypass the redemption path entirely by selling the claim token into third-party automated market makers
  \item the yield paid to depositors originates in transaction fees from businesses using an accelerated cross-border payment rail, in practice short-duration advances originated by a single named counterparty
  \item the redemption round and the price refresh happen only because scheduled off-chain jobs run and call them
  \item the senior claim may not exceed four times the junior claim, enforced on deposit and again in the ordering of redemptions
  \item the same address alone sets the borrower's credit limit and the yield in basis points, and may revise both afterwards by raising the limit, changing the yield, extending the term or waiving a late fee
  \item that address is required to hold junior-tranche capital in the very pool whose credit it underwrites, checked as a liquidity requirement
  \item the financed obligation is minted as a non-fungible token recording an amount, a currency code, a creation date, a maturity date, a paid amount and a creator, by whoever holds the minting role
\end{itemize}
\end{remark}

\begin{remark}
The witness adds a servicing discretion, a guardian and fee distribution, and drops the attestation and the accrual index. Servicing is the corpus's own candidate and it fits exactly one act here, that a single address declares a borrower in default and the loss it declares flows straight into the seniority waterfall; the index goes because neither claim moves by one, the retail return being a stated rate and a point balance. The attestation goes because there is none to name: the financed obligation is minted by whoever holds the minting role, recording an amount, a currency code, a creation date and a maturity date, with no payer, no document hash, no jurisdiction and no attester --- so among the alternatives the undercollateralised-credit term leaves open there is none the code supports, the term stays open, the construction is rejected and a prohibition arms. The larger departure is which product is being decomposed: the corpus element set describes the permissioned tranched pools, whose code is open and frozen, while every unit of the capital sits in a permissionless retail product on another chain whose program is closed-source and returns unverified from the verified-builds registry, so the witness carries both and says which obligations belong to which.
\end{remark}

\subsection{Pendle}\label{sub:cat:yld:pendle}

A depositor brings a yield-bearing token that already exists somewhere else and wraps it in a uniform interface over that source. The wrapper splits, in accounting-asset units, into a principal claim dated to a fixed maturity and a yield claim on everything the source pays until then. The yield is the source's own, and it reaches a holder by whichever leg was bought: the principal claim taken below par and redeemed at par, the yield claim taking the variable interest and the reward tokens, the pool claim taking swap fees, the emissions the gauge streams to that pool. Nobody chooses where a depositor's capital goes --- the wrapper is a passive adapter over one named upstream and the depositor picks the maturity and the market --- and the only discretion the design ever held directed emissions, exercised by an escrow that now sits in the protocol's own deprecated manifest beside its voting controller, replaced by a plain stake with no lock weighting. Exit before maturity is a sale into the pool or the order book, or burning both legs back into the wrapper; at maturity the principal claim redeems alone.

\begin{measurement}\label{meas:cat:yld:pendle}
Take $X = \{Em,Ep,Fd,Ft,Gp,Ix,Py,Rd,Sh,Up,Wq\}$, with canonical form
$\mathrm{ex}(X) = \{Em,Fd,Ft,Gp,Ix,Py,Sh,Up,Wq\}$; the elements \{Ep,Rd\} are derived rather than chosen.
It is admissible: no requirement term is open, every element is warranted, and it arms no prohibition. It discharges 11 of the 21 recorded obligations; 10 are residue. No composite of at most three corpus protocols equals it; 2 corpus protocols are contained in it, and together they supply every element except \{Fd,Wq\}.
It is not minimal: \{Gp,Up\} can be removed with every obligation still covered and admissibility intact.
\end{measurement}

\begin{remark}
Pendle is the one member of this category whose residue is \emph{mechanism} rather than discretion: every obligation left uncovered is an on-chain formula, and not one of them is a choice anybody makes. The pool-pricing symbols all name static invariants, so the same reserves imply the same price forever, and this market's do not --- its stiffness and its intercept are functions of the time remaining, and the intercept is rewritten before each trade so that the rate the previous trade left is the rate the next trade meets. The accounting symbols all describe a claim whose notional is preserved or grows, and one of the two claims here shrinks to nothing by design while the index that measures it ratchets upward and pushes every drawdown onto the other. So the gap is not that this venue was read carelessly; it is that the vocabulary cannot make a curve a function of a clock, nor a claim a function of what remains of its own life, nor a fee a spread in rate space rather than in price space. That the residue takes this shape is itself the result: at Pendle no agent chooses anything, and its uncovered obligations look nothing like those of the applications where one does.

The obligations no element discharges are these.
\begin{itemize}\setlength{\itemsep}{0pt}
  \item the accrual index is a monotonic high-water mark: if the source exchange rate falls the yield claim stops accruing entirely, and the drawdown is borne by the principal claim at redemption
  \item the yield claim's notional shrinks to zero by construction as the remaining accrual window closes
  \item after maturity the yield claim's index freezes at a snapshot, only already-accrued interest remains claimable, and every reward token subsequently arriving on that expired claim is swept to the treasury
  \item the principal claim is priced against the wrapper by a pool whose invariant is a logit curve in implied-rate space, whose stiffness parameter grows without bound and whose intercept decays as time-to-expiry shrinks, so identical pool state prices differently on a different day
  \item the curve's intercept is recomputed before every single trade so that the pre-trade implied rate equals the implied rate left by the previous trade
  \item the swap fee is quoted as an interest-rate spread whose multiplier is exponential in years-to-expiry, so the same trade pays less notional the nearer maturity is
  \item anyone may create a new maturity, constrained only to a fixed 24-hour calendar grid, and anyone may create additional markets on an existing principal claim that differ only in curve parameters
  \item a staker may instead exit immediately by paying a percentage of the staked amount, at a rate the owner may set up to one hundred per cent
  \item the yield claim accrues off-chain-attested, non-token reward streams from the source protocol, and the same five per cent protocol fee is levied on them
  \item execution authority over privileged functions is granted per (caller, function selector, target) triple rather than per role or per contract
\end{itemize}
\end{remark}

\begin{remark}
The witness is exhibited rather than derived, and it adds terms the corpus record omits, each carried by a mechanism in the source: a pro-rata claim on the market pool, whose minimum liquidity is locked permanently at initialisation; a queued exit on the staked governance token, which offers a cooldown or a priced jump ahead of it; and fee distribution, since every swap fee is split between the treasury and the liquidity providers and a further levy is taken from every yield claim including its non-token rewards. Nothing is dropped, and the contested-register vote-escrow marker is not carried at all: the escrow it names is deprecated in the protocol's own mainnet manifest alongside its voting controller and both fee distributors, and its replacement contains no lock-time weighting and no gauge vote, so a construction that emitted it would be decomposing a retired subsystem. Among the claim-accounting alternatives the separation term leaves open, the index is the one the code supports --- the wrapper's redemption value tracks a stored source rate rather than a rebasing balance --- and the pro-rata share is carried separately, for the pool claim rather than for the split. This is the only application in the corpus that carries the separation at all.
\end{remark}

\subsection{Spark Savings}\label{sub:cat:yld:spark}

A depositor brings the system's stablecoin and receives a share token whose nominal balance never moves and whose redemption value rises. The rate it earns is a written number held on the token itself, compounding per second into a stored accumulator; it is not derived from a utilisation curve, an auction or any market, and the only bound inside the token is that it be non-negative. The payment is funded by minting internal debt against the system's surplus account, unconditionally and with no read of revenue, reserves or any solvency condition, so the saver is paid first and the shortfall is left behind as protocol debt. Nobody allocates the depositor's capital, because the claim holds no assets --- it is a claim on a balance sheet --- and the deployment decision sits one level up, where a named multisig moves the system's own capital across an enumerated set of venues under caps that refill continuously with elapsed time, watched by a separate multisig that can freeze it without being able to use it. Exit is redemption at the accumulator at any time, or a swap against the reserve asset through a peg module at no fee.

\begin{measurement}\label{meas:cat:yld:spark}
Take $X = \{Aw,Ex,Gp,Ix,Ps,Rd,Sh,Sr,Tg,Up,Xf,Xm\}$, with canonical form
$\mathrm{ex}(X) = \{Aw,Ex,Gp,Ix,Ps,Rd,Sh,Sr,Tg,Up,Xf,Xm\}$, every element being primitive in it.
It is admissible: no requirement term is open, every element is warranted, and it arms no prohibition. It discharges 12 of the 21 recorded obligations; 9 are residue. No composite of at most three corpus protocols equals it; 4 corpus protocols are contained in it, and together they supply every element except \{Ex\}.
It is not minimal: \{Sh,Ix,Sr,Ex\} can be removed with every obligation still covered and admissibility intact.
\end{measurement}

\begin{remark}
What the construction cannot say here is most of what makes this a savings product rather than an accounting identity. The rate is a number a named agent chooses at its own discretion; the code holds a floor, a ceiling, a maximum step per update and a cooldown between updates, and it holds nothing that computes the number, triggers a change, or defines what a wrong rate would be. The same shape appears again one level up over venues instead of rates, with an enumerated action set standing where the floor and ceiling stood and a capacity regenerating with elapsed time standing where the step stood. Both are this category's mandate and neither has a symbol: the vocabulary's authority terms name a delay, a proxy and a pause, which is to say the facilities, never the \emph{envelope}. The remainder is the funding path, and no incentive, loss or stability term can express it, because the payment does not respond to a deficit --- it creates one.

The obligations no element discharges are these.
\begin{itemize}\setlength{\itemsep}{0pt}
  \item the interest owed is funded by minting internal protocol debt against the system's surplus/deficit account, unconditionally, with no read of revenue, reserves or any solvency condition
  \item the rate paid is a written number, not derived from a utilisation curve, an auction or any market; the only bound inside the savings token is that it be non-negative
  \item the enumerated agent chooses the value of the live savings rate at its own discretion, with no objective function, trigger, target or derivation anywhere in the code
  \item the agent may write only parameters drawn from an enumerated set of rate identifiers, each with its own configuration record
  \item the value written must lie between a governance-set floor and ceiling, live two per cent and thirty per cent
  \item a single update may move the rate by at most a governance-set step, live four percentage points, and no update may occur within a cooldown of the previous one, live sixteen hours; the update then takes effect immediately
  \item the remote token contains no accrual state and no drip function at all; its value is extrapolated locally from the forwarded tuple and its redeemability is limited by local peg-module inventory rather than by the mainnet balance sheet
  \item that role may act only through an enumerated set of twenty-nine venue-specific actions — mint, swap, deposit into a named vault, subscribe to a named fund, bridge by a named route — and cannot construct a new one
  \item each enumerated action carries a maximum amount and a slope that regenerates spent capacity continuously with elapsed time
\end{itemize}
\end{remark}

\begin{remark}
The witness adds a permission gate, an external feed and cross-domain message verification, and drops nothing. The gate is carried twice over, for two different agents: the rate may be written only by addresses enumerated in an on-chain set, and the savings token records the rate-setting module among its own authorised writers on mainnet, while the deployment of the system's capital is confined to one named role that holds custody of everything deployed. The verification and the feed enter together, because the remote representation is not the claim --- only the mainnet contract holds the accumulator and can mint against the balance sheet, while each remote chain receives a forwarded tuple over its canonical bridge, accepts it only after an application-level bound on the rate, and settles against local peg-module inventory. Among the control-plane alternatives the timelock is kept, and it is worth being exact about what it gates: governance's own parameter changes wait for it, and the agent's move inside the envelope does not.
\end{remark}

\subsection{BTCB}\label{sub:cat:bri:btcb}

A holder deposits bitcoin to an exchange and then withdraws it over BNB Chain, and what arrives is a BEP-20 the exchange issues. There is no source-domain event at all: the debit is a row in an internal ledger and the credit is a transfer of tokens the owner key has already minted. Redemption is depositing the token back and withdrawing bitcoin, which is an exchange withdrawal subject to account standing, network status and unilateral suspension. What a holder trusts is an externally-owned account that can mint without limit, a self-published reserve address whose issuer discloses that the amounts may not match and states no bound on the discrepancy, and an exchange that can refuse redemption with no on-chain act.

\begin{measurement}\label{meas:cat:bri:btcb}
Take $X = \{Rd,Xf\}$, with canonical form
$\mathrm{ex}(X) = \{Rd,Xf\}$, every element being primitive in it.
It is not admissible: it is ungrounded. It discharges 2 of the 15 recorded obligations; 13 are residue. No composite of at most three corpus protocols equals it; no corpus protocol is contained in it.
\end{measurement}

\begin{remark}
The construction is small because the system is weak, and that is the category's finding in a line. Issuance is unbounded and unilateral from one key with no quorum, no delay and no published process, and the control group has no symbol whose weakest reading is that weak --- so on the corpus arrays this entry is a subset of both other custodial wrappers and reads as the mildest of them, when on the evidence it has by a wide margin the weakest control plane in the category. The one dimension a holder can check runs the other way: with no code to freeze or seize a balance, the worst-governed wrapper gives the strongest on-chain guarantee, and the vocabulary states it only by omitting an element. Reserve claims stand twice over the same custodian, collateral for the pegged tokens and the exchange's own customer liabilities, published on separate pages with a disclaimer between them; the attestation element is the only symbol available for either, so using it for both would assert that a self-published address list and a proven liability snapshot are the same kind of object.

The obligations no element discharges are these.
\begin{itemize}\setlength{\itemsep}{0pt}
  \item What secures the peg is Binance's statement that it holds an equivalent amount of bitcoin in multi-signature cold storage; there is no message, no proof, no verifier and no threshold of any kind.
  \item Nothing is locked in a contract on either side: the issuer describes the mechanism as a swap in which 'their Bitcoin remains locked and can be redeemed by depositing BTCB back', with the lock being an internal custody arrangement rather than an escrow.
  \item Access to both mint and redemption runs through an exchange account, so only KYC'd customers in eligible jurisdictions can enter or exit, and the exchange can freeze an account and refuse redemption with no on-chain action whatsoever.
  \item Issuance is unbounded and unilateral: mint(uint256) is onlyOwner, owner() = getOwner() = 0xf68a4b64162906eff0ff6ae34e2bb1cd42fef62d, and eth\_getCode on that address returns empty, so a single externally-owned account can mint arbitrarily against a multi-billion-dollar liability.
  \item The on-chain artefact of a redemption is a self-burn -- burn(uint256) burns the caller's own balance -- with nothing in the transaction linking it to a bitcoin payout or to a redeeming party.
  \item The entire on-chain governance surface is transferOwnership and renounceOwnership: the contract is Binance's standard BEP20Token template (solc 0.5.16, Apache-2.0, verified Exact Match, not a proxy) with no pause, no blacklist, no forced transfer, no upgrade path, no timelock, no multisig and no registry.
  \item The authority over that surface is one key: not a quorum, not a contract, not a delay, with no published internal process binding its use.
  \item The issuer publishes a Proof of Collateral page for its pegged tokens reporting BTC at 100\%, with Proof of Assets of 68,200 BTC at a single published bitcoin address, against wrapped supplies of 2,900 on Ethereum, 128 on BEP2 and 65,222 on BEP20, plus Pioneer Burn 24 BTC and Token Recovered 49 BTC, covering 97 tokens.
  \item The issuer discloses that the peg may not hold exactly: 'the amount of locked BTC may not be precisely the same as the wrapped BTC due to the delays caused by the extensive auditing process', with no bound stated on the discrepancy.
  \item One reserve pool is drawn against by three pegged representations of the same asset on three different ledgers -- Ethereum, BEP2 and BEP20 -- so no single representation's backing is separable.
  \item Two reserve claims exist over the same custodian and must not be added together: 68,200 BTC of pegged-token collateral on the Proof of Collateral page, and 640,780.255 BTC of exchange customer reserve on the Proof of Reserves page (monthly Merkle-tree and zk-SNARK snapshot, latest 01/07/26 00:00 UTC at BTC block 956,140, ratio 100.08\%), the first page carrying an explicit disclaimer that it is not the exchange's proof of reserve.
  \item No independent auditor attests to the BTCB collateral specifically; no such report was found.
  \item A holder's transfers cannot be blocked and a holder's balance cannot be seized on-chain: the token has no freeze, no blacklist and no forced-transfer function.
\end{itemize}
\end{remark}

\begin{remark}
The corpus gives this token an emergency guardian and the witness drops it: the verified source and its function list carry mint, burn, transfer, approve and the ownership calls, and no pause. The access gate is dropped too, since there is no on-chain permission or identity check of any kind and reading the owner check on minting as one would make the element true of every ownable contract ever deployed. The freeze marker is dropped on the same evidence, the blacklist and forced transfer the corpus records being absent from the code. Both dropped powers are real and both are exercised inside a company where no contract can observe them; what survives is the wrap, the reserve listing and the redemption claim --- the smallest construction in the category, and the one the checker admits without complaint.
\end{remark}

\subsection{Coinbase Bridge}\label{sub:cat:bri:coinbase}

A customer withdrawing bitcoin from an exchange account to a supported network receives a wrapped ERC-20 instead of the asset, and the exchange keeps the asset. The crossing is a side effect of a transfer rather than a transaction the user submits, and sending the token back to an account unwraps it. There is no source-domain event to verify and, for all but one of the source assets, no source chain that could run a verification contract. What a holder trusts is that the issuer holds the underlying and keeps unwrapping, that a redemption right conditioned on account standing and on jurisdiction is not withdrawn, and that the single keys holding the pause, the blacklist, the minter administration and the proxy admin are not turned against them.

\begin{measurement}\label{meas:cat:bri:coinbase}
Take $X = \{Aw,Fz,Gp,Rd,Up,Xf\}$, with canonical form
$\mathrm{ex}(X) = \{Aw,Fz,Gp,Rd,Up,Xf\}$, every element being primitive in it.
It is admissible: no requirement term is open, every element is warranted, and it arms no prohibition. It discharges 6 of the 18 recorded obligations; 12 are residue. No composite of at most three corpus protocols equals it; no corpus protocol is contained in it.
\end{measurement}

\begin{remark}
The powers here are the strongest in the category and the custody of them the weakest, and only the first half is in the construction. The keys that can replace the implementation, blacklist a holder and appoint a minter are plain externally-owned accounts; the prohibition written for exactly this shape --- an upgrade path under immediate single-key control --- therefore cannot be evaluated, because its antecedent is unstateable. The only quantitative bound on any authority among the custodial wrappers, a rate limit on minting with programmatic replenishment, has no element either. Redemption is a jurisdiction vector rather than a right, the same token being redeemable in one country and not in another, enforced at an account layer no contract in the construction can see. And the measured size of the system is the issuer's own published file, which the attestation element names as an attestation and cannot name as a circularity.

The obligations no element discharges are these.
\begin{itemize}\setlength{\itemsep}{0pt}
  \item What secures the peg is that Coinbase holds the underlying asset and continues to unwrap on demand; there is no message, no proof, no verifier set and no threshold anywhere in the system.
  \item No contract escrows the underlying on any source chain, and five of the six source assets live on chains that cannot run a verification contract at all -- Bitcoin, XRP, Dogecoin, Cardano and Litecoin -- while the sixth, MegaETH, can.
  \item The conversion is a side effect of a transfer rather than a transaction the user submits: the user asked to withdraw BTC, and the system delivered cbBTC; sending the wrapped token back to a Coinbase account auto-unwraps it.
  \item The redemption right is conditioned on jurisdiction and asset: cbBTC cannot be sent or received by accounts in Canada, Japan and Georgia, and the other five wrapped assets are blocked in over a hundred jurisdictions including most of the EEA, Japan, Singapore, Hong Kong and New York State.
  \item The issuer publishes a live per-address proof of reserves -- 96,065.18 BTC of reserve against 96,054.41 cbBTC of supply, 20 reserve addresses, refreshed 8/4/2026 8:45 PM, split 48,338.78 / 44,528.809 / 3,108.551 / 78.272 across four networks -- and the Ethereum-side supply reads 48,338.77958531 cbBTC on chain, matching the first bucket exactly.
  \item Third parties treat the issuer's own reserve JSON as ground truth for the protocol's size: DefiLlama's adapter calls getConfig('coinbase-cbbtc-proof-of-reserves', 'https://www.coinbase.com/cbbtc/proof-of-reserves.json') and sums the declared addresses.
  \item Minting is rate-limited with programmatic allowance replenishment: a MintForwarder contract lets a minter continuously mint up to N tokens over M time without a cold masterMinter key having to top the allowance back up.
  \item The holders of these powers are single keys: eth\_getCode returns empty for both owner and masterMinter, so there is no multisig contract and no timelock anywhere in the admin surface.
  \item Authority is partitioned across keys: the owner can change owner, pauser, blacklister, masterMinter and oracle, and cannot change the proxy owner, so upgrade authority sits behind a separate key from role administration.
  \item The contract that holds every balance is not the issuer's own code: the repository states the token 'will be an exact duplicate of centre-tokens [FiatTokenV2\_1]' and the proxy an exact duplicate of the centre-tokens proxy, and the issuer's repository has not been touched since 2023-05-12 and never mentions cbBTC or any of the newer wrapped assets by name.
  \item The custodian of the underlying is also the largest venue on which the underlying trades, so the party whose failure would break the peg is the same party whose order book would price the break.
  \item One reserve pool backs four networks' supply of cbBTC simultaneously (48,338.78 Ethereum, 44,528.809, 3,108.551 and 78.272 on the others), reconciled only on the issuer's page.
\end{itemize}
\end{remark}

\begin{remark}
The upgrade element is right and load-bearing --- the token is a genuine proxy whose implementation slot resolves to a verified fiat-token implementation, and nothing in the other custodial wrappers is comparable --- so the witness keeps it. It also promotes the freeze from a marker to a firm assignment on a live named blacklister role whose accessors answer, which no other custodial wrapper in this category can offer. Both the freeze and the permission gate are recorded against that obligation, because both genuinely discharge it and the vocabulary gives no rule for choosing; the checker duly reports the freeze as removable, so the one power that distinguishes this issuer from the other two reads as redundant. Nothing in the witness records that the holders of these powers are individual keys, because the vocabulary has nothing to record it with.
\end{remark}

\subsection{Hyperliquid Bridge}\label{sub:cat:bri:hyperliquid}

A user transfers native USDC into an escrow contract on Arbitrum and is credited on the destination chain's own ledger; nothing is minted, and what the user holds is a balance rather than a token. Both directions are authorised by signatures from validators holding a supermajority of staking power --- the set that already secures the destination chain --- so the bridge inherits the chain's trust model rather than adding one to it. A withdrawal is deducted first and released only after a dispute period, during which named lockers may vote the bridge shut. The escrow has no owner and no upgrade path: every privileged action, including re-seating the validator set, runs on validator signatures, and the cold keys that do the re-seating belong to the same principals as the hot keys that sign ordinary flow.

\begin{measurement}\label{meas:cat:bri:hyperliquid}
Take $X = \{Gp,Gs,Vl,Wq,Xf,Xm\}$, with canonical form
$\mathrm{ex}(X) = \{Gp,Gs,Vl,Wq,Xf,Xm\}$, every element being primitive in it.
It is not admissible: it leaves the requirement term $L21\!:\!Au$ open; \{Gs\} has no consumer present. It discharges 7 of the 19 recorded obligations; 12 are residue. No composite of at most three corpus protocols equals it; no corpus protocol is contained in it.
\end{measurement}

\begin{remark}
This is the only entry in the category whose safety is not a legal promise, and the construction cannot say so. That the verifying set is endogenous --- that it is the destination chain's own validator set --- is the most important property a native bridge has, and the verification element asserts the same thing here as it would of a lone external attester. The reserve needs no attestation because it is a contract balance anyone may read, and the vocabulary expresses that only by omitting the attestation element, which is the string an unaudited custodian also produces: two opposite epistemic states, one representation. The threshold is then evaluated over a distribution no element carries, one operator family holding enough stake to veto every withdrawal and every validator-set update on its own, so a supermajority rule over a concentrated set and a supermajority rule over an even one decompose identically.

The obligations no element discharges are these.
\begin{itemize}\setlength{\itemsep}{0pt}
  \item The set that signs withdrawals is the same stake-weighted validator set that secures the destination chain, so holding bridged funds adds no trust assumption beyond using Hyperliquid at all.
  \item A deposit below the 5 USDC minimum is not credited and is lost permanently.
  \item A posted withdrawal is released only after a 200-second dispute period, during which lockers may vote to lock the bridge and stop it; disputePeriodSeconds() = 200 and blockDurationMillis() = 350 are on-chain readable.
  \item Nothing is minted on the destination: what the user receives is a balance in HyperCore's ledger, not a transferable token, so there is no destination-side representation to hold, trade or bridge onward.
  \item Two supplies of the same asset coexist on the destination chain with different trust models and are fungible to the user: \$421.9m that came through this escrow and \$5.48b issued natively on Hyperliquid via CCTP.
  \item The escrow has no owner and no upgrade path: implementation(), admin() and owner() all revert, and every privileged action is authorised by validator signatures rather than by an admin key.
  \item The same principals hold two key tiers with different powers: hot keys sign routine deposits and withdrawals, cold keys invalidate withdrawals and re-seat the validator set; hotValidatorSetHash() and coldValidatorSetHash() are distinct on-chain state.
  \item Two named roles exist that are neither owners nor governance: lockers, who can vote the bridge into a lock, and finalizers, who submit the transaction that releases a withdrawal after the dispute period.
  \item Validator stake is not slashable: the documentation states that there is currently no automatic slashing implemented, so the only consequences for a validator are jailing by peer vote and the seven-day queue, and no path exists from a bridge failure to a loss of stake.
  \item The two-thirds threshold is evaluated over a concentrated distribution: five Hyper Foundation validators hold 212,971,320 of 436,005,726 HYPE, or 48.85\%, so one operator family is far above the one third needed to veto any withdrawal or validator-set update and needs only a further sixth of stake to act alone.
  \item Validator-set updates are themselves subject to the dispute period and are ordered by an on-chain epoch counter, epoch() = 7.
  \item No reserve attestation exists or is needed: the escrow is a contract balance that anyone can read directly, and the issuer publishes nothing about it.
\end{itemize}
\end{remark}

\begin{remark}
The corpus carries a staked backstop and the witness drops it: the documentation states that no automatic slashing is implemented, so no path runs from a bridge failure to a loss of stake and there is nothing for the element to name. The queue element is not dropped but moved. The corpus attaches it to the dispute period, where it is wrong --- that delay exists so an honest validator can veto a fraudulent withdrawal, not to ration scarce liquidity among exiting claimants --- and the genuine queue in this system is the unstaking delay with its cap on concurrent exits, which governs the capital that constitutes the verifier set. The gas element is carried from the candidate register for the finalizer that pays the destination transaction, and it is precisely what makes this construction inadmissible: the law behind it demands a user-granted delegation, and the user grants nothing.
\end{remark}

\subsection{LayerZero V2}\label{sub:cat:bri:layerzero}

Not a bridge but the messaging layer bridges are built on, so its guarantees are conditional on the application above it. A packet leaves an immutable endpoint under a per-pathway nonce, is verified on the destination by the verifier networks \emph{the application itself chose} under a threshold it set, and is delivered by an executor that fronts the destination gas and is repaid on the source chain. The token half belongs to the application: under one standard the supply is burned and reminted and nothing is escrowed, under the other the original is locked in a lockbox and a representation circulates. What a holder trusts is therefore not this protocol but a configuration held by an application owner, changeable without delay or notice, which may name a single company as the sole verifier.

\begin{measurement}\label{meas:cat:bri:layerzero}
Take $X = \{Au,Gs,Xf,Xm\}$, with canonical form
$\mathrm{ex}(X) = \{Gs,Xf,Xm\}$; the elements \{Au\} are derived rather than chosen.
It is admissible: no requirement term is open, every element is warranted, and it arms no prohibition. It discharges 6 of the 18 recorded obligations; 12 are residue. No composite of at most three corpus protocols equals it; no corpus protocol is contained in it.
\end{measurement}

\begin{remark}
The residue is the product. Verification here is a per-application parameter, so the element carried for it asserts the same of a lone company and of a broad independent set; when a pathway's configured set was a single verifier network and the nodes behind it were compromised, a destination minted against no source burn, and the decomposition of that pathway is identical before and after. Inherited security --- what an application gets when it configures nothing --- has no symbol in either of its failure directions, the bricked pathway or the unguarded one. Nor does the separation the design rests on: the verifiers hold safety, the executor holds liveness, and the loss of each has a different consequence. The two supply topologies this one protocol offers have opposite attack surfaces, the same user-visible interface, and one element between them.

The obligations no element discharges are these.
\begin{itemize}\setlength{\itemsep}{0pt}
  \item Verification security is a per-application parameter rather than a protocol property: two applications on the same endpoint pair can run entirely different verifier sets and thresholds, and one of them may be a single company.
  \item An application that does not set its configuration inherits protocol defaults, which the documentation warns are placeholders that 'may be Dead DVNs that prevent message delivery' -- a Dead DVN is deployed on each chain, 0x747c741496a507e4b404b50463e691a8d692f6ac on Ethereum -- while the historical default in the other direction was a 1-of-1 LayerZero Labs DVN.
  \item Safety on a pathway is exactly the honesty of the configured verifier set: on 2026-04-18, with rsETH configured to a single verifier -- the LayerZero Labs DVN -- and two LayerZero-hosted RPC nodes compromised, approximately \$292m / 116,500 rsETH was minted on a destination without a matching source burn.
  \item Safety and liveness are held by two different role types: DVNs decide whether a message is true, Executors decide whether it arrives, and the loss of each has a different consequence -- a corrupted DVN set mints falsely, an absent Executor merely stalls.
  \item The two topologies have opposite attack surfaces and the same user-visible interface, so the question 'how much can be stolen' has a different answer per token on one protocol.
  \item Mint and burn authority, and custody of any lockbox, belong to the application's own contract and its owner, not to LayerZero, so the party that must be trusted for supply conservation is one the protocol does not name.
  \item The Endpoint's code cannot be replaced -- 24,005 bytes of direct non-proxy bytecode, documented as immutable and permissionless -- and it is nevertheless owned: owner() = 0xbe010a7e3686fdf65e93344ab664d065a0b02478, a custom multisig with threshold() = 3 and nonce() = 611, carrying no secondsTimeLocked and no getMinDelay.
  \item There is no protocol-level pause: paused() and pause() both revert on EndpointV2, so no party can stop message flow by pausing the endpoint.
  \item The Endpoint owner controls which message libraries are registered -- getRegisteredLibraries() returns 4 entries, including a blockedLibrary at 0x1ccbf0db9c192d969de57e25b3ff09a25bb1d862 -- and the protocol-level defaults, so delivery on a pathway can be stopped or redirected by configuration.
  \item An application's send and receive configuration -- its verifier set, its threshold and its libraries -- can be changed at any time by the application's owner, with no delay and no notice, including while messages are in flight.
  \item Fees may optionally be denominated in a protocol token, and this is a per-chain parameter: lzToken() is address(0) on Ethereum, so ZRO fee payment is not enabled there while it is elsewhere.
  \item A received message can itself trigger further cross-domain calls through lzCompose, so a transfer composes horizontally into a sequence of domain-crossing actions.
\end{itemize}
\end{remark}

\begin{remark}
The corpus decomposes this protocol with an upgrade proxy, a guardian pause and a timelock, and the witness drops all three on direct reads of the endpoint: the deployed code carries no implementation slot, the pause accessors revert, and the owning multisig executes on signature count with no delay accessor present. What is mutable is the registered library set and the protocol defaults, which is a different mechanism from mutable code and has no element at all. In their place the witness carries the delegation an application owner may grant over its own security stack, and the executor's cross-domain gas sponsorship from the candidate register --- the same shape the Hyperliquid finalizer takes, which is the case for reading that element across domains rather than only as an account-abstraction paymaster.
\end{remark}

\subsection{Binance Wallet}\label{sub:cat:int:binance}

A wallet front-end that routes its users' orders, with no public repository anywhere and no primary source for the quote mechanism the corpus assigns it. Its position in the category is owed to the order flow it originates rather than to any mechanism it operates.

\begin{measurement}\label{meas:cat:int:binance}
Take $X = \{Ag,Aw,Up\}$, with canonical form
$\mathrm{ex}(X) = \{Ag,Aw,Up\}$, every element being primitive in it.
It is admissible: no requirement term is open, every element is warranted, and it arms no prohibition. It discharges 3 of the 19 recorded obligations; 16 are residue. No composite of at most three corpus protocols equals it; 2 corpus protocols are contained in it, and together they supply every element except \{Aw,Up\}.
\end{measurement}

\begin{remark}
This venue's entire economic position is that it owns the flow. The vocabulary has four symbols for how an order is matched and none for who owns it, so the thing that makes this a top-five venue is the thing the model cannot see.

The obligations no element discharges are these.
\begin{itemize}\setlength{\itemsep}{0pt}
  \item the user is shown a single recommended provider rather than a route, and the rule by which that provider was chosen is not stated
  \item whether any counterparty ever signs a firm quote that this system then fills is not stated at any primary source; the published wording ('CeFi and DeFi exchanges') is equally consistent with routing into an exchange-operated on-chain pool
  \item the user sets a slippage tolerance beside a displayed Minimum Received, and bears every adverse move between signature and inclusion; all swaps are final once confirmed
  \item the user pays network gas on every swap, including on swaps that fail, and no fee sponsorship is offered at any point
  \item the user's transactions are shielded from the public mempool by default, by a private-RPC programme operated by the host chain and its partners rather than by this application
  \item the settlement contract's bytecode is verified against a source blob the operator uploaded to a block explorer, while no public, versioned or reviewable codebase for it exists anywhere
  \item the wallet holds the user's signing key as three MPC shares rather than as a single secret, so the party that routes the trade is also a party to producing the signature
  \item the swap sits inside an application that already holds the user's KYC-verified exchange account and balances, so the party that owns the demand also chooses the venue, and no consideration for that choice is disclosed in either direction
  \item the right to fill is not auctioned, not tendered and not priced: there is no solver set, no bidding, no eligibility rule and no published list of the providers that supply the other side
  \item the application claims to find the most competitive price from thousands of options on CeFi and DeFi exchanges, with no benchmark named, no measurement published and no disclosure of the alternatives that were rejected
  \item trading fees on all swaps were waived for eight months, scoped to particular products and to backed-up keyless addresses only, with imported wallets excluded
  \item limit orders placed in the Alpha product execute against on-chain markets and their fee was cut from 0.15 percent to 0.01 percent on one chain
  \item an exchange Spot or Funding balance can be turned into an on-chain trade with no wallet in the loop, and the order is auto-cancelled if the minimum received falls below the amount displayed
  \item a token bought this way cannot be withdrawn or transferred out of the operator's environment until the operator lists it on its own exchange
  \item the operator disclaims being the counterparty and disclaims responsibility for the third-party applications into which it routes the user's trade
  \item the same router is deployed independently on twelve chains; no value crosses a domain and no message is verified between domains
\end{itemize}
\end{remark}

\begin{remark}
The witness drops both the quote element and the sponsored-fee element the corpus carries, because no primary source supports either and the gasless claim is contradicted by the venue's own documentation. Dropping an unsupported element is the same discipline as adding a supported one; here it leaves the construction very thin, and that thinness is the evidence.
\end{remark}

\subsection{Jupiter}\label{sub:cat:int:jupiter}

A meta-aggregator whose own documentation lists other routers --- including two members of this category --- as competitors inside its routing. It is not a peer of the others so much as a layer above some of them.

\begin{measurement}\label{meas:cat:int:jupiter}
Take $X = \{Ag,Aw,Gs,In\}$, with canonical form
$\mathrm{ex}(X) = \{Ag,Aw,Gs,In\}$, every element being primitive in it.
It is not admissible: it leaves the requirement term $L21\!:\!Au$ open. It discharges 4 of the 20 recorded obligations; 16 are residue. No composite of at most three corpus protocols equals it; 2 corpus protocols are contained in it, and together they supply every element except \{Au,Aw,Gs,In\}.
\end{measurement}

\begin{remark}
That this venue routes to routers is the category's structural finding and has no expression here. Composition in this paper is union of element sets between peers; one protocol consuming another as a component is a containment relation the carrier cannot state. The yield category found the same gap one level down, where a curator allocates across protocols.

The obligations no element discharges are these.
\begin{itemize}\setlength{\itemsep}{0pt}
  \item quotes are solicited from independent third-party routing protocols as well as from venues, and the winning quote may come from a rival aggregator rather than from a pool
  \item candidate routes are simulated on-chain before one is chosen, and are ranked on predicted executed price rather than on quoted price
  \item quote sources that perform badly are automatically sidelined by the router without any governance act
  \item the market maker retains the right to decline execution after the taker has accepted its quote
  \item two mutually exclusive API paths exist with different guarantees: a managed path that assembles and lands the transaction, charges a platform fee and can use RFQ, and a raw-instruction path that is composable, carries no platform fee and can only use the in-house router
  \item the signed transaction is submitted through infrastructure the protocol owns, a validator and RPC of its own, with a minimum tip, rather than through a public endpoint or a third-party relay
  \item the protocol publishes a measured average positive slippage of +0.63 basis points as a product property, that is, a claim that executed prices beat quoted prices
  \item a platform fee is deducted inside the quote at a rate that depends on the asset classes being traded and rises to 50 basis points for any token listed within the previous 24 hours
  \item the fee is collected in a priority-ordered choice among the pair's tokens rather than in a fixed leg
  \item an integrator sets a referral fee between 50 and 255 basis points on its users' trades and the protocol keeps 20 percent of it; if the integrator's fee account is uninitialised the trade still executes and the integrator simply gets nothing
  \item resting limit, OCO and recurring orders are held in a custodial vault operated by a third-party key manager, stored off-chain and private by default, rather than in an account the user controls
  \item a keeper operated by the protocol fires those orders when a US dollar price condition is met rather than a pool rate, sets the slippage for each slice itself, and the output amount is explicitly not guaranteed
  \item the protocol assigns tokens a quality score and a verification status of its own devising, which condition how a token is surfaced and what it costs to trade
  \item the routing engine is distributed as a closed binary and no source for the deployed aggregator program is published; a published audit inventory stands in its place
  \item the router owns venues on the other side of its own routing: a perpetuals exchange, a lending market with flash loans, a liquid-staking token, a launchpad with bonding curves, a prediction market and a validator, so it can internalise the flow it receives
  \item the protocol occupies four positions in the order-flow market at once: it owns the default trade surface on its chain, it buys routing from rival protocols, it sells the right to fill to market makers including the last-look privilege, and it sells distribution to integrators at a stated split
\end{itemize}
\end{remark}

\begin{remark}
The witness drops the quote element, because the venue's fills are subject to a last look and a quote that can be withdrawn at fill time is not the firm quote the element names. It reassigns the intent element from a component that has become custodial to the one that still answers to the definition.
\end{remark}

\subsection{KyberSwap}\label{sub:cat:int:kyberswap}

A router that has become a venue owner: it now holds pools on two major exchanges where it is the only permitted taker, on a revenue split. The corpus records it as a bare router, which its current arrangement contradicts.

\begin{measurement}\label{meas:cat:int:kyberswap}
Take $X = \{Ag,Au,Ep,Fd,In,Rf,Xf\}$, with canonical form
$\mathrm{ex}(X) = \{Ag,Au,Ep,Fd,In,Rf,Xf\}$, every element being primitive in it.
It is admissible: no requirement term is open, every element is warranted, and it arms no prohibition. It discharges 6 of the 20 recorded obligations; 14 are residue. No composite of at most three corpus protocols equals it; 6 corpus protocols are contained in it, and together they supply every element except \{Au,Ep,Fd,Xf\}.
It is not minimal: \{Fd,Ep,Xf\} can be removed with every obligation still covered and admissibility intact.
\end{measurement}

\begin{remark}
Exclusivity over a venue is not the same obligation as routing to one, and the vocabulary has a symbol only for the second. A router that owns the pool it routes to, and is the only party allowed to trade against it, is a different object from a router; the model writes them the same way.

The obligations no element discharges are these.
\begin{itemize}\setlength{\itemsep}{0pt}
  \item one of the four routed liquidity classes is the operator's own limit-order product, so the router is also a venue operator on its own routes
  \item the user sets a slippage tolerance between 0 and 2000 basis points and a deadline, and settlement reverts unless the amount received clears the resulting minimum
  \item the venue is re-decided on-chain at execution time, per hop, by comparing candidate pools against their state in that block and switching away from the one named in the calldata
  \item that re-selection is aimed at specific adversaries: pools that have been sandwiched or front-run in the same block, market makers that have widened their spread, and liquidity withdrawn just in time
  \item the user pays no fee for the swap, and the operator's revenue is instead the surplus of tokens above the ESTIMATED output, together with unconsumed dust from partial fills, while the user is protected only to the confirmed MINIMUM
  \item the operator holds pools, deployed as hooks on third parties' automated market makers, whose swaps are restricted in the pool's own code so that only the operator's aggregator may take from them
  \item for each such trade the aggregator computes the best route with those pools and the best route without them, and treats the difference as the additional value the exclusivity created
  \item an off-chain signer produces a signed fair market price for that single trade, which is carried in the transaction and validated by the hook against the pool's actual output
  \item when the pool can do better than the signed price the taker is given the signed price and the pool keeps the excess; when it does worse the taker gets the actual output and nothing is kept
  \item takers of the operator's limit orders pay a fee keyed to a six-band volatility classification of the token, from 0.01 percent for super-stable pairs to 1 percent for super-high-volatility pairs
  \item several fee recipients may be paid on a single trade, each with its own amount, its own side of the trade and its own receiving address, supplied as parallel lists
  \item the terms a caller is offered depend on the wallet address in the request, which the documentation states is used to reach exclusive pools and rates
  \item the pool simulators that the route search runs on are published for all 129 liquidity sources, and third-party venue teams add their own venues to the router by pull request
  \item the operator's commercial asset is the venue rather than the user: it charges its users nothing, buys exclusivity over pools it did not build, pays the liquidity providers 70 percent of the arbitrage that exclusivity captures, and takes its own revenue from surplus at settlement
\end{itemize}
\end{remark}

\begin{remark}
The witness adds four elements the corpus omits --- a quote mechanism, an intent path, a fee distribution and an epoch gate --- each supported at the venue's own documentation. The corpus record predates all of them.
\end{remark}

\subsection{LiquidMesh}\label{sub:cat:int:liquidmesh}

A router with no public repository and no identified operating entity, whose implementation contract is unverified on the explorer for the chain it works on. It reaches the corpus on volume alone. The corpus records it as decomposing to a single element, and on the mechanism of routing that is correct.

\begin{measurement}\label{meas:cat:int:liquidmesh}
Take $X = \{Ag,Au,Aw,Up\}$, with canonical form
$\mathrm{ex}(X) = \{Ag,Au,Aw,Up\}$, every element being primitive in it.
It is admissible: no requirement term is open, every element is warranted, and it arms no prohibition. It discharges 6 of the 20 recorded obligations; 14 are residue. No composite of at most three corpus protocols equals it; 2 corpus protocols are contained in it, and together they supply every element except \{Au,Aw,Up\}.
\end{measurement}

\begin{remark}
The corpus records this venue and one other as having identical decompositions. They do --- and their residue is disjoint. This one sells latency and takes the other side of the flow it routes; the other owns pools where it is the only permitted taker. Two systems can agree on every symbol the vocabulary has and share no economic mechanism at all, which is a sharper statement of non-identifiability than a collision on its own.

The obligations no element discharges are these.
\begin{itemize}\setlength{\itemsep}{0pt}
  \item the route is computed off-chain on the operator's own servers and handed to the caller as pre-built calldata; no route search, quote or comparison of alternatives occurs on-chain and none is recorded
  \item the party that authored the route neither settles the trade nor bears price risk: the user signs and broadcasts an ordinary chain-native transaction and any adverse move between quote and inclusion is the user's
  \item settlement reverts unless the amount received meets a minimum return fixed in the signed calldata
  \item a quote is refused outright when the computed price impact exceeds the operator's threshold, returning error 40106 rather than a worse price
  \item the signed transaction may be handed back to the operator for broadcast through private channels rather than the public mempool, at one of two privacy grades on EVM (public, private) or three on Solana (off, reduced, secure)
  \item the builders and validators that receive the flow are forbidden by agreement from auctioning the user's transaction, and the flow is not disclosed to third-party searchers
  \item the operator extracts MEV by backrunning the very transactions it built, and returns the majority of the proceeds to the trader; on Solana it instead buys priority through staked-validator and Jito paths and returns a portion of the saved tips
  \item the guarantee sold is inclusion rather than price: quote latency, quote freshness for newly created tokens, landing rate and uptime are the published product properties
  \item a commission in basis points, taken from either the input or the output token at the caller's election, is charged on the trade and paid to the integrating wallet; the protocol defines the lever and collects nothing itself
  \item the only benchmark offered against the returned route is the operator's own midPrice and priceImpactPct in the same response; no external reference price, no price-improvement measurement and no record of the alternatives considered is published
  \item no party other than the user may fill a LiquidMesh order: there is no solver set, no auction, no RFQ tier, no maker and no filler of any kind
  \item the token being traded is itself screened for scam and honeypot characteristics before a quote is issued
  \item a bonding-curve launchpad is a first-class venue class, with dedicated trade types in the quote API and its own token-manager contract on BSC
  \item the same service is offered across EVM chains, Solana, Sui and Tron as independent per-chain routers; no value crosses a domain and no message is verified between domains
\end{itemize}
\end{remark}

\begin{remark}
The witness adds a delegated-execution scope and a permission gate on evidence of the settlement path, and keeps the routing element the corpus assigns. It adds nothing for the venue's economics, because nothing in the vocabulary answers to them.
\end{remark}

\subsection{OKX DEX}\label{sub:cat:int:okx}

An exchange-operated aggregator that now runs a sealed auction over its own affiliated solvers, and publishes a conflict-of-interest policy conceding that it may display a route which is not the best available.

\begin{measurement}\label{meas:cat:int:okx}
Take $X = \{Ag,Au,Aw,Gs,In,Rf,Xf\}$, with canonical form
$\mathrm{ex}(X) = \{Ag,Aw,Gs,In,Rf,Xf\}$; the elements \{Au\} are derived rather than chosen.
It is not admissible: it arms $X19*$. It discharges 7 of the 19 recorded obligations; 12 are residue. No composite of at most three corpus protocols equals it; 7 corpus protocols are contained in it, and together they supply every element except \{Au,Gs,Xf\}.
\end{measurement}

\begin{remark}
A venue that admits in writing that its displayed route may not be the best one is telling us the benchmark matters. There is no element for the benchmark, none for who may enter the auction, and none for the affiliation between the auctioneer and the bidders. The construction also arms a prohibition, which is the only place in this category where the algebra objects to anything at all.

The obligations no element discharges are these.
\begin{itemize}\setlength{\itemsep}{0pt}
  \item the signed intent is offered to competing solvers who bid for the right to fill it in an auction lasting about two seconds, and the best bid wins
  \item the bids in that auction are sealed, so no solver can see another's price before committing
  \item only solvers on an operator-maintained allowlist may fill an intent, subject to onboarding requirements, performance standards and anti-manipulation controls set by the operator
  \item solvers operated by or affiliated with the operator compete in the operator's own auction for the operator's own order flow, and the operator publishes a policy saying so
  \item the system, not the user, decides which of two incompatible execution regimes runs, on trade size, token classification and a user-chosen speed mode, and which party bears price risk between signature and fill changes with that decision
  \item the operator's own route may be presented as the default selection or visually emphasised in the interface, including where its quoted received amount is not the highest available
  \item a referrer commission is read out of the last words of the settlement calldata behind magic flags, capped by the router at 3 percent (10 percent on Solana), charged on one side of the trade only, deducted before execution, and paid to the integrator rather than to the operator
  \item the same routing engine is sold to other order-flow owners, powering swaps inside a rival wallet, and is supplied as one competing quote engine inside a rival aggregator
  \item the identity of the settlement contract is not a constant: routers may be replaced on upgrade and integrators are instructed to use whatever address the API returns rather than hardcoding one
  \item a self-hosted routing client is distributed as a closed binary behind an API key and a whitelist form, computes routes locally from streamed pool state, and will additionally search for the most profitable round-trip across a set of intermediate tokens
  \item the splitting algorithm, the private-market-maker set, the solver allowlist and the auction's conduct are all unverifiable from outside, and the canonical settlement repository the operator once published has been withdrawn
  \item the operator stands on both sides of the order-flow market at once: it owns users through an exchange and a wallet, owns a chain, sells its engine to rivals who own users, sells the right to fill through an auction it also bids in, and hands the fee lever to integrators
\end{itemize}
\end{remark}

\begin{remark}
The witness adds an intent element on evidence of the auction, and explicitly refuses the batch-auction element: many solvers bidding on one order is not a uniform price over a batch, and the two are different mechanisms that the vocabulary distinguishes. Refusing a nearby symbol is what keeps the remaining identity claims in this category meaningful.
\end{remark}

\subsection{BlackRock BUIDL}\label{sub:cat:rwa:buidl}

A qualified purchaser subscribes through a placement agent and receives a share in a British Virgin Islands company, recognised there as a professional fund and recorded on chain as a token. The company owes the money, and its sponsor states, in the same paragraph as the dollar target, that it has no legal obligation to support the fund. There is no price feed of any kind: the unit is assumed at par, and yield is declared off chain against a snapshot of who held the token at a wall-clock instant, then paid the following business day by minting new tokens. A holder redeems for cash on the same business day if the request reaches the transfer agent in good order inside the morning-to-afternoon window --- and the transfer agent, not the chain, decides whether it did.

\begin{measurement}\label{meas:cat:rwa:buidl}
Take $X = \{Aw,Fz,Gp,Rd,Sh,Up\}$, with canonical form
$\mathrm{ex}(X) = \{Aw,Fz,Gp,Rd,Sh,Up\}$, every element being primitive in it.
It is admissible: no requirement term is open, every element is warranted, and it arms no prohibition. It discharges 6 of the 19 recorded obligations; 13 are residue. No composite of at most three corpus protocols equals it; no corpus protocol is contained in it.
\end{measurement}

\begin{remark}
The strongest sentence in the category is the fund's own: a redemption effected on chain before the deadline may not be recorded on the fund's books in time, and in that event the shares are not redeemed on that date. The books are a distinct artefact from the chain and are authoritative over it, and the register keeper supplies the rest itself --- one control book and one master securityholder file, singular across chains, and on a fork the transfer agent \emph{determines} which chain is the authoritative record --- so the register is not a ledger but a determination by a named company, exercised over ledgers, and held under a rolling annual contract the fund may terminate on short notice. Nothing in the vocabulary reaches the register, so nothing reaches its tenure either. Every element this construction carries names the wrapper --- who may hold the token, who may halt it, who may replace its code --- and not one names the thing wrapped, which here is a claim on a company whose custodian, prime broker and auditor are absent from the one filing obliged to name them, and whose offering document sits with a regulator that publishes its existence and not its contents. The single-key upgrade path is likewise a fact about the \emph{holder} of a power rather than about the power: the control element says the implementation is mutable, and cannot say who may mutate it, how fast, or with what independence.

The obligations no element discharges are these.
\begin{itemize}\setlength{\itemsep}{0pt}
  \item the fund targets a constant one-dollar share value with no price feed of any kind on chain, warns in the same sentence that it may fail, and disclaims any sponsor obligation to support it
  \item yield is declared off chain against a snapshot of who held the token at three in the afternoon eastern time, and paid the following business day by minting new tokens to them
  \item the on-chain ledger is keyed on a legal identity rather than on an address: wallets are attached to an investor identifier, and the investor carries attributes that expire and that point at an off-chain document
  \item the gate tests two distinct legal predicates at once, accredited-investor status under the securities law and qualified-purchaser status under the investment-company law, and the tighter of the two is what caps the fund's investor base
  \item the holder must know whether the chain is the register of ownership or a mirror of a transfer agent's book, and here a redemption completed on chain before the deadline is not a redemption: if the transfer agent has not yet determined the transfer to be final and recorded it in the fund's books, the shares are not redeemed that day and roll to the next redemption date
  \item the registrar of that authoritative book is a distinct SEC-registered entity from the broker-dealer that sells the fund, and it outsources its own recordkeeping to a legacy stock transfer agent
  \item the party who holds that replacement power is one private key: the proxy owner is an externally owned account, with no multisig at the proxy layer and no delay of any kind
  \item the same fund issues a legally distinct share class on each chain, priced with a different management fee, against one off-chain register, and movement between chains is a third-party service the fund disclaims
  \item the holder must know who owes them the money and what recourse they have, and here the obligor is a BVI company whose sponsor has expressly disclaimed any obligation to support it and whose gating powers are unknowable
  \item the holder must know whether their claim survives the fund's insolvency, and here it is unsecured equity: no security interest, no collateral agent, no bankruptcy-remoteness claim and no perfection step anywhere
  \item the holder must know what the reserve holds and who custodies it, and here the only public statement of the portfolio is a marketing phrase and the custodian could not be established at a primary source
  \item no independent attestation of the fund's assets is published anywhere obtainable: there is no holdings feed, no reserve report, no oracle and no covenant
  \item the holder must know who may change what the money is invested in, and here the manager can change duration, issuer mix and repo exposure at will with nothing on chain constraining it
\end{itemize}
\end{remark}

\begin{remark}
The witness is the smallest in the category, and most of its departures from the corpus record are removals. The attestation element goes because no independent statement of the fund's assets is published anywhere obtainable, and because the sworn adviser filing that must name a custodian, a prime broker and an auditor names none of them. The peg-swap element goes because the par-conversion facility was always a separate commercial arrangement of another company and the fund's own page now disclaims third-party liquidity generally; the cross-domain elements go because each chain carries a legally distinct share class, priced with its own management fee, against one off-chain register, and movement between chains is a disclaimed third-party service rather than a mechanism of the fund. Against that, the witness arms a prohibition the corpus recorded as unarmed --- the token's proxy exposes an owner and a target and nothing else, and that owner is an externally owned account which may replace the implementation in a single call, with no delay and no second signature.
\end{remark}

\subsection{Centrifuge}\label{sub:cat:rwa:centrifuge}

An investor subscribes an asset into a vault attached to a share class of a pool and receives a share token that rises in price rather than in balance. The protocol owes nothing: the obligor is an off-chain fund --- for the flagship share class, a British Virgin Islands segregated portfolio company recognised there as a professional fund --- and that fund describes the token as \emph{prima facie} evidence of ownership, in the same way a share certificate would be. The net asset value is computed off chain and enters the protocol as a single number one party posts, with the freshness timestamp supplied by that same party and no dispute window, challenge bond or deviation bound anywhere in the call path; it reaches every other chain only when somebody pays gas to push it. Subscription and redemption are asynchronous --- a request is placed, the issuer approves a batch, shares are issued or revoked, the batch is fulfilled, and only then may the investor claim --- with the issuer setting each batch's execution price separately from the published one.

\begin{measurement}\label{meas:cat:rwa:centrifuge}
Take $X = \{Au,Aw,Ep,Ex,Fz,Gp,Rd,Sh,Sv,Tg,Up,Xf,Xm\}$, with canonical form
$\mathrm{ex}(X) = \{Au,Aw,Ep,Ex,Fz,Gp,Rd,Sh,Sv,Tg,Up,Xf,Xm\}$, every element being primitive in it.
It is admissible: no requirement term is open, every element is warranted, and it arms no prohibition. It discharges 11 of the 22 recorded obligations; 11 are residue. No composite of at most three corpus protocols equals it; 2 corpus protocols are contained in it, and together they supply every element except \{Au,Ep,Ex,Fz,Rd,Sh,Sv\}.
It is not minimal: \{Ex,Rd,Sv,Tg,Up\} can be removed with every obligation still covered and admissibility intact.
\end{measurement}

\begin{remark}
The register here is a different legal device again: the token is evidence of ownership under the fund's own law, which is to say rebuttable evidence of an entry in a book kept elsewhere --- and the book is kept by a fund administrator the fund discloses as a role and never names. That pattern runs through the whole disclosure: prime broker, custodian, administrator, each present as a function and absent as a party. On chain the residue is emptier than expected rather than richer, because a defaulted underlying asset has no representation of any kind --- no delinquency state, no roll rate, no write-off flag, no impairment call --- so the only way a loss reaches a holder is a lower number in the next price the manager posts. The construction the residue most needs to state is the wrapper on the wrapper: a restricted share wrapped into a freely transferable token whose economic exposure is bearer while the right to mint or redeem stays behind the identity gate, so a holder of the outer token has a real claim and no path to the obligor.

The obligations no element discharges are these.
\begin{itemize}\setlength{\itemsep}{0pt}
  \item that price is authoritative on one chain and a message-derived copy everywhere else, propagated only when somebody pays gas for it, and a per-chain manager may override the local copy with a different number
  \item eligibility carries an expiry: membership is stored as a valid-until timestamp packed into the same storage slot as the balance, so a holder ceases to be eligible with no actor doing anything
  \item a restricted share can be wrapped into a freely transferable token that trades on top of it, while the right to mint or redeem remains behind the identity gate
  \item the global pause is deliberately partial: it halts all cross-chain messaging in both directions and leaves local deposits, redemptions and withdrawals running
  \item there is no seniority anywhere in the protocol: multiple share classes exist, they can be named and permissioned differently, and nothing allocates a loss between them
  \item a defaulted underlying asset has no on-chain representation of any kind: no delinquency state, no roll rate, no write-off flag and no impairment call exist in the protocol
  \item the holder must know who owes them the money, and here the protocol owes them nothing: the obligor is an off-chain fund and Centrifuge is infrastructure
  \item the holder must know whether the chain is the register, and here the token is prima facie evidence of ownership under BVI law in the same way a share certificate would be
  \item the holder must know whether their claim survives the issuer's insolvency, and here a bankruptcy-remote structure is asserted by the issuer and no security trustee is named anywhere
  \item the holder must know what the reserve holds and who custodies it, and here the fund discloses the roles and names none of the parties
  \item the holder must know who may change what the money is invested in, and here nothing on chain records who selected an asset or on what criteria
\end{itemize}
\end{remark}

\begin{remark}
The corpus record carries a tranche element and the witness drops it, on a search of the protocol's own source: seniority, subordination and waterfall appear nowhere in it, and the only occurrences of the word tranche are in an audit document, a test-properties file and legacy migration scripts. The documentation's junior-and-senior phrasing describes a configuration idiom --- name share classes differently, permission them differently, price them differently --- and there is no loss-allocation logic, no seniority ordering and no trigger that redirects anything, so subordination here lives in offering documents and reaches the chain only as the manager choosing to mark one class down before another. In its place the witness carries a freeze-and-recovery element the corpus omitted, although this protocol ships more of that machinery than the fund tokens the corpus does credit with it; a redemption element, because redemption is a first-class operation with its own transfer-hook regime; and a delegated-execution element, for the merkle-rooted policies that restrain a strategist by target contract, function selector and parameter value. Epoch-gated exit replaces the queue: this is a request, approve and fulfil batch cycle priced by the issuer, and the queue element belongs to the private-credit member of this category, where it is exactly right.
\end{remark}

\subsection{Circle USYC}\label{sub:cat:rwa:usyc}

An investor subscribes a stablecoin and receives a token equal in number to the shares they hold in a Cayman Islands mutual fund. The fund owes the money; the company the investor deals with is a Bermuda token administrator, and the sponsor disclaims every guarantee to the fund while listing it as a subsidiary in an exhibit to the same annual report whose notes treat it as an unconsolidated variable interest entity. The price is struck once per business day off chain from a prime brokerage account, divided by the supply and published on chain behind a standard aggregator interface, so value accrues in the price and never in the balance. A holder may burn the token at the issuer's teller and be paid at the latest reported price at any hour --- but the teller escrows nothing, an operator pushes liquidity in at the moment of the transaction, a holder who needs more is directed to an email address, and unconditional instant convertibility exists as a separately priced tier.

\begin{measurement}\label{meas:cat:rwa:usyc}
Take $X = \{At,Aw,Ex,Fz,Gp,Ix,Rd,Sh,Sv,Up,Xf\}$, with canonical form
$\mathrm{ex}(X) = \{At,Aw,Ex,Fz,Gp,Ix,Rd,Sh,Sv,Up,Xf\}$, every element being primitive in it.
It is admissible: no requirement term is open, every element is warranted, and it arms no prohibition. It discharges 9 of the 20 recorded obligations; 11 are residue. No composite of at most three corpus protocols equals it; 4 corpus protocols are contained in it, and together they supply every element except \{Sv\}.
It is not minimal: \{At,Ex,Fz,Gp,Sv\} can be removed with every obligation still covered and admissibility intact.
\end{measurement}

\begin{remark}
This is the register-of-record gap in its clearest published form. Ownership sits as shares at a transfer agent and the tokens are the issuer's own \emph{digital twins} of them, with the fund administrator moving shares in the registry \emph{to reflect} movements in the token. Operationally the token leads and the register follows; legally the register is still the record, and a formalisation that treats the token as the ledger misstates what the holder owns, while one that treats it as a mirror misstates where transactions sequence. The residue also swallows the fund's dominant actual use --- pledged at a tri-party custodian as exchange margin, immobilised and still accruing --- along with the identity of every party that holds the assets, which no regulator's record reaches: no adviser files a private-fund schedule for this fund, and the Cayman register is served behind a challenge the fetch could not pass. The mandate over those assets exists in legal form and stops short, naming an agent and a domain and supplying no counterparty limit, no concentration limit and no holdings feed.

The obligations no element discharges are these.
\begin{itemize}\setlength{\itemsep}{0pt}
  \item a subscription arriving after two in the afternoon eastern time is priced at a forecast of the next price report rather than at any observed value
  \item redemptions execute at a different price from subscriptions, with a time-varying spread between the two legs
  \item no reserve stands behind that promise: the teller holds no USDC of its own, liquidity is pushed in by an operator at transaction time, and a holder needing more than the day's posted amount is told to contact the issuer
  \item unconditional instant convertibility exists but must be bought: it is a separately priced per-client tier carrying additional fees
  \item the token's primary use is not investment but margin: it sits immobilised with a tri-party custodian, pledged against derivatives positions at an exchange, while continuing to accrue
  \item a ten per cent performance fee on yield plus four and three basis points on subscription and redemption are taken by the manager off the net asset value
  \item the holder must know who owes them the money, and here the fund owes it while a different company in a different jurisdiction is the party they deal with, and three primary sources disagree about which of the two is the issuer
  \item the holder must know whether the chain is the register, and here the token is explicitly a digital twin of a transfer agent's share register that the administrator moves afterwards to match
  \item the holder must know whether their claim survives the issuer's insolvency, and here they hold an unsecured equity interest with no security agreement, no collateral agent and no perfection step of any kind
  \item the holder must know what the reserve holds and who holds it, and here the assets sit at a prime broker in repo positions whose counterparties are never named and whose composition is never published
  \item the holder must know who may change what the money is invested in, and here the manager sets the mix between bills and repo with no published mandate, band or limit
\end{itemize}
\end{remark}

\begin{remark}
The corpus records this fund and the other treasury fund as identical decompositions while writing down different element sets for them, and the witness makes the difference explicit rather than repairing the arrays. A peg-swap element does not belong here: the price is a rising net asset value and the teller is a net-asset-value mint and redeem with a deliberately asymmetric pair of legs, which is the structural opposite of a par swap. The witness adds an accrual element, because the claim grows in the price and not in the balance; a discretion element, because the issuer determines the price, may delay it, forecasts it after the afternoon cut-off and decides how much instant liquidity exists on a given day; and a cross-domain element, because the token is natively issued on several chains and almost all of it lives on the chain of the single venue that accepts it as margin. Sanctions screening is carried by the external-feed element rather than by the access gate, because the gate is a bit set once at onboarding and the screen is a third party that can invalidate a valid wallet with no act by the issuer, the sponsor or the holder.
\end{remark}

\subsection{Maple Finance}\label{sub:cat:rwa:maple}

A lender deposits a stablecoin into a pool and receives a vault share whose redemption value rises as borrowers repay. The claim runs against one segregated portfolio of a segregated portfolio company, ring-fenced by statute from the company's other portfolios, and recourse is limited to the assets in that portfolio alone. The pool's capital is lent to named institutional borrowers as dated loan contracts on terms agreed off chain, and is also allocated to futures basis trades, other market-neutral strategies and liquidity provision on other protocols, with no published cap, venue whitelist or leverage limit. Value is the pool's assets less its unrealised losses: interest accrues mechanically from the loan schedules, and a loss is recognised when a human declares it and may be reversed by another. Exit is a first-in first-out queue served as liquidity appears, and the protocol publishes both a normal case and a far longer contractual bound.

\begin{measurement}\label{meas:cat:rwa:maple}
Take $X = \{Aw,Ct,Fd,Ft,Gp,Pl,Sh,Sl,Sv,Tg,Up,Wq,Xf,Xm\}$, with canonical form
$\mathrm{ex}(X) = \{Aw,Fd,Ft,Gp,Pl,Sh,Sl,Sv,Tg,Up,Wq,Xf,Xm\}$; the elements \{Ct\} are derived rather than chosen.
It is not admissible: it leaves the requirement term $L1\!:\!Ex{\mid}Tp{\mid}At$ open. It discharges 10 of the 22 recorded obligations; 12 are residue. No composite of at most three corpus protocols equals it; 2 corpus protocols are contained in it, and together they supply every element except \{Ct,Fd,Ft,Pl,Sh,Sl,Sv,Wq\}.
\end{measurement}

\begin{remark}
Maple is the only member of this category whose register of record is the token ledger. There is no transfer agent, no share registry and no off-chain book: the pool contract \emph{is} the ledger of the lender's claim, and what that costs is a recourse limited by statute to one portfolio's assets. The residue holds everything around the claim instead --- that the solvency test runs off chain against price feeds nothing on chain can see and is enforced by an operations team selling through desks, that the delay protecting the upgrade path is set by the same authority it restrains and may be shortened retroactively, that the permission check passes everyone when its criteria were never configured, and that the real instant exit is a third-party automated market maker the protocol does not control. The bounded delegate mandate appears here in its on-chain form and is the most complete instance in the category, and the two parts a checker would have to evaluate are the defective ones: the cap bounds the pool rather than exposure to any borrower, venue or strategy, and the rate-of-change limit is revocable by the party it binds.

The obligations no element discharges are these.
\begin{itemize}\setlength{\itemsep}{0pt}
  \item the solvency test itself runs off chain: an operations team watches three price feeds around the clock, gives a borrower twenty-four hours to top up, and then sells the collateral through over-the-counter desks under the loan agreement
  \item a manager may declare a paper loss that immediately reduces what an exiting lender receives, may reverse it, and the protocol states the purpose is to incentivise current lenders to remain
  \item the first-loss module exists in the code and is deliberately empty, and the staked governance token was never slashable and no longer earns anything
  \item the instant exit is not a protocol facility at all: a lender who will not wait sells the share token into third-party automated market makers the protocol does not control
  \item the permission check returns true by default: if a pool's criteria bitmap was never set, everyone passes, and the protocol documents this as intentional
  \item a pool may be made permissionless once and can never be made permissioned again
  \item the delay protecting those upgrades is set by the same authority the delay is meant to restrain, can be changed at any time, and the change applies retroactively to calls already scheduled
  \item the holder must know who owes them the money, and here the claim runs to a pool contract that holds loan contracts funded to named institutions under bilateral agreements whose governing law and recourse are not published
  \item the holder must know whether the chain is the register, and here it is: the pool contract is the authoritative ledger of the lender's claim, and Maple is the only one of the five whose investor register is on chain
  \item the holder must know whether their claim is perfected against the borrower's estate, and here no security trustee, control agreement or perfection step is evidenced anywhere, only possession of collateral by a contract
  \item the holder must know what backs the claim and who holds it, and here the answer is unusually good: collateral is custodied by the loan contracts on chain and the ratio is published, while the disposal of it is off chain
  \item the holder must know who may change what the money is lent against, and here one delegate selects the counterparties, sets the rate and the size, and no concentration limit is enforced on chain
\end{itemize}
\end{remark}

\begin{remark}
The corpus record carries a staked backstop and protocol-funded emissions, and the witness drops both. The first-loss module exists in the code under the protocol's own heading calling it unused, the governance stake was never slashable and its rewards were sunset by a governance proposal, and both emission channels are closed --- the lender programme's final claim window has passed, and the surviving lender incentives are funded by third parties, which is precisely what protocol-funded emissions are not. With the backstop gone, socialised loss over the pro-rata claim is the only loss-absorption mechanism the protocol has, so the witness adds it, and it is what closes the loss-absorption term the corpus record left open. The witness also adds the cross-domain pair the corpus omitted, and keeps the collateral test as a first-class element rather than a derived one: collateral is posted to and repossessed by the loan contracts on chain, and only the price monitoring and the act of selling are off chain.
\end{remark}

\subsection{Ondo Finance}\label{sub:cat:rwa:ondo}

One sponsor issues claims that are not of the same legal kind, and the difference is larger than the wrapper they share. A subscriber to the yield token is a \emph{lender}: the token evidences the issuer's debt, a named collateral agent holds a first-priority security interest for the holders, and recourse is contractually capped at what that collateral recovers. A subscriber to the fund token becomes a limited partner in a Delaware partnership; a subscriber to the tokenised equities holds a structured note of a British Virgin Islands vehicle governed by Swiss law, whose holders are entitled to damages and not to performance. Price is set a different way in each: a rate that compounds deterministically on chain between administrative changes, a net asset value pushed by hand once a business day from an estimate trued up the next, and a signed off-chain quote that expires in seconds and is consumed by the trade it prices. Redemption is atomic for a holder the registry knows, and otherwise a wire to a non-US bank account, a shared capacity cap, or a diligence check a secondary buyer can fail.

\begin{measurement}\label{meas:cat:rwa:ondo}
Take $X = \{At,Aw,Ex,Fz,Gp,Ix,Ps,Rb,Rd,Rf,Sh,Sv,Tg,Up,Xf,Xm\}$, with canonical form
$\mathrm{ex}(X) = \{At,Aw,Ex,Fz,Gp,Ix,Ps,Rb,Rd,Rf,Sh,Sv,Tg,Up,Xf,Xm\}$, every element being primitive in it.
It is admissible: no requirement term is open, every element is warranted, and it arms no prohibition. It discharges 13 of the 23 recorded obligations; 10 are residue. No composite of at most three corpus protocols equals it; 13 corpus protocols are contained in it, and together they supply every element except \{Rb,Rf,Sv\}.
It is not minimal: \{Ex,Sv,Tg,Up,Xf,Xm\} can be removed with every obligation still covered and admissibility intact.
\end{measurement}

\begin{remark}
The residue is where the sponsor's own record stops agreeing with itself. Its live documentation names two obligors for the same token, and its terms of service keep the Delaware company alive as a distinct contracting party while naming issuers for the other two products and none for that one, so the first thing a holder must know --- who owes them the money --- is settled by no public source. The security interest is asserted and the control agreements are confirmed, and the financing statement that would make the interest good against third parties cannot be checked at all: Delaware publishes no free debtor-name index, and a certified search must be bought from an authorised searcher. That is a documented absence rather than a search failure, and it is this category's problem in its sharpest form --- the wrapper is readable by anyone with an RPC endpoint, and the step that makes the claim survive an insolvency is readable only by whoever pays for it. The one published quantitative bound on discretion anywhere in the category is here, and it bounds the sponsor's credit support rather than the manager's mandate.

The obligations no element discharges are these.
\begin{itemize}\setlength{\itemsep}{0pt}
  \item a subscriber's dollars buy a token that evidences the issuer's debt to the subscriber as a lender, not a proportional interest in any pool
  \item the atomic path is not self-contained: every tokenised-equity mint or redeem requires a live artefact from the issuer's attestation service, and settlement into a stablecoin depends on a swapper holding inventory
  \item the instant channel is capacity-limited, fifty million dollars globally and twenty-five million per investor per rolling twenty-four hours, and is additionally limited by a third party's stablecoin availability
  \item the USDY gate is installed, wired into the transfer hook and deliberately switched off, and one role-holder can re-arm it in a single transaction with no delay
  \item the holder must know who owes them the money and what they can do about it, and here the issuer's own live documentation names two different obligors for the same token while capping recourse at whatever collateral is recovered
  \item the holder must know whether the chain is the register of ownership or a mirror of somebody's book, and across Ondo's three products the answer is three different things
  \item the holder must know whether their claim is secured, perfected, and able to survive the issuer's insolvency, and here a first-priority security interest is asserted while the steps that would make it good are either private or absent
  \item the holder must know what the reserve holds and who holds it, and here the disclosure is a self-published portfolio page with no independent composition attestation, whose fund product is itself two tokenisation layers away from any Treasury
  \item the holder must know who may change what the money is invested in, and here the manager states an unconstrained discretion to change it with no target allocation, no band and no concentration limit disclosed anywhere
  \item over-collateralisation is provided as sponsor equity under a contractual floor tested at quarter end, and there is no second class of token, no subordinated holder and no distribution waterfall
\end{itemize}
\end{remark}

\begin{remark}
The corpus record carries a tranche element for this sponsor and the witness drops it. The junior capital is real --- a contractual floor of sponsor equity above the tokens outstanding, tested at quarter end --- but it is equity standing behind the whole issue rather than a subordinated class of tokens: there is no second token, no junior holder and no distribution waterfall in normal operation, and a priority ordering appears only on enforcement. In its place the witness carries a peg-swap element the issuer names in its own address register, a signed-quote element for a price that arrives as an unforgeable expiring capability rather than as a readable feed, and a discretion element for a rate the issuer administers against a document the holder cannot read. The pro-rata claim element is scoped to the fund token alone; applying it to the debt token would make a creditor into an equity holder, which is the distinction the security agreement exists to preserve.
\end{remark}

\subsection{Aevo}\label{sub:cat:opt:aevo}

An order-book venue in the same lineage as Derive and, the corpus recorded, identical to it. They are not identical: they differ in index construction, in loss-absorption topology, and in whether auto-deleveraging exists at all. This venue has it and its counterpart does not.

\begin{measurement}\label{meas:cat:opt:aevo}
Take $X = \{Ad,Ct,Ex,Fd,Li,Ob,Op,Pf,Rf,Xf,Xm\}$, with canonical form
$\mathrm{ex}(X) = \{Ad,Fd,Ob,Op,Pf,Rf,Xf,Xm\}$; the elements \{Ct,Ex,Li\} are derived rather than chosen.
It is admissible: no requirement term is open, every element is warranted, and it arms no prohibition. It discharges 10 of the 21 recorded obligations; 11 are residue. No composite of at most three corpus protocols equals it; no corpus protocol is contained in it.
It is not minimal: \{Pf,Xf,Xm\} can be removed with every obligation still covered and admissibility intact.
\end{measurement}

\begin{remark}
The residue that remains is the category's own: one symbol carries the whole of options, and the six axes it collapses --- exercise style, settlement in cash or in kind, who underwrites, whether premium is upfront or streamed, how margin is computed, and whether the price comes from a model or a book --- are each unnameable.

The obligations no element discharges are these.
\begin{itemize}\setlength{\itemsep}{0pt}
  \item Exercise is never elected: the option is European and pays automatically at the 08:00 UTC expiry.
  \item The payoff is paid in USDC as a cash difference against the Aevo index; no underlying is delivered.
  \item The premium is paid in full at execution.
  \item Both sides of every option are user accounts; the peer-to-pool layer that gave the protocol its name — the Ribbon DeFi Options Vaults, which sold weekly options into an auction — was exploited on 2025-12-12 for roughly \$2.7M (about 32\% of the affected vault) following an oracle upgrade, and has been halted and is being decommissioned with a compensation mechanism, leaving the exchange unaffected.
  \item Collateral may be posted in several accepted assets under a cross-margin collateral framework, including aeUSD with its own deposit, redemption and composition regime, plus a spot-convert feature.
  \item Under Portfolio Margin the requirement is Scenario Margin plus Floor Margin: fifteen scenarios formed from three spot-movement levels crossed with five IV-shift positions, at maximum spot movement +/-20\% and IV shift +50\% / -25\% in the initial ETH phase, with initial margin 1.25 times maintenance.
  \item Floor Margin = Net Short Exposure(Option) x Unit Floor Margin, at 0.01 ETH per unit in the initial ETH phase, added to the scenario result.
  \item If the book cannot absorb the seized position, the insurance fund takes it at a markup, and a liquidation fee is charged on insurance-fund fills as well.
  \item The three loss-absorption venues escalate in a fixed order — book workout, then insurance fund, then auto-deleveraging — and each transfers the position on a different price basis: limit orders, then a markup, then mark price.
  \item The settlement-price observation window at expiry is not documented anywhere in the exchange's documentation.
  \item Mark prices and per-instrument greeks including implied volatility are computed by the exchange and published through /markets; no pricing model, surface parameterisation or margin-model source is published.
\end{itemize}
\end{remark}

\begin{remark}
This is one of the two rejections in the category that turn out to be decomposition artefacts. The corpus left the loss-absorption term open and the venue was rejected for it; the witness restores the collateral test, which is an order-gating engine rather than a derived consequence, and the auto-deleveraging mechanism, which has its own documentation page and which the corpus omitted. The term closes and the rejection disappears. Nothing about the protocol changed.
\end{remark}

\subsection{Derive}\label{sub:cat:opt:derive}

The only member of the category the corpus decomposition admitted. Options are quoted on a book, margined against a portfolio rather than position by position, and settled in cash against an index. The construction departs from the corpus in one place that matters: the venue's security module is an owner-funded treasury with no stake and no slashing, so the staked-backstop element is not supported and socialised loss carries the obligation instead.

\begin{measurement}\label{meas:cat:opt:derive}
Take $X = \{Ct,Em,Ex,Fd,Ix,Li,Ob,Op,Pf,Rf,Sl,Up,Xf,Xm\}$, with canonical form
$\mathrm{ex}(X) = \{Em,Fd,Ix,Ob,Op,Pf,Rf,Sl,Up,Xf,Xm\}$; the elements \{Ct,Ex,Li\} are derived rather than chosen.
It is admissible: no requirement term is open, every element is warranted, and it arms no prohibition. It discharges 12 of the 22 recorded obligations; 10 are residue. No composite of at most three corpus protocols equals it; no corpus protocol is contained in it.
It is not minimal: \{Fd,Ix,Pf,Xf,Xm\} can be removed with every obligation still covered and admissibility intact.
\end{measurement}

\begin{remark}
What the venue prices --- a volatility surface, and the model that turns it into a quote --- has no element at all, and it is the thing the venue competes on.

The obligations no element discharges are these.
\begin{itemize}\setlength{\itemsep}{0pt}
  \item Exercise is never elected: the option is European, cannot be exercised early, and pays at expiry with no action by the holder.
  \item The payoff is paid in USDC as a cash difference; no underlying is ever delivered.
  \item The payoff is fixed against a 30-minute, one-minute-sampled TWAP of the index (SETTLEMENT\_TWAP\_PERIOD = 30), and over that final window the option's mark migrates linearly from the forward price into the running average.
  \item Each feed carries a confidence score between 0 and 1, and low confidence raises the account's initial margin through an Oracle Contingency term.
  \item Implied volatility is supplied as a five-parameter SVI curve per expiry, with log-moneyness bounded at four standard deviations, total implied volatility capped at 24.0 and total variance at 144.0, produced by Block Scholes and posted on-chain.
  \item Order entry, matching and ordering are performed by a centralized operator; only settlement, collateral and liquidation are on-chain, and a user may exit unilaterally through a block-explorer escape hatch if the matching engine or interface fails.
  \item The premium is paid in full at execution, as a single transfer, and never streams.
  \item Under Portfolio Margin 2 an account's maintenance margin is portfolio mark-to-market plus a maxLoss taken as the minimum over four scenario families — 23 joint spot/volatility shocks over spot -18\%..+18\% in 4.5\% steps, a damped tail family at roughly +/-20-25\%, two skew scenarios and a forward-basis family at +/-4.5\% — plus collateral-haircut, perp, option and oracle contingencies; and a PM account may hold derivatives on only one base asset, under per-asset open-interest caps.
  \item If the auction leaves the account short, the Security Module pays the deficit from an owner-funded subaccount that a whitelisted set of modules may draw on.
  \item Instruments are listed on a rolling calendar — 2 daily, 3 weekly and 3 monthly expiries, auto-replaced on expiry — with strike increments that are a function of the option's delta and time to expiry (ETH 40-60 delta: \$25 at under a day rising to \$200 beyond 35 days).
\end{itemize}
\end{remark}

\begin{remark}
Among the alternatives the loss-absorption term leaves open, the witness takes the one the code supports rather than the one the corpus recorded. That is the whole content of the departure: an element is carried when a mechanism answers to it, and a treasury that cannot be slashed does not answer to a bond.
\end{remark}

\subsection{Hegic}\label{sub:cat:opt:hegic}

A peer-to-pool venue: buyers take the other side of a single pool, and the pool underwrites every position. Like Rysk it closes its solvency obligation at write time, by an inequality on the pool's balance sheet rather than by a test on an individual position.

\begin{measurement}\label{meas:cat:opt:hegic}
Take $X = \{Bs,Ep,Ex,Fd,Gp,Op,Sh\}$, with canonical form
$\mathrm{ex}(X) = \{Bs,Ep,Ex,Fd,Gp,Op,Sh\}$, every element being primitive in it.
It is not admissible: it leaves the requirement term $L1\!:\!Ct$ open. It discharges 7 of the 15 recorded obligations; 8 are residue. No composite of at most three corpus protocols equals it; no corpus protocol is contained in it.
\end{measurement}

\begin{remark}
A balance-sheet inequality checked when a contract is written and a threshold test evaluated while it is open are different mechanisms with the same purpose. The vocabulary has a name for the second and none for the first, so a venue that cannot become insolvent in the relevant sense is recorded as one that has not said how it avoids it.

The obligations no element discharges are these.
\begin{itemize}\setlength{\itemsep}{0pt}
  \item The Hegic Operational Treasury is the writer of every position; there is no counterparty user anywhere in the system.
  \item At purchase the treasury locks negativePNL — the maximum possible payoff of that position — out of its USDC and adds the premium to lockedPremium; the reservation is released by unlock at expiry.
  \item Every write is admitted only if totalLocked + lockedPremium <= totalBalance() + coverPool.availableForPayment(), checked on the write itself, and refused otherwise.
  \item DISCHARGED BY CONSTRUCTION: there is no margin engine, no health factor, no liquidator and no liquidation path for any position on either side.
  \item The buyer realises the payoff by calling payOff before expiry, which pays the strategy's net P\&L in USDC; exercise is permitted only inside exerciseWindowDuration before expiry, and an in-the-money position not claimed in that window pays nothing, its reservation returning to the treasury.
  \item The premium is computed on-chain at purchase by that instrument's own PriceCalculator: either a single admin-set impliedVolRate scalar, or five admin-set int256 polynomial coefficients over the holding period; strikes are scaled to spot at purchase by a priceCorrectionRate.
  \item Composed payoffs — straddles, strangles, straps, strips, spreads, inverse butterflies and inverse condors — are each their own deployed contract with its own pricer and its own limit, rather than compositions a user assembles from legs.
  \item DEFAULT\_ADMIN\_ROLE can withdraw an arbitrary amount from the Operational Treasury to an arbitrary address, set the implied-volatility scalar and the polynomial coefficients, set period limits and the benchmark, set the next epoch's conversion price, set the payoff pool and close epochs — with no timelock in the code read, and against plain non-proxy contracts.
\end{itemize}
\end{remark}

\begin{remark}
The narrower of the two by-construction cases, and the difference is instructive. This venue does carry a first-loss layer --- staked capital that absorbs losses ahead of the pool --- so the loss-absorption term closes and only the collateral-test term stays open. One rejection where the corpus recorded two, and the remaining one is again an obligation that cannot arise.
\end{remark}

\subsection{Panoptic}\label{sub:cat:opt:panoptic}

A venue that manufactures option-like payoffs out of concentrated liquidity positions rather than writing contracts, so its collateral, its pricing and its liquidation all run through an automated market maker. It is described as oracle-free, which is true of external feeds and false of prices.

\begin{measurement}\label{meas:cat:opt:panoptic}
Take $X = \{Cl,Ct,Fd,Gp,Ix,Li,Op,Pl,Sh,Sl,Sr,Tp\}$, with canonical form
$\mathrm{ex}(X) = \{Cl,Fd,Gp,Ix,Li,Op,Pl,Sh,Sl,Sr,Tp\}$; the elements \{Ct\} are derived rather than chosen.
It is admissible: no requirement term is open, every element is warranted, and it arms no prohibition. It discharges 13 of the 23 recorded obligations; 10 are residue. No composite of at most three corpus protocols equals it; no corpus protocol is contained in it.
It is not minimal: \{Sh\} can be removed with every obligation still covered and admissibility intact.
\end{measurement}

\begin{remark}
That an option can be assembled from a liquidity range, rather than written as a contract, is the venue's whole design and has no expression here. The vocabulary has a symbol for the payoff and a symbol for the range, and no way to say that one is being built out of the other.

The obligations no element discharges are these.
\begin{itemize}\setlength{\itemsep}{0pt}
  \item A position has no strike and no expiry: it is a tick range, up to four legs, encoded in one ERC-1155 tokenId, and it never expires and is never exercised at a price.
  \item The payoff is realised continuously as the underlying range position's composition changes with price; closing a position burns it and returns tokens, and there is no terminal fixing to determine.
  \item The rate of streamia is the Uniswap swap fee the removed liquidity would have generated, multiplied by a spread multiplier that grows with how much of that specific chunk has been removed (VEGOID = 8, nu = 1/VEGOID).
  \item The counterparty to a buyer is whoever's liquidity chunk was removed, the funding for both sides comes from a shared pool, and the payoff shape is the AMM's own.
  \item Collateral requirements scale with pool utilisation up to SATURATED\_POOL\_UTIL = 9,000,000, and leverage of up to 10x buying and 5x selling holds only under normal conditions and changes in response to pool utilisation.
  \item A shortfall remaining after liquidation costs and bonus is absorbed first by a haircut on the account's unsettled premia.
  \item A seller may force-exercise a buyer's out-of-range long position to reclaim the liquidity, paying the exercised user a fee that decays exponentially with distance from strike (FORCE\_EXERCISE\_COST = 30,000 scaled units, floored at ONE\_BPS = 1000).
  \item PanopticFactoryV3 and PanopticFactoryV4 will open an options market on any Uniswap V3 or V4 pool permissionlessly, on any ERC20 pair.
  \item Every user action funnels through dispatch() / dispatchFrom(), which will revert when the current tick is further than MAX\_TWAP\_DELTA\_DISPATCH from the TWAP tick.
  \item The protocol lends against, and prices with, the same Uniswap pool its instruments are written on, and mitigates the resulting manipulation channel with a median ring buffer, a four-EMA ensemble, a clamp on median movement, the dispatch TWAP-delta bound and safe mode.
\end{itemize}
\end{remark}

\begin{remark}
The second artefact. The corpus carried no price source for this venue at all, having read oracle-free as price-free, and demoted its collateral test --- a substantial audited risk engine --- to a derived consequence. The witness restores both: an on-chain time-weighted price is a price, and a risk engine is a threshold test. The rejection disappears with them.
\end{remark}

\subsection{Circle USDC}\label{sub:cat:fiat:usdc}

The most disclosed issuer of the five: a public company, a monthly reserve report, a named auditor, and cross-chain issuance by burn and mint rather than by wrapping. Its on-chain surface is nonetheless the same shape as the others'.

\begin{measurement}\label{meas:cat:fiat:usdc}
Take $X = \{At,Au,Aw,Fz,Gp,Rd,Up,Xf,Xm\}$, with canonical form
$\mathrm{ex}(X) = \{At,Au,Aw,Fz,Gp,Rd,Up,Xf,Xm\}$, every element being primitive in it.
It is admissible: no requirement term is open, every element is warranted, and it arms no prohibition. It discharges 9 of the 24 recorded obligations; 15 are residue. No composite of at most three corpus protocols equals it; 3 corpus protocols are contained in it, and together they supply every element except \{Au,Fz\}.
It is not minimal: \{Xf,Xm\} can be removed with every obligation still covered and admissibility intact.
\end{measurement}

\begin{remark}
The obligation the token \emph{is} --- who may redeem, at what minimum, at what fee, with what latency --- is stated in the issuer's terms and nowhere in the model. That is the category's whole point.

The obligations no element discharges are these.
\begin{itemize}\setlength{\itemsep}{0pt}
  \item USDC is a transferable receipt for a claim the issuer asserts is a property claim rather than a contractual one: Circle states it holds only bare legal title to the reserve and no beneficial interest or property rights, so the reserve should not be considered property of its bankruptcy estate -- while conceding in the same filing that no court has tested this, that holders could face automatic-stay delays even if it holds, and that a US or French court could decide otherwise.
  \item The issuers of record are Circle Internet Financial, LLC for US holders and Circle Internet Financial Europe SAS for EEA holders, so which entity owes the holder is decided by the holder's residence.
  \item The issuer may decline to process any issuance or redemption without prior notice, may block addresses and freeze the USDC temporarily or permanently, and warns that a blocked holder may forfeit any rights associated with the USDC including the ability to redeem.
  \item The reserve at 2026-06-30 is US\$73,344,909,176 against US\$73,268,560,097 of USDC in circulation, of which US\$61,917,347,895 sits in the Circle Reserve Fund -- a Rule 2a-7 government money market fund managed by BlackRock Advisors, LLC, custodied at BNY, available only to Circle, whose equity Circle Internet Financial, LLC holds 100\% of 'on behalf of USDC holders' -- holding nine CUSIPs of Treasuries maturing 2026-07-02 to 2026-08-20, US\$52.53bn of overnight Treasury repo and US\$1.00bn of cash.
  \item The remaining US\$11,382,256,398 is cash held at regulated financial institutions in for-benefit-of accounts, primarily with banks designated by the Financial Stability Board as global systemically important banks; no individual bank is named in the reserve report.
  \item Circulation is not supply: it is total supply across 36 named USDC Approved Blockchains less Tokens Allowed But Not Issued (1,380,790,392 at 2026-06-30, an artefact of how USDC is implemented on Algorand, Hedera, Polkadot Asset Hub and Solana), less Access Denied Tokens (123,219,920), less Circle Gateway pending burns (94,358), with a further 993,225 permanently frozen on the deprecated FLOW deployment.
  \item An access denial does not extinguish the liability: the tokens remain owed until the segregated reserve funds are transferred to the relevant law enforcement agency, or until the denial is reversed and a redemption is made.
  \item Mint authority is a delegated QUANTITY rather than a role: masterMinter() is 0xE982615d461DD5cD06575BbeA87624fda4e3de17, a 7,667-byte contract named MinterAdmin, and the implementation exposes configureMinter, removeMinter and minterAllowance, so each minter holds an allowance that minting consumes.
  \item The proxy admin is 0x807a96288A1a408dBC13DE2b1d087d10356395d2, an externally owned account with no code, and there is no timelock anywhere in the control plane.
  \item Authority is genuinely divided across five distinct roles -- owner 0xFcb19e6a322b27c06842A71e8c725399f049AE3a, masterMinter, blacklister, pauser and a rescuer -- which is the only real role separation among the five; every one of them resolves to a single externally owned account or, for masterMinter, to a contract.
  \item Circle Gateway, launched July 2025, provides a unified USDC balance instantly accessible across supported blockchains, holding balances in Circle smart contracts and producing the pending-burns deduction in the reserve report.
  \item The issuer is licensed by many authorities and supervised as a payments and virtual-asset business rather than as a bank: a NYDFS BitLicense, state money-transmitter licences and FinCEN MSB registration in the United States; an AMF DASP registration and an ACPR electronic-money-institution licence in France; an ADGM FSRA licence; and a Bermuda Monetary Authority DABA licence.
  \item In December 2025 the issuer received OCC conditional approval to establish Circle National Trust, which is to manage the USDC reserve and hold a first-priority perfected security interest in it as collateral trustee for the benefit of USDC holders.
  \item The float accrues entirely to the issuer and its distributors: FY2025 reserve income was US\$2,636,822 thousand and distribution and transaction costs were US\$1,661,549 thousand, so 63\% of the reserve yield was paid to distribution partners, principally Coinbase under a Collaboration Agreement with an initial three-year auto-renewing term; holders receive nothing.
  \item On-chain the token is plain balance accounting: 6 decimals, version() returns '2', currency() returns 'USD', and Ethereum supply is 49,536,777,415.220915 at 2026-08-04 against attested all-chain circulation of 73,268,560,097.
\end{itemize}
\end{remark}

\begin{remark}
The witness adds a delegated-execution element the corpus omits, because the contract implements signed transfer authorisations that satisfy every term of the governing requirement, and drops the peg-swap element for the reason given above. Disclosure quality does not change the decomposition; it changes only what a reader can check.
\end{remark}

\subsection{Global Dollar USDG}\label{sub:cat:fiat:usdg}

The only issuer of the five that runs a timelock, and the only one whose reward distribution happens on-chain rather than as an off-chain arrangement. Both facts are absent from the corpus record.

\begin{measurement}\label{meas:cat:fiat:usdg}
Take $X = \{At,Au,Aw,Ep,Fd,Fz,Gp,Rd,Sh,Tg,Up\}$, with canonical form
$\mathrm{ex}(X) = \{At,Au,Aw,Ep,Fd,Fz,Gp,Rd,Sh,Tg,Up\}$, every element being primitive in it.
It is admissible: no requirement term is open, every element is warranted, and it arms no prohibition. It discharges 10 of the 23 recorded obligations; 13 are residue. No composite of at most three corpus protocols equals it; no corpus protocol is contained in it.
It is not minimal: \{Ep,Fd,Sh\} can be removed with every obligation still covered and admissibility intact.
\end{measurement}

\begin{remark}
And here the category's sharpest fact appears, which is a fact about a \emph{party}. A single externally-owned account is the freeze key for this issuer and for one other, and is simultaneously the sole proposer, executor and canceller of the timelock --- which does not cover the freeze. The vocabulary records that a timelock is present and that a freeze is present. It cannot record that they are the same key, that the key is one account rather than a quorum, or that the delay does not apply to the power that matters. This is the party sort, exhibited rather than conjectured.

The obligations no element discharges are these.
\begin{itemize}\setlength{\itemsep}{0pt}
  \item USDG is an electronic-money claim with two obligors selected by the holder's residence: for EEA holders the issuer is Paxos Issuance Europe Oy, a Finnish electronic money institution supervised by the Finnish Financial Supervisory Authority under LEI 743700KYSSTKZYGEUF50 and Business ID 3236886-2; for everyone else the issuer is Paxos Digital Singapore Pte. Ltd., a Major Payments Institution supervised by the Monetary Authority of Singapore.
  \item Each supply controller's minting is rate-limited on-chain: SupplyControl carries per-controller limit configuration through contracts/lib/RateLimit.sol, with updateLimitConfig and getRemainingMintAmount, under SUPPLY\_CONTROLLER\_MANAGER\_ROLE, SUPPLY\_CONTROLLER\_ROLE and TOKEN\_CONTRACT\_ROLE.
  \item The reserve is at least 100\% of outstanding EU-attributed USDG in high-quality, low-risk assets such as cash and short-term government debt, legally segregated from the issuer's own funds and from the custodian's funds, held with qualified custodians, authorised by the Finnish Financial Supervisory Authority and structured for bankruptcy remoteness, with MiCA additionally requiring a portion to sit with European banking partners.
  \item In insolvency holders are entitled to a proportionate share of the segregated reserve assets, with their claims taking precedence over unsecured creditors, distributed under Finnish insolvency law and overseen by the insolvency administrator and the Finnish Financial Supervisory Authority.
  \item USDG is not covered by any investor compensation or deposit guarantee scheme: white paper field D.11 is false for both Directive 97/9/EC and Directive 2014/49/EU, though the underlying segregated accounts are.
  \item The governing law is the law of Finland and the competent court is the District Court of Helsinki.
  \item Reward entitlement survives a transfer of the token: updateSharesWithRewardPreservation, \_transferWithDifferentPayouts and \_updateWalletWithPayoutGroup re-attribute reward shares between payout groups when balances move.
  \item The money funding the on-chain rewards ledger is off-chain reserve income paid under partner agreements on three distinct bases -- average USDG balance held, net USDG minted, and USDG transaction volume -- with an explicit disclaimer that nothing displayed constitutes a commitment or obligation to pay any particular amount.
  \item Beneath the proxy the token carries a second, finer mutability surface: a facet router with setFacet(bytes4,address) and batchSetFacet(FacetCut[]) under DEFAULT\_ADMIN\_ROLE, a facets mapping, and a fallback() that delegatecalls facets[msg.sig] and reverts with FacetNotFound when the selector is unmapped.
  \item A separate and much shorter delay governs handover of the admin role itself: defaultAdminDelay() returns 10800, three hours, on transferring DEFAULT\_ADMIN\_ROLE, alongside beginDefaultAdminTransfer, acceptDefaultAdminTransfer, cancelDefaultAdminTransfer, changeDefaultAdminDelay and rollbackDefaultAdminDelay.
  \item USDG is deployed on Ethereum as an ERC-20, Solana as an SPL Token-2022, and Ink, X Layer and Robinhood Chain as ERC-20s, on distributed ledgers the white paper confirms the issuer does not operate.
  \item Ethereum supply is 500,509,895.398634 USDG at 6 decimals on 2026-08-04, a fraction of all-chain supply.
  \item The contracts are published under MIT licence and audited: github.com/paxosglobal/usdg-contract and github.com/paxosglobal/paxos-token-contracts carry an audits directory holding Halborn on PaxosTokenV2 and Zellic on USDG Rewards and PAXG V2, with the USDG README naming Zellic and Trail of Bits.
\end{itemize}
\end{remark}

\begin{remark}
The witness adds the timelock, the on-chain reward ledger with its share accounting and epoch gate, and the guardian --- elements the corpus assigned to nobody in this category, one of them established by reading the live facet table rather than the documentation.
\end{remark}

\subsection{PayPal USD}\label{sub:cat:fiat:pyusd}

An issuer whose token is a thin contract over a large regulated arrangement, and whose supervisory regime changed during the period the corpus covers.

\begin{measurement}\label{meas:cat:fiat:pyusd}
Take $X = \{At,Au,Aw,Fz,Gp,Rd,Up\}$, with canonical form
$\mathrm{ex}(X) = \{At,Au,Aw,Fz,Gp,Rd,Up\}$, every element being primitive in it.
It is admissible: no requirement term is open, every element is warranted, and it arms no prohibition. It discharges 8 of the 22 recorded obligations; 14 are residue. No composite of at most three corpus protocols equals it; no corpus protocol is contained in it.
It is not minimal: \{Rd\} can be removed with every obligation still covered and admissibility intact.
\end{measurement}

\begin{remark}
This issuer shares its freeze key with another in the same category, held by one account. Two issuers with different obligors, different regulators and different reserves have a single common point of control, and the model has no way to say so.

The obligations no element discharges are these.
\begin{itemize}\setlength{\itemsep}{0pt}
  \item PayPal is not the issuer: PYUSD is issued by Paxos Trust Company, N.A., and PayPal's own 10-Q calls Paxos 'a third-party issuer (the PYUSD Issuer)' and states the claim precisely -- each token of PYUSD held by PayPal represents a contractual right to redeem with the third-party issuer of PYUSD for one U.S. dollar.
  \item The obligor is not the counterparty: a PayPal or Venmo consumer's counterparty is PayPal under PayPal's user agreement, while the obligor of last resort is Paxos under Paxos's stablecoin terms -- two hops, two contracts, one token.
  \item The dominant use path never touches a blockchain: most PYUSD activity is a row change in PayPal's or Venmo's custodial database, and the token contract is a settlement layer for a payment network whose real ledger is invisible to it.
  \item Reserves are US dollar deposits, US treasuries and cash equivalents held in segregated, bankruptcy-remote accounts at a national trust bank.
  \item The issuer's own publications do not agree on who supervises it: paxos.com/pyusd states the issuer is Paxos Trust Company, N.A. subject to strict regulatory oversight by the Office of the Comptroller of the Currency, while the PYUSD transparency page still describes 'Paxos' prudential regulator, the New York State Department of Financial Services' and its Dollar-Backed Stablecoins guidance as the body approving the examiner's appointment.
  \item A second regulator sits over the distributor rather than the issuer: PayPal is regulated by NYDFS as a virtual currency business and is a digital asset service provider under the GENIUS Act.
  \item The token's default admin -- which holds reclaimToken(), setSupplyControl() and role administration -- is a bare externally owned account, owner() equal to defaultAdmin() equal to 0x3aF3E85f4f97DE7Ad0F000B724fB77Fe5FFc024b, and the only delay in the system, defaultAdminDelay() of 10800 seconds, applies to TRANSFERRING that role and not to using it.
  \item PYUSD shares the Paxos implementation base with USDG -- PaxosTokenV2 over BaseStorage, EIP2612, EIP3009 and AccessControlDefaultAdminRulesUpgradeable, with the same SupplyControl and RateLimit machinery -- but has NO rewards module and NO facet router.
  \item PYUSD is deployed on five documented mainnets -- Ethereum, Solana, Arbitrum, Polygon PoS and X Layer -- each with its own token contract AND its own Supply Control contract address, which is the signature of independent native issuance rather than a bridged representation.
  \item With the cross-domain elements withdrawn, PYUSD's on-chain behaviour reduces to transfer accounting, a permission gate on issuance, a freeze-and-wipe power, a pause, a signed-authorisation family and an upgradeable implementation behind a single key.
  \item PayPal pays a variable-rate PYUSD reward to eligible US users who hold a PYUSD balance in the PayPal app, funded by PayPal as distributor rather than by Paxos as issuer or by any on-chain protocol.
  \item The same instrument is reversible on one side of the custodial boundary and final on the other: dispute reversal and chargeback exist as primitives in PayPal's ledger and cannot exist in the token.
  \item PayPal itself holds PYUSD on its own balance sheet, carried at amortised cost and classified alongside cash and cash equivalents.
  \item On-chain the token is plain balance accounting: 6 decimals and Ethereum supply 1,832,790,045.376448 at 2026-08-04, against the corpus's US\$2.68B all-chain figure.
\end{itemize}
\end{remark}

\begin{remark}
The witness removes the cross-domain elements the corpus assigns, because no primary source supports them for this issuer. Removing an unsupported element is the same discipline as adding a supported one, and it is what makes the remaining identity claims in this category mean anything.
\end{remark}

\subsection{Tether USDT}\label{sub:cat:fiat:usdt}

The largest asset in the corpus. The contract is small and old: a freeze list, a mutable implementation, an attestation the issuer publishes and a redemption right available to a whitelisted few. Everything determining whether the token is worth a dollar sits outside it.

\begin{measurement}\label{meas:cat:fiat:usdt}
Take $X = \{At,Aw,Fz,Gp,Rd,Up\}$, with canonical form
$\mathrm{ex}(X) = \{At,Aw,Fz,Gp,Rd,Up\}$, every element being primitive in it.
It is admissible: no requirement term is open, every element is warranted, and it arms no prohibition. It discharges 6 of the 22 recorded obligations; 16 are residue. No composite of at most three corpus protocols equals it; no corpus protocol is contained in it.
It is not minimal: \{Aw\} can be removed with every obligation still covered and admissibility intact.
\end{measurement}

\begin{remark}
A second pass established two facts about this issuer that no element can hold. The freeze list has been used against thousands of addresses and has destroyed a substantial balance outright; and in almost nine years the fee parameters and the pause have never once been exercised. A power used constantly and a power never used are recorded identically, because the vocabulary names the facility and not its history.

The obligations no element discharges are these.
\begin{itemize}\setlength{\itemsep}{0pt}
  \item USDT is a transferable receipt for an unsecured, on-demand claim against a single private company, Tether International, S.A. de C.V., incorporated in El Salvador and authorised by the Comision Nacional de Activos Digitales under DASP registration PSAD-0028; the attested accounts classify the outstanding tokens as a refund liability under IFRS 9.
  \item The claim is to face value only: holders are not entitled to any increase in the value of the reserves in excess of the face value of the tokens, so the reserve surplus of US\$4,109,529,196 at 2026-06-30 belongs to the shareholder and the holder is senior at par with no upside.
  \item The issuer may delay or suspend redemption where it determines a Prohibited Use, where directed by any government, where the account is under investigation, or on suspicion of fraud, and maintains sole discretion to approve or reject applications to become a customer at all.
  \item The reserve at 2026-06-30 is US\$187,751,426,411 of total assets against US\$183,641,897,215 of liabilities, of which US\$183,622,105,630 is digital tokens issued: US\$114.96bn of Treasury bills under 90 days, US\$25.62bn of repo, US\$18.84bn of LBMA gold bars, US\$5.80bn of bitcoin, US\$3.76bn of public equities, US\$5.24bn of other investments and US\$13.45bn of secured loans.
  \item Cash and bank deposits are US\$40,307,440, which is 0.02\% of the reserve, and roughly US\$23.6bn (12.6\%) is gold, bitcoin, equities, other investments and secured loans -- assets whose realisable value in a redemption run is a judgement.
  \item No bank, depositary or custodian is named anywhere in the primary record: the offering document says only that the issuer 'primarily holds the Reserves with various third parties including banks and licensed financial institutions', and the assurance report names none.
  \item The attester expressly declines the solvency judgement: no assurance is given on the going-concern assessment, the Notes are outside the engagement, and the valuations 'do not reflect unexpected and extraordinary market conditions, or the case of key custodians or counterparties experiencing substantial illiquidity'.
  \item The regulator is El Salvador's CNAD; disputes are routed to confidential BVI-law arbitration seated in London before a sole arbitrator; there is no deposit insurance, no trust, no stated security interest and no segregation covenant.
  \item On-chain the token is plain balance accounting: TetherToken at 0xdAC17F958D2ee523a2206206994597C13D831ec7, Solidity 0.4.18, 6 decimals, Ethereum supply 92,056,579,574.486405 at 2026-08-04, inheriting Pausable, StandardToken and BlackList from a single Ownable with no role separation whatsoever.
  \item Issuance and destruction are treasury operations, not user-facing: issue(uint) and redeem(uint) are onlyOwner and mint to and burn from the owner's own balance; no contract function observes, values or is constrained by the reserve, and the contract contains no reference to a reserve.
  \item The trigger for freezing is a law-enforcement, regulatory or government request, and the issuer reserves the right to freeze tokens in external wallets for which it does not hold private keys; the issuer claims cooperation with more than 340 agencies in 65 countries and more than US\$4.4bn frozen since launch, including more than US\$344M across two addresses in coordination with OFAC on 2026-04-23.
  \item The owner can impose a fee on every transfer of the instrument: setParams(uint newBasisPoints, uint newMaxFee) sets a per-transfer basis-point rate capped by maximumFee; both read zero at 2026-08-04.
  \item Mint, burn, freeze, destroy, pause, the transfer-fee dial and the upgrade are all onlyOwner and owner() is one address, 0xC6CDE7C39eB2f0F0095F41570af89eFC2C1Ea828, a MultiSigWallet with 6 owners and required() == 3, with no timelock of any kind.
  \item Supply moves between Ethereum, Tron, Solana, Avalanche, TON and other chains by the issuer burning on one and issuing on another under off-chain authorisation, over a chain set determined by the issuer alone; there is no bridge contract, no message and no verifier in this path.
  \item The issuer's answer to United States regulation was to create a different token with a different obligor: USAT, launched 2026-01-27, issued by Anchorage Digital Bank under OCC supervision with Cantor Fitzgerald as reserve custodian under the GENIUS Act, with separate reserves, separate issuance and separate redemption rails from USDT.
  \item There is no canonical public source repository: the authority for USDT-on-Ethereum is the verified bytecode, verified on Blockscout on 2019-04-18 as a partial (metadata-hash) match with no license declared, and the only GitHub artefact is an unofficial three-commit mirror that itself points back to a block explorer.
\end{itemize}
\end{remark}

\begin{remark}
The witness drops the peg-swap element the corpus assigns, because no on-chain module performs a par exchange against a reserve; redemption is a legal process conducted off-chain with the issuer. An element is carried when a mechanism answers to it, and here nothing on-chain does.
\end{remark}

\subsection{USD1}\label{sub:cat:fiat:usd1}

The newest issuer in the corpus, and the one the corpus recorded as decomposing identically to Tether. It does not: it performs genuine burn-and-mint issuance across domains, which the corpus record omits entirely.

\begin{measurement}\label{meas:cat:fiat:usd1}
Take $X = \{At,Au,Aw,Fz,Gp,Rd,Up,Xf,Xm\}$, with canonical form
$\mathrm{ex}(X) = \{At,Au,Aw,Fz,Gp,Rd,Up,Xf,Xm\}$, every element being primitive in it.
It is admissible: no requirement term is open, every element is warranted, and it arms no prohibition. It discharges 9 of the 23 recorded obligations; 14 are residue. No composite of at most three corpus protocols equals it; 3 corpus protocols are contained in it, and together they supply every element except \{Au,Fz\}.
It is not minimal: \{Aw,Rd,Xf,Xm\} can be removed with every obligation still covered and admissibility intact.
\end{measurement}

\begin{remark}
What survives the correction is the more interesting claim. Two issuers can differ in every fact that determines whether a holder is paid --- obligor, regulator, reserve, custodian --- and still differ by nothing the vocabulary records until a cross-domain path is added.

The obligations no element discharges are these.
\begin{itemize}\setlength{\itemsep}{0pt}
  \item The obligor is not the brand: the June 2026 examination states that USD1 is issued and redeemed by BitGo Bank \& Trust, N.A., a national trust bank chartered and supervised by the Office of the Comptroller of the Currency, while the USD1 brand and certain associated trademarks are owned and controlled by World Liberty Financial, Inc. and SC Financial Technologies, LLC.
  \item The claim is contractual and confers no interest in any asset: 'USD1 holders' redemption rights are contractual and do not constitute a security interest in, or direct property interest in, any specific reserve asset.'
  \item No redemption minimum, fee schedule or settlement time is published: the governing terms describe only submission through the issuer's platform, compliance review, and settlement timing.
  \item The reserve at 2026-06-30 is US\$4,634,586,966 against 4,634,410,387 redeemable tokens outstanding -- a surplus of US\$176,579 -- held as US\$1,178,613,044 in demand deposit accounts and US\$3,455,973,922 in a government money market fund at NAV under CUSIP 31607A703, with zero in the issuer's own bank account.
  \item Cash sits in demand-deposit and money-market-deposit accounts at US commercial banks regulated by the Federal Reserve, in segregated accounts titled to BitGo Bank \& Trust, N.A. for the benefit of USD1 stablecoin holders, with the attestation noting that balances at times may exceed the FDIC limit of US\$250,000; no bank is named.
  \item Eligible reserve assets are contractually limited to cash, cash equivalents, short-term US Treasuries, fully Treasury-collateralised reverse repurchase agreements and government money market funds.
  \item The issuer is also the custodian: BitGo Bank \& Trust, N.A. issues, redeems and holds the reserve.
  \item The seizure machinery is armed and, at the report date, unused: the June 2026 report records that with respect to USD1 there are no time-locked tokens, test tokens, or access-restricted tokens in the reporting period.
  \item Upgrade authority and token authority are held by two different quorums with no delay on either: the ProxyAdmin's owner is a Gnosis Safe at 0x6a8dc6dbf909f542f2536edf8676d52630cf59e1 with 5 owners and threshold 3, while the token's owner is TokenGovernor at 0xEE9b1a09aEdaCED9dCDa74964EA447FEB93861c2 whose DEFAULT\_ADMIN is a different Safe at 0x0d190b74308669e8f4fe7dbce3466169b285289d with 6 owners and threshold 3.
  \item Token authority is decomposed into nine named roles in TokenGovernor -- MINTER\_ROLE, BRIDGE\_MINTER\_OR\_BURNER\_ROLE, BURNER\_ROLE, FREEZER\_ROLE, UNFREEZER\_ROLE, PAUSER\_ROLE, UNPAUSER\_ROLE, RECOVERY\_ROLE and CHECKER\_ADMIN\_ROLE -- with DEFAULT\_ADMIN\_ROLE reserved for drainFrozenAccount, batchDrainFrozenAccounts and executeTokenFunction; minters and burners are the admin Safe plus two further addresses, freezers the admin Safe plus one, pausers the admin Safe plus one.
  \item A compliance checker can be installed at any time and is not installed: TokenGovernor exposes setChecker() for an IChecker, and getChecker() reads the zero address at 2026-08-04.
  \item Ownership cannot be dropped: renounceOwnership() reverts.
  \item Bridge-driven supply is a different authority from treasury-driven supply: BRIDGE\_MINTER\_OR\_BURNER\_ROLE is separate from MINTER\_ROLE and BURNER\_ROLE and is held by different addresses.
  \item The token has 18 decimals rather than the category's usual 6, and Ethereum supply is 1,529,515,526.739243049392059492 at 2026-08-04 against attested all-chain outstanding of 4,634,410,387 across six chains, of which BNB Smart Chain is the larger leg at 1,802,654,137.
\end{itemize}
\end{remark}

\begin{remark}
The witness adds the two cross-domain elements on evidence of a live canonical-issuance path and a named bridge role. The corpus identity claim depended on their absence, and it does not survive them.
\end{remark}

\subsection{Azuro}\label{sub:cat:prd:azuro}

A bettor stakes collateral on an outcome at odds quoted and locked at the moment of the bet, and holds the position as a transferable token. There is no counterparty on the other side: a single protocol-wide pool underwrites every bet, and the odds move against incoming flow because each bet shifts the virtual funds from which they are computed. The payout is fixed at bet time and claimed once the event resolves; a cancelled event refunds the stake instead, and the pool will also quote a price to buy an open bet back before the event concludes. The party that priced the event also resolves it --- a permissioned data provider creates the event, seeds and reprices the odds, sets the risk budget, and then declares which outcome won, resolving in good faith against a bookmaker's published settlement rules. The protocol's documentation names its DAO as arbiter of last resort in a dispute and specifies no mechanism for that arbitration anywhere on-chain, so the dispute path exists on paper and nowhere else.

\begin{measurement}\label{meas:cat:prd:azuro}
Take $X = \{Au,Aw,Fd,Gp,Gs,Rd,Rl,Up\}$, with canonical form
$\mathrm{ex}(X) = \{Aw,Fd,Gp,Gs,Rd,Rl,Up\}$; the elements \{Au\} are derived rather than chosen.
It is admissible: no requirement term is open, every element is warranted, and it arms no prohibition. It discharges 7 of the 19 recorded obligations; 12 are residue. No composite of at most three corpus protocols equals it; no corpus protocol is contained in it.
It is not minimal: \{Gs,Gp\} can be removed with every obligation still covered and admissibility intact.
\end{measurement}

\begin{remark}
What the witness cannot say is almost everything that makes this a betting protocol rather than an exchange. A pool that is principal to every trade, writing a non-linear event-contingent payoff, has no element; nor has a curve that prices a \emph{state} against the pool's own exposure rather than an asset against an asset; nor has a payoff quoted and locked for one bettor while that curve keeps moving for the next. The resource lock the witness does carry is the closest good fit in the whole vocabulary and is still not exact, because what is reserved is a worst-case loss across an entire outcome vector and it doubles as the pricing seed. The accounting is residue too: a provider's claim is a leaf with an entry position rather than a proportional share of a pot, and the vocabulary has no accounting form for it. What remains is the commercial arrangement --- early buy-back of a live claim, correlated exposure across combined legs, a frontend as a first-class revenue claimant, and an order of loss-bearing declared in a fee schedule rather than encoded as a structure.

The obligations no element discharges are these.
\begin{itemize}\setlength{\itemsep}{0pt}
  \item a bettor's position is a transferable NFT carrying a payout struck against the pool; no operation splits pool collateral into complementary outcome claims and none recombines them
  \item a single protocol-wide liquidity pool is the counterparty to every bet, writing a non-linear event-contingent payoff, so no bettor ever faces another bettor
  \item the price of an outcome is the normalised inverse of the pool's virtual fund for that outcome, loaded with a spread allocated inversely to the outcome's probability under a total-margin constraint, and every bet moves the virtual funds so the odds worsen against incoming flow
  \item the odds struck at the moment of the bet are locked for that bettor and fix the payout, while the curve continues to move for everyone after
  \item each liquidity deposit is a separate leaf of a segment tree over a 2K+1 array, a Condition's profit and loss is distributed across a leaf range, and the updates are lazy — annotated at parent nodes and pushed down only when a provider withdraws
  \item a liquidity provider is locked for seven days from deposit and may withdraw at will thereafter
  \item one permissioned party creates and cancels events, seeds and reprices the sell-side odds, sets the reinforcement and margin, and then declares which outcome won
  \item a Data Provider must post collateral up to the total Reinforcement it expects to use across all its active markets
  \item the pool quotes a price to buy back an open bet before the event concludes, extinguishing a live contingent claim early
  \item a combination bet multiplies the odds of several legs into one payout while the pool's exposure to the combination is not the product of its exposures to the legs
  \item an app's share is the LESSER of 70\% of the revenue it generated and a monthly cap of 50\% of the spread times the bet size summed over its bets
  \item when the pool loses, the liquidity providers', Data Provider's and DAO's rewards go negative, while an app's share can only ever be positive or zero
\end{itemize}
\end{remark}

\begin{remark}
The corpus record carries $Rl$ with no $Au$, which makes this the only protocol in the corpus rejected on a \emph{warrant} --- an element present that nothing consumes --- rather than on a requirement alone. The witness repairs that with $Au$ and $Gs$ on evidence: the live protocol routes every bet through a relayer that executes the bettor's signed order and draws sponsored fees from a paymaster, both documented in the architecture and both absent from the public contract repository, which covers a superseded version. $Ex$ and $Sv$ are dropped because there is no feed, no aggregation and no medianizer, and because servicing discretion is written for a loan rather than for a party that sets the price and then decides who won; $Sh$ is dropped because a provider's claim is a leaf of a segment tree whose value depends on which events resolved during its residence; $Em$ is dropped because what is documented is a share of the pool's profit and loss, which is $Fd$. Dropping the truth elements leaves this venue with no expressible resolution mechanism at all, and that is the finding rather than a defect in the witness.
\end{remark}

\subsection{Grove Finance}\label{sub:cat:prd:grove}

Grove is the other member the sampling frame could not file elsewhere, and it is an allocator and a savings product under one name --- a distinction that is load-bearing. The allocator mints USDS against a Sky allocation vault under a mandate Sky governance approves, and deploys it into credit, most visibly a tokenised collateralised-loan-obligation fund alongside on-chain lending, liquidity and yield venues across several chains; the capital is Sky's, the discretion is Grove's, and the depositor in the economically meaningful sense is a DAO. A retail user who supplies a stablecoin to the savings product receives sUSDS, whose exchange rate accrues at the Sky Savings Rate, plus non-transferable points carrying no claim --- that user is not buying Grove's credit book, and treating this as a deposit-taking vault is an error. Where the money goes is decided by Grove and approved through Sky's weekly governance cycle after a public forum round, and the standing on-chain controls are a rate limit on every controller operation and a freezer role able to revoke the relayer and halt the machine. Recourse therefore runs to Sky governance and to that freezer rather than to a depositor: the savings user's only remedy is redemption, and Grove's own governance token has Sky's pause proxy as its sole authorised admin.

\begin{measurement}\label{meas:cat:prd:grove}
Take $X = \{At,Aw,Fd,Gp,Ix,Ps,Sh,Tg,Xf\}$, with canonical form
$\mathrm{ex}(X) = \{At,Aw,Fd,Gp,Ix,Ps,Sh,Tg,Xf\}$, every element being primitive in it.
It is admissible: no requirement term is open, every element is warranted, and it arms no prohibition. It discharges 8 of the 20 recorded obligations; 12 are residue. No composite of at most three corpus protocols equals it; 1 corpus protocol is contained in it, and together they supply every element except \{Fd,Ix,Ps\}.
It is not minimal: \{Sh,Ix\} can be removed with every obligation still covered and admissibility intact.
\end{measurement}

\begin{remark}
The residue is the mandate, the custody and the flow control. Nothing names an authorisation granted by one governance system over capital it owns and discretion it does not exercise, nor a governance system whose root of authority is another one, nor a credit line whose underwriting is a ballot. Nothing names custody at all: every asset sits in a stateless proxy that holds no logic and forwards calls, which is the inverse of a mutable implementation behind a fixed address and carries the opposite risk. The rate limit --- a linearly refilling per-key bound on how fast a privileged actor may move value --- is the most load-bearing on-chain control in the allocator, and the vocabulary has no flow limiter at all; the gap is general rather than local to this protocol. The largest object here post-dates the corpus entirely: a standing facility that sells immediate stablecoin liquidity to holders of slow-settling tokenised funds, warehousing the timing risk for a fee while the fund's own settlement cycle continues unchanged behind it, sold to counterparties who are not the protocol's depositors at all.

The obligations no element discharges are these.
\begin{itemize}\setlength{\itemsep}{0pt}
  \item a strategy may be funded only after Sky governance approves it through the Atlas Edit weekly cycle following a public forum round in which Grove answers voter questions; only approved strategies move to on-chain deployment
  \item the capital enters by minting USDS against a Sky allocation vault rather than by any user transferring assets in — the deposit is a governance-authorised credit line
  \item every asset the allocator controls sits in one minimal stateless proxy that holds no logic, and every movement of those assets is a call the proxy makes on a controller's instruction; funds never leave it
  \item only holders of a relayer role may invoke a fund-moving function, and the role is exercised by an off-chain planner that submits the transactions
  \item every controller operation is bounded by a rate limit that refills linearly at a configured slope up to a maximum, keyed per asset, per destination or per domain
  \item positions are taken in vaults whose deposits and redemptions do not settle synchronously: a request is made on-chain now and the claim is fulfilled later out of an off-chain process
  \item the economically decisive transfer is not between chains but from a chain into a fund administrator's register of record, where the subscription becomes an entry in a Cayman or Luxembourg vehicle's books
  \item the credit the capital ultimately buys is a tranche of a collateralised loan obligation, which from the allocator's position is a single opaque net asset value with no visible waterfall
  \item users also accrue Grove Points in proportion to their activity, which are non-transferable, have no monetary value and confer no claim or entitlement
  \item GROVE holders stake into stGROVE and vote directly or by delegate, but the token contract's only authorised admin is Sky's pause proxy, so administrative authority over the governance token belongs to another protocol's governance
  \item the Basin facility commits up to \$1 billion a day of immediate stablecoin liquidity to eligible holders of tokenised funds, without purchasing or taking ownership of the underlying assets and without changing the fund's own redemption procedures, settlement cycles or transfer restrictions
  \item in that timelock the external issuer holds the proposal right, Grove governance executes after the delay, and a security multisig may cancel
\end{itemize}
\end{remark}

\begin{remark}
The corpus records $Sh$ without siting it. The witness sites it on the savings share and pairs it with $Ix$, because that share's exchange rate accrues by an index whose rate is a governance parameter of another protocol --- and the siting is what matters most here, since the share is a claim on Sky's savings rate rather than on Grove's book. $Ps$ and $Fd$ are added on evidence, for the peg-swap module through which the allocator converts between its stablecoins at par and for the standing liquidity facility's on-chain fee accrual to authorised issuers, while nothing is claimed for what Grove earns on the Sky mandate, which is published at no primary source and is not guessed at. $Gp$ is retained under protest and $Au$ is declined: the freezer does not pause a system but revokes an actor, leaving the machine running with nobody able to act and restoration requiring a role to be re-granted rather than a flag cleared, and $Au$ describes a principal delegating their own authority under a session policy, which no user here has done.
\end{remark}

\subsection{Kalshi}\label{sub:cat:prd:kalshi}

A user opens an account at a regulated exchange and buys a contract on a stated future fact, at a price that reads directly as a probability. At expiration the contract pays a fixed settlement value to the side whose payout criterion is met; a scalar contract apportions that value between long and short rather than paying it whole, and either way the payment is a clearing-house bookkeeping entry against a segregated settlement account. The outcome is determined by the exchange itself. An internal committee interprets the contract's terms with full discretion and its determination is final and not subject to review, so a disagreement about an outcome is not adjudicated but \emph{decided}; where the fact is indeterminable the contract may settle at its own last traded price, and once trading has begun the venue may designate a different settlement source for a contract that is already open.

\begin{measurement}\label{meas:cat:prd:kalshi}
Take $X = \{Ad,Aw,Bs,Ct,Fd,Gp,Pf,Rd,Sl\}$, with canonical form
$\mathrm{ex}(X) = \{Ad,Aw,Bs,Fd,Gp,Pf,Rd,Sl\}$; the elements \{Ct\} are derived rather than chosen.
It is not admissible: it leaves the requirement terms $L1\!:\!Ex{\mid}Tp{\mid}At$, $L4\!:\!Ex$, $L4\!:\!Li$ open. It discharges 9 of the 22 recorded obligations; 13 are residue. No composite of at most three corpus protocols equals it; 1 corpus protocol is contained in it, and together they supply every element except \{Ad,Bs,Ex,Fd,Gp,Li,Pf,Sl\}.
\end{measurement}

\begin{remark}
What the construction cannot say is the institution. Novation extinguishes a bilateral obligation and replaces it with obligations against a clearing house, and nothing in the vocabulary is a counterparty-substitution operator: $Rd$ names a right to be paid and carries no obligor, and the obligor is the entire content of novation. Beyond it, the operator's own capped first-loss contributions and the callable assessments the clearing house may levy on surviving members --- unfunded until called, capped, and refundable out of later recoveries --- have no symbol either, so the vocabulary reconstructs the mutualised middle of this waterfall and neither end of it. The rest of the residue is the product: a payout criterion written in prose, a settlement source the venue may replace mid-contract, a listing that exists because a regulator did not object, and customer money whose investment income belongs to its custodian. This is the most precisely specified default waterfall in the corpus, and that precision is what makes its inexpressibility a result rather than an omission.

The obligations no element discharges are these.
\begin{itemize}\setlength{\itemsep}{0pt}
  \item orders rest in a continuous two-sided limit order book operated by the exchange, quoted in sub-penny dollar prices and published as bids only
  \item the long and short of one contract are complementary claims on one escrowed dollar, and that complementarity exists only as a price identity — a YES bid at X is a NO ask at \$1.00 minus X — with no object that can be split into or merged out of the pair
  \item on acceptance of a matched trade the clearing house is substituted as seller to the buying participant and buyer to the selling participant, the original bilateral contract is extinguished, and the parties are released from their obligations to each other
  \item open margined positions are marked and settled in cash at least once each business day, with unscheduled intra-day calls available
  \item a scalar contract apportions the \$1.00 between long and short as a function of where the Expiration Value falls, rather than paying it wholly to one side
  \item the operator commits two capped tranches of its own capital to the default waterfall, one ahead of the mutualised fund and one behind it, each capped at the greater of \$20 million or 10\% of the guaranty fund
  \item when the fund is exhausted the clearing house may call assessments on every surviving member, pro rata to their deposit requirement over the preceding three months, capped at 200\% for one default and 550\% for multiple defaults in six months, and refundable pro rata out of later recoveries
  \item the mutualised default fund is partitioned by Contract Segment between the fully-collateralised and margined product classes, and the collateral of non-defaulting customers in each account class may not be applied to a defaulting member's obligations
  \item the outcome of a contract is determined by an internal committee of the exchange with full discretion, within 24 hours, and its determination is final and not subject to review
  \item where the Expiration Value is indeterminable the contract may be settled at its own last traded price
  \item after the first day of trading the venue may designate a new Source Agency and Underlying for a listed contract and change its specifications
  \item a new market exists because the venue self-certified its terms with the regulator under CEA 5c(c) and the regulator did not object within the window
  \item customer funds are held in United States dollars only, segregated by account class under CFTC Reg. 1.20 or Part 22 with the clearing house holding a continuing first-priority security interest and UCC control over them, investable only in US Treasuries, agency securities, reverse repo on those and government money-market funds up to two-year maturity — and the clearing house retains all profit from that investment not paid to members
\end{itemize}
\end{remark}

\begin{remark}
The corpus records $Aw$, $Ct$ and $Rd$. The witness adds $Bs$, $Sl$ and $Ad$ because the clearing rulebook specifies a default waterfall in more detail than any other document in the corpus --- a mutualised guaranty fund, variation-margin gains haircutting, and tear-ups --- and adds $Pf$, $Fd$ and $Gp$ because the venue now lists perpetual futures with a funding rate, publishes a fee schedule computed from the contract's own price, and holds emergency powers under both the exchange and the clearing rulebook. Among the alternatives the loss-absorption term leaves open, the rulebook names its own layers in its own order, so the witness follows the rulebook rather than the nearest DeFi analogue. The construction is rejected, and the rejection is correct: every element that could close the truth term requires a feed, a price series or an independent attester, and the truth source here is a committee.
\end{remark}

\subsection{Polymarket}\label{sub:cat:prd:polymarket}

A user deposits a reserve-backed dollar token and buys outcome shares in a market whose resolution criteria are written in prose beneath the order book. A share pays a unit of collateral if its outcome occurs and nothing otherwise, so its price reads as a probability; the user may also mint the complementary pair directly out of collateral and merge it back, which is what makes the venue fully collateralised with no counterparty. The outcome is proposed to an optimistic oracle by a whitelisted proposer under a bond, and stands if nobody disputes it within the challenge window. A dispute resets the question and re-requests a price, a further dispute escalates to a token-holder vote, and orthogonally to both an admin may flag the market and, after a fixed safety period, write the payout vector directly. The same brand also runs a regulated contract market on which a company officer determines the outcome at sole discretion and may reverse a determination for obvious error, so the same question resolves by bonded assertion on one venue and by officer discretion on the other.

\begin{measurement}\label{meas:cat:prd:polymarket}
Take $X = \{Au,Aw,Ct,Fd,Gp,Gs,Oa,Ob,Ps,Rd,Up\}$, with canonical form
$\mathrm{ex}(X) = \{Aw,Ct,Fd,Gp,Gs,Oa,Ob,Ps,Rd,Up\}$; the elements \{Au\} are derived rather than chosen.
It is admissible: no requirement term is open, every element is warranted, and it arms no prohibition. It discharges 12 of the 20 recorded obligations; 8 are residue. No composite of at most three corpus protocols equals it; 2 corpus protocols are contained in it, and together they supply every element except \{Fd,Ps\}.
\end{measurement}

\begin{remark}
The residue begins with the primitive that defines the sector. Splitting collateral along the \emph{state} space of a condition, under an enforced disjoint and non-degenerate partition and recursively against further conditions, has no element; the only splitting element in the vocabulary partitions a claim along \emph{time} and takes no partition argument at all. That absence propagates into the matching engine, because a crossing pair of buys is settled by minting the complementary pair out of collateral and a crossing pair of sells by merging one back, so the order book and the state partition are the same machine and the vocabulary can name neither. The rest is the market itself: the prose question that every dispute is a dispute about, the completion operator over a family of mutually exclusive markets and the private buffer that finances it, and --- on the regulated leg --- a central counterparty that novates every trade and holds no guaranty fund and no financial resource package at all, against which the loss-absorption symbols reconstruct nothing. A fully collateralised clearing house is pure novation and custody, and that is the case the vocabulary has least to say about, not most.

The obligations no element discharges are these.
\begin{itemize}\setlength{\itemsep}{0pt}
  \item one unit of collateral is partitioned into one YES and one NO outcome token, and returning both recombines them into that unit, under an enforced invariant that the partition is pairwise disjoint, that no index set is empty and that none is the whole state space
  \item two of the three settlement modes of the matching engine are that primitive: two crossing buy orders mint a complementary pair out of collateral, and two crossing sell orders merge a pair back into collateral, with only a crossing buy against a sell settling by transfer
  \item the question being resolved is prose: a natural-language resolution criterion posted under the order book, against which every dispute is a dispute about text
  \item an admin may flag a market, which pauses it and starts a fixed one-hour safety period, after which the admin writes the payout vector directly, overriding the oracle entirely
  \item for a family of binary markets of which exactly one resolves true, a holder may exchange a bundle of NO positions for collateral plus the residual YES, financed before resolution out of a wrapped-collateral buffer the adapter holds itself; if the family's mutual-exclusivity constraint is violated the family cannot fully resolve and funds are frozen
  \item the exchange wraps that collateral a second time as PMCT through a collateral adapter with dedicated onramp and offramp contracts, so an adapter can hold liquidity outside the conditional-token accounting
  \item on the US venue every matched trade is novated: the clearing house is immediately substituted as seller to the purchaser and purchaser to the seller
  \item on the US venue a company officer runs a Contract Outcome Review Process, determines the final outcome at sole discretion, and may reverse a determined outcome in the case of obvious error, with company liability to any one participant capped at \$2,500 a day, \$5,000 a month and \$50,000 a year
\end{itemize}
\end{remark}

\begin{remark}
The corpus canonical form omits $Ps$, and the collateral has since changed: it is now a token backed at par by a reserve, convertible in either direction with no fee and enforced on-chain, which is what the peg-swap element names exactly rather than approximately. The witness adds $Ct$ and $Fd$ on the same evidential footing --- the pre-trade collateral check the affiliated clearing house performs before it will novate, documented in the rulebook, and the per-order fee the exchange charges at match time, documented in the exchange repository. $Oa$ is kept for resolution although its shape now fits badly, because the ladder it names --- assert, bond, dispute, escalate --- is genuinely present; what it cannot express is that proposals are permissioned per request while disputes remain open to anyone, and that asymmetry is the structurally interesting half. Among the alternatives the truth term leaves open, this is the only honest one: there is no feed and no price series, and an attestation would misname a bonded assertion about prose.
\end{remark}

\subsection{Steakhouse Financial}\label{sub:cat:prd:steakhouse}

This is not a prediction market. It sits in this category because it fits no conventional one, and the category is where the sampling frame put what it could not file elsewhere. A depositor puts a stablecoin into a vault and receives shares; the vault does not lend but allocates the deposit across isolated lending markets on a base protocol, so what the depositor buys is a manager and what they hold is the credit and oracle risk of markets they did not choose. A named firm decides which markets the capital may enter, at what ceiling for each, and in what order deposits fill and withdrawals drain, and takes a share of the interest for doing so. The depositor's recourse is \emph{preventive} and not \emph{compensatory}: they may exit while liquidity lasts, and they hold veto weight over a critical curator action in proportion to their deposit, exercisable inside a timelock explicitly sized to let them withdraw or veto before capital moves --- but after a loss there is no addressee, the terms disclaiming any fiduciary relationship, denying that the firm is a broker, intermediary, agent, custodian or adviser, and capping total liability at a nominal sum.

\begin{measurement}\label{meas:cat:prd:steakhouse}
Take $X = \{Ct,Fd,Gp,Im,Rd,Sh,Tg\}$, with canonical form
$\mathrm{ex}(X) = \{Fd,Gp,Im,Rd,Sh,Tg\}$; the elements \{Ct\} are derived rather than chosen.
It is not admissible: it leaves the requirement terms $L1\!:\!Ex{\mid}Tp{\mid}At$, $L1\!:\!Li{\mid}Ad{\mid}Sl{\mid}Bs$ open. It discharges 7 of the 17 recorded obligations; 10 are residue. No composite of at most three corpus protocols equals it; no corpus protocol is contained in it.
\end{measurement}

\begin{remark}
The unnamed obligations are the mandate and everything that gives it shape. Nothing names the delegation of investment discretion over other people's capital, and nothing names the supply cap that is the vault's own risk instrument --- a portfolio limit on a manager rather than a solvency test on a borrower, defaulting to nothing so that a market is unusable until an affirmative act makes it usable. The depositor-weighted veto has no element, and neither has its absence: the vocabulary cannot distinguish a mandate with recourse from one without, so a mandate with recourse and one without decompose identically and the single most important fact about the instrument is invisible. The fee is residue in its mechanism if not in its fact, being paid by minting fresh shares to the recipient --- dilution, which cannot fail for want of assets and does not appear in the asset accounting at all. Above all, where the capital goes is a function of a published, versioned rating document, which is the prediction venues' rulebook problem in a different sector and is the reason both halves of this category belong in one place.

The obligations no element discharges are these.
\begin{itemize}\setlength{\itemsep}{0pt}
  \item a named firm decides which markets the capital may enter, in what order they fill and drain, and rebalances at will, for a share of the interest generated
  \item each market carries a supply cap defaulting to zero, so no market may receive any capital until the curator explicitly caps it up, and a vault may enable at most thirty markets
  \item deposits fill along an ordered supply queue up to each cap and withdrawals drain along an ordered withdraw queue, so the identity of the market receiving the next dollar is a curator-set ordering
  \item capital not deployed into a market sits as idle supply in a market with null oracle and null collateral addresses, immediately withdrawable
  \item vault users become vault guardians in proportion to the size of their deposit, and may veto a critical curator action inside the timelock window
  \item authority is split four ways — an appointer, a strategy-setter who sets caps and fees, bounded executors who may only move funds within the strategy-setter's limits, and sentinels who may only reduce
  \item the fee is paid by minting new vault shares to the fee recipients, so it is borne as dilution of every other holder rather than as a transfer of assets
  \item the curator disclaims being a broker, intermediary, agent, custodian or advisor, disclaims any fiduciary relationship or obligation, caps total liability at \$1,000, and routes disputes to Cayman arbitration with a class-action waiver
  \item before a market may be listed at all, at least two members of the risk team run a quality-control review across six named risk vectors, and every vault passes a due-diligence questionnaire with at least two internal and one external reviewer
  \item where the capital goes is a function of a published, versioned three-layer rating — an asset rating over issuer, credit and operational pillars, a platform rating, and a market rating over oracle, liquidity, price fluctuation, LLTV and credit-enhancement criteria — which is entirely off-chain
\end{itemize}
\end{remark}

\begin{remark}
The corpus records $Wq$; the witness drops it, because the queues here are queues of \emph{venues} in priority order and no claimant ever waits in them, and adds $Rd$ instead, because the redemption right is unconditional in the contract even though its exercise is not. The rest is retained --- the share, the isolated markets the vault is a client of, the inherited threshold test, the asymmetric timelock that delays a raised cap and lets a lowered one take effect at once, the guardian and the fee. Among the alternatives the mandate leaves open, no element is chosen for the mandate itself: forcing servicing discretion onto exposure selection would name a mechanism this firm does not have, and declining is the honest move. The construction is rejected for want of a price source and of a loss-absorption mechanism, and that is right --- this vault computes no price and absorbs no loss, both being properties of the markets underneath it, and the carrier has no way to record an obligation discharged one level down.
\end{remark}

\section{Evidence}\label{sec:evidence}

Each obligation recorded in this supplement carries a source. The sources are
listed here by application, in the order the obligations appear, with the
access date on which each was consulted. A source consulted for several
obligations of one application is listed once.

\subsection*{Curve}
\begin{itemize}\setlength{\itemsep}{0pt}\small
\item \texttt{https://docs.curve.finance/pdf/whitepapers/whitepaper\_stableswap.pdf} (2026-08-04)
\item \texttt{https://docs.curve.finance/pdf/whitepapers/whitepaper\_cryptoswap.pdf} (2026-08-04)
\item \texttt{https://docs.curve.finance/developer/amm/cryptoswap-in-depth} (2026-08-04)
\item \texttt{https://eth.blockscout.com/api/v2/smart-contracts/0xec977F46467a3021785Cff88894886E617abd65b} (2026-08-04)
\item \texttt{https://docs.curve.finance/developer/amm/legacy/stableswap-overview} (2026-08-04)
\item \texttt{https://docs.curve.finance/developer/amm/factory/tricrypto-ng/deployer-api} (2026-08-04)
\item \texttt{https://docs.curve.finance/developer/amm/factory/stableswap-ng/deployer-api} (2026-08-04)
\item \texttt{https://docs.curve.finance/developer/gauges/gauges/liquidity-gauge-v6} (2026-08-04)
\item \texttt{https://docs.curve.finance/developer/curve-dao/governance/overview} (2026-08-04)
\item \texttt{https://docs.curve.finance/developer/curve-dao/voting-escrow/} (2026-08-04)
\item \texttt{https://docs.curve.finance/developer/fees/overview} (2026-08-04)
\item \texttt{https://docs.curve.finance/developer/fees/cow-swap-burner} (2026-08-04)
\item \texttt{https://docs.curve.finance/governance} (2026-08-04)
\end{itemize}

\subsection*{Fluid}
\begin{itemize}\setlength{\itemsep}{0pt}\small
\item \texttt{https://fluid.guides.instadapp.io/dex-protocol/protocol-details.md} (2026-08-04)
\item \texttt{https://docs.fluid.instadapp.io/integrate/dex-v2-swaps.html} (2026-08-04)
\item \texttt{https://fluid.guides.instadapp.io/dex-protocol/smart-collateral.md} (2026-08-04)
\item \texttt{https://fluid.guides.instadapp.io/dex-protocol/smart-debt.md} (2026-08-04)
\item \texttt{https://fluid.guides.instadapp.io/liquidity-layer/protocols-on-fluid.md} (2026-08-04)
\item \texttt{https://fluid.guides.instadapp.io/liquidity-layer/unifying-liquidity.md} (2026-08-04)
\item \texttt{https://fluid.guides.instadapp.io/vault-protocol/advanced-liquidation-mechanism/vault-liquidation.md} (2026-08-04)
\item \texttt{https://fluid.guides.instadapp.io/vault-protocol/advanced-liquidation-mechanism.md} (2026-08-04)
\item \texttt{https://docs.fluid.instadapp.io/} (2026-08-04)
\item \texttt{https://fluid.guides.instadapp.io/dex-protocol/oracles.md} (2026-08-04)
\item \texttt{https://fluid.io/blog/protocol-introducing-fluid-dex-v2} (2026-08-04)
\item \texttt{https://github.com/Instadapp/fluid-governance} (2026-08-04)
\item \texttt{https://fluid.guides.instadapp.io/dex-protocol/fees.md} (2026-08-04)
\end{itemize}

\subsection*{PancakeSwap}
\begin{itemize}\setlength{\itemsep}{0pt}\small
\item \texttt{https://docs.pancakeswap.finance/earn/pancakeswap-pools.md} (2026-08-04)
\item \texttt{https://docs.pancakeswap.finance/trade/pancakeswap-infinity/pool-types/infinity-clamm-and-lbamm.md} (2026-08-04)
\item \texttt{https://developer.pancakeswap.finance/contracts/v3/addresses} (2026-08-04)
\item \texttt{https://docs.pancakeswap.finance/trade/pancakeswap-infinity/key-features.md} (2026-08-04)
\item \texttt{https://developer.pancakeswap.finance/contracts/infinity/overview/accounting-layer-vault} (2026-08-04)
\item \texttt{https://developer.pancakeswap.finance/contracts/infinity/overview} (2026-08-04)
\item \texttt{https://developer.pancakeswap.finance/contracts/infinity/faq/pancakeswap-infinity-vs-uniswap-v4} (2026-08-04)
\item \texttt{https://raw.githubusercontent.com/pancakeswap/infinity-core/main/src/ProtocolFeeController.sol} (2026-08-04)
\item \texttt{https://docs.pancakeswap.finance/protocol/cake-tokenomics.md} (2026-08-04)
\item \texttt{https://raw.githubusercontent.com/pancakeswap/infinity-core/main/src/pool-cl/CLPoolManagerOwner.sol} (2026-08-04)
\item \texttt{https://docs.pancakeswap.finance/welcome-to-pancakeswap/vecake-sunset} (2026-08-04)
\end{itemize}

\subsection*{Raydium}
\begin{itemize}\setlength{\itemsep}{0pt}\small
\item \texttt{https://docs.raydium.io/security/oracle-and-token-risks} (2026-08-04)
\item \texttt{https://docs.raydium.io/products/clmm/ticks-and-positions} (2026-08-04)
\item \texttt{https://docs.raydium.io/reference/program-addresses} (2026-08-04)
\item \texttt{https://docs.raydium.io/products/cpmm/fees} (2026-08-04)
\item \texttt{https://docs.raydium.io/products/clmm/fees} (2026-08-04)
\item \texttt{https://api-v3.raydium.io/main/cpmm-config} (2026-08-04)
\item \texttt{https://docs.raydium.io/ray/protocol-fees} (2026-08-04)
\item \texttt{https://docs.raydium.io/ray/treasury} (2026-08-04)
\item \texttt{https://docs.raydium.io/products/clmm/accounts} (2026-08-04)
\item \texttt{https://docs.raydium.io/reference/changelog/2026-07-22-amm-v4-openbook-removal} (2026-07-22)
\item \texttt{https://docs.raydium.io/products/launchlab/bonding-curve} (2026-08-04)
\item \texttt{https://docs.raydium.io/products/launchlab/overview} (2026-08-04)
\item \texttt{https://docs.raydium.io/user-flows/burn-and-earn} (2026-08-04)
\item \texttt{https://docs.raydium.io/security/admin-and-multisig} (2026-08-04)
\item \texttt{https://verify.osec.io/status/CPMMoo8L3F4NbTegBCKVNunggL7H1ZpdTHKxQB5qKP1C} (2026-08-04)
\end{itemize}

\subsection*{Uniswap}
\begin{itemize}\setlength{\itemsep}{0pt}\small
\item \texttt{https://raw.githubusercontent.com/Uniswap/v2-core/master/contracts/UniswapV2Factory.sol} (2026-08-04)
\item \texttt{https://raw.githubusercontent.com/Uniswap/v3-core/main/contracts/UniswapV3Factory.sol} (2026-08-04)
\item \texttt{https://developers.uniswap.org/llms.mdx/docs/protocols/v4/concepts/flash-accounting} (2026-08-04)
\item \texttt{https://developers.uniswap.org/llms.mdx/docs/protocols/protocol-fee/overview} (2026-08-04)
\item \texttt{https://developers.uniswap.org/llms.mdx/docs/protocols/v2/concepts/oracles} (2026-08-04)
\item \texttt{https://developers.uniswap.org/llms.mdx/docs/protocols/v4/concepts/hooks} (2026-08-04)
\item \texttt{https://github.com/Uniswap/v4-core/blob/v4.0.0/README.md} (2026-08-04)
\item \texttt{https://raw.githubusercontent.com/Uniswap/v4-core/v4.0.0/src/libraries/Hooks.sol} (2026-08-04)
\item \texttt{https://developers.uniswap.org/llms.mdx/docs/ecosystem/governance/technical-reference} (2026-08-04)
\item \texttt{https://raw.githubusercontent.com/Uniswap/v4-core/v4.0.0/src/PoolManager.sol} (2026-08-04)
\item \texttt{https://github.com/Uniswap/protocol-fees} (2026-08-04)
\item \texttt{https://blog.uniswap.org/unification} (2026-08-04)
\item \texttt{https://vote.uniswapfoundation.org/proposals/93} (2026-08-04)
\end{itemize}

\subsection*{Aave V3}
\begin{itemize}\setlength{\itemsep}{0pt}\small
\item \texttt{https://aave.com/docs/developers/smart-contracts/pool} (2026-08-04)
\item \texttt{https://raw.githubusercontent.com/aave-dao/aave-v3-origin/main/src/contracts/misc/DefaultReserveInterestRateStrategyV2.sol} (2026-08-04)
\item \texttt{https://raw.githubusercontent.com/aave-dao/aave-v3-risk-stewards/main/src/contracts/examples/EthereumExample.sol} (2026-08-04)
\item \texttt{https://raw.githubusercontent.com/aave-dao/aave-v3-origin/main/src/contracts/protocol/libraries/logic/LiquidationLogic.sol} (2026-08-04)
\item \texttt{https://raw.githubusercontent.com/aave-dao/aave-v3-risk-stewards/main/src/contracts/RiskSteward.sol} (2026-08-04)
\item \texttt{https://raw.githubusercontent.com/aave-dao/aave-v3-origin/main/src/contracts/protocol/libraries/configuration/ReserveConfiguration.sol} (2026-08-04)
\item \texttt{https://raw.githubusercontent.com/aave-dao/aave-v3-origin/main/src/contracts/protocol/configuration/ACLManager.sol} (2026-08-04)
\item \texttt{https://raw.githubusercontent.com/aave-dao/aave-v3-risk-stewards/main/src/interfaces/IRiskSteward.sol} (2026-08-04)
\item \texttt{https://raw.githubusercontent.com/aave-dao/aave-v3-origin/main/src/contracts/protocol/configuration/PoolAddressesProvider.sol} (2026-08-04)
\item \texttt{https://raw.githubusercontent.com/aave-dao/aave-governance-v3/main/src/contracts/payloads/PayloadsControllerCore.sol} (2026-08-04)
\end{itemize}

\subsection*{JustLend V1}
\begin{itemize}\setlength{\itemsep}{0pt}\small
\item \texttt{https://api.github.com/repos/justlend/justlend-protocol/git/trees/main?recursive=1} (2026-08-04)
\item \texttt{https://raw.githubusercontent.com/justlend/justlend-protocol/main/contracts/BaseJumpRateModelV2.sol} (2026-08-04)
\item \texttt{https://raw.githubusercontent.com/justlend/justlend-protocol/main/contracts/Comptroller.sol} (2026-08-04)
\end{itemize}

\subsection*{Maple}
\begin{itemize}\setlength{\itemsep}{0pt}\small
\item \texttt{https://api.github.com/repos/maple-labs/pool-v2/git/trees/main?recursive=1} (2026-08-04)
\item \texttt{https://raw.githubusercontent.com/maple-labs/pool-v2/main/contracts/MaplePoolManager.sol} (2026-08-04)
\item \texttt{https://raw.githubusercontent.com/maple-labs/fixed-term-loan/main/contracts/MapleLoanStorage.sol} (2026-08-04)
\end{itemize}

\subsection*{Morpho}
\begin{itemize}\setlength{\itemsep}{0pt}\small
\item \texttt{https://raw.githubusercontent.com/morpho-org/morpho-blue/main/src/Morpho.sol} (2026-08-04)
\item \texttt{https://raw.githubusercontent.com/morpho-org/morpho-blue/main/README.md} (2026-08-04)
\item \texttt{https://raw.githubusercontent.com/morpho-org/morpho-blue/main/src/libraries/ConstantsLib.sol} (2026-08-04)
\item \texttt{https://raw.githubusercontent.com/morpho-org/morpho-blue-irm/main/src/adaptive-curve-irm/AdaptiveCurveIrm.sol} (2026-08-04)
\item \texttt{https://raw.githubusercontent.com/morpho-org/morpho-blue-irm/main/src/adaptive-curve-irm/libraries/ConstantsLib.sol} (2026-08-04)
\item \texttt{https://raw.githubusercontent.com/morpho-org/metamorpho/main/src/MetaMorpho.sol} (2026-08-04)
\end{itemize}

\subsection*{SparkLend}
\begin{itemize}\setlength{\itemsep}{0pt}\small
\item \texttt{https://api.github.com/repos/sparkdotfi/sparklend-v1-core/git/trees/dev?recursive=1} (2026-08-04)
\item \texttt{https://raw.githubusercontent.com/sparkdotfi/sparklend-advanced/master/src/RateTargetKinkInterestRateStrategy.sol} (2026-08-04)
\item \texttt{https://raw.githubusercontent.com/sparkdotfi/sparklend-advanced/master/src/SSRRateSource.sol} (2026-08-04)
\item \texttt{https://raw.githubusercontent.com/sparkdotfi/sparklend-advanced/master/src/RateTargetBaseInterestRateStrategy.sol} (2026-08-04)
\item \texttt{https://api.github.com/repos/sparkdotfi/spark-alm-controller/git/trees/master?recursive=1} (2026-08-04)
\item \texttt{https://raw.githubusercontent.com/sparkdotfi/sparklend-cap-automator/master/src/CapAutomator.sol} (2026-08-04)
\item \texttt{https://raw.githubusercontent.com/sparkdotfi/spark-alm-controller/master/src/RateLimits.sol} (2026-08-04)
\item \texttt{https://github.com/sparkdotfi/sparklend-advanced} (2026-08-04)
\item \texttt{https://api.github.com/orgs/sparkdotfi/repos?per\_page=100} (2026-08-04)
\end{itemize}

\subsection*{Ethena}
\begin{itemize}\setlength{\itemsep}{0pt}\small
\item \texttt{https://docs.ethena.fi/protocol-overview/peg-arbitrage-mechanism.md} (2026-08-04)
\item \texttt{https://docs.ethena.fi/overview/how-usde-works.md} (2026-08-04)
\item \texttt{https://docs.ethena.fi/technical-design/minting-usde/mint-and-redeem-contract-v2.md} (2026-08-04)
\item \texttt{https://docs.ethena.fi/backing-assets/overview.md} (2026-08-04)
\item \texttt{https://docs.ethena.fi/protocol-overview/reserve-fund.md} (2026-08-04)
\item \texttt{https://docs.ethena.fi/protocol-overview/risks/exchange-failure-risk.md} (2026-08-04)
\item \texttt{https://docs.ethena.fi/backing-custody-and-security/overview/off-exchange-settlement-in-detail.md} (2026-08-04)
\item \texttt{https://docs.ethena.fi/technical-design/use-of-oracles.md} (2026-08-04)
\item \texttt{https://docs.ethena.fi/technical-design/staking-usde.md} (2026-08-04)
\item \texttt{https://docs.ethena.fi/protocol-overview/rewards-mechanism.md} (2026-08-04)
\item \texttt{https://docs.ethena.fi/technical-design/staking-usde/staking-key-functions.md} (2026-08-04)
\item \texttt{https://docs.ethena.fi/technical-design/key-trust-assumptions/matrix-of-multisig-and-timelocks.md} (2026-08-04)
\item \texttt{https://docs.ethena.fi/technical-design/key-addresses.md} (2026-08-04)
\end{itemize}

\subsection*{Liquity}
\begin{itemize}\setlength{\itemsep}{0pt}\small
\item \texttt{https://docs.liquity.org/v2-faq/borrowing-and-liquidations.md} (2026-08-04)
\item \texttt{https://github.com/liquity/bold} (2026-08-04)
\item \texttt{https://docs.liquity.org/liquity-v1/faq/general.md} (2026-08-04)
\item \texttt{https://docs.liquity.org/v2-faq/redemptions-and-delegation} (2026-08-04)
\item \texttt{https://docs.liquity.org/liquity-v1/faq/stability-pool-and-liquidations.md} (2026-08-04)
\item \texttt{https://docs.liquity.org/liquity-v1/faq/recovery-mode.md} (2026-08-04)
\item \texttt{https://docs.liquity.org/v2-faq/lqty-staking.md} (2026-08-04)
\item \texttt{https://docs.liquity.org/v2-faq/general.md} (2026-08-04)
\end{itemize}

\subsection*{Lista}
\begin{itemize}\setlength{\itemsep}{0pt}\small
\item \texttt{https://docs.bsc.lista.org/for-developer/collateral-debt-position/mechanics} (2026-08-04)
\item \texttt{https://docs.bsc.lista.org/for-developer/multi-oracle/multi-oracle-standard.md} (2026-08-04)
\item \texttt{https://docs.bsc.lista.org/for-developer/collateral-debt-position/smart-contract.md} (2026-08-04)
\item \texttt{https://docs.bsc.lista.org/introduction/collateral-debt-position-lisusd/lisusd/algorithmic-market-operations-amo.md} (2026-08-04)
\item \texttt{https://docs.bsc.lista.org/introduction/collateral-debt-position-lisusd/collateral/loan-liquidation} (2026-08-04)
\item \texttt{https://docs.bsc.lista.org/introduction/collateral-debt-position-lisusd/lisusd/stable-pool-price-stability-module-psm} (2026-08-04)
\item \texttt{https://docs.bsc.lista.org/introduction/collateral-debt-position-lisusd/lisusd/d3m-direct-deposit-module.md} (2026-08-04)
\item \texttt{https://docs.bsc.lista.org/introduction/collateral-debt-position-lisusd/lisusd/lisusd-saving-rate-lsr.md} (2026-08-04)
\item \texttt{https://docs.bsc.lista.org/introduction/collateral-debt-position-lisusd/lisusd.md} (2026-08-04)
\item \texttt{https://docs.bsc.lista.org/governance/lista/governance.md} (2026-08-04)
\item \texttt{https://docs.bsc.lista.org/introduction/collateral-debt-position-lisusd/collateral/classic-collateral-options.md} (2026-08-04)
\end{itemize}

\subsection*{Sky}
\begin{itemize}\setlength{\itemsep}{0pt}\small
\item \texttt{https://developers.skyeco.com/protocol/core/vat/} (2026-08-04)
\item \texttt{https://developers.skyeco.com/protocol/core/osm/} (2026-08-04)
\item \texttt{https://developers.skyeco.com/deep-dives/rate-mechanism/} (2026-08-04)
\item \texttt{https://developers.skyeco.com/protocol/rates/jug/} (2026-08-04)
\item \texttt{https://github.com/sky-ecosystem/sp-beam} (2026-08-05)
\item \texttt{https://ethereum-rpc.publicnode.com} (2026-08-04)
\item \texttt{https://github.com/sky-ecosystem/dss-emergency-spells} (2026-08-05)
\item \texttt{https://developers.skyeco.com/protocol/vaults/collateral-liquidation/} (2026-08-04)
\item \texttt{https://developers.skyeco.com/protocol/core/vow/} (2026-08-04)
\item \texttt{https://github.com/sky-ecosystem/dss-flappers/blob/dev/README.md} (2026-08-04)
\item \texttt{https://developers.skyeco.com/protocol/tokens/susds/} (2026-08-04)
\item \texttt{https://developers.skyeco.com/protocol/liquidity/litepsm/} (2026-08-04)
\item \texttt{https://github.com/sky-ecosystem/usds} (2026-08-05)
\item \texttt{https://chainlog.skyeco.com/api/mainnet/active.json} (2026-08-04)
\item \texttt{https://github.com/sky-ecosystem/dss-allocator} (2026-08-05)
\item \texttt{https://developers.skyeco.com/protocol/governance/chief/} (2026-08-04)
\item \texttt{https://developers.skyeco.com/protocol/governance/pause/} (2026-08-04)
\item \texttt{https://developers.skyeco.com/protocol/tokens/stusds/} (2026-08-04)
\end{itemize}

\subsection*{USDD}
\begin{itemize}\setlength{\itemsep}{0pt}\small
\item \texttt{https://docs.usdd.io/developers/core-contracts/vat.md} (2026-08-04)
\item \texttt{https://docs.usdd.io/introduction/core-features.md} (2026-08-04)
\item \texttt{https://docs.usdd.io/developers/oracle.md} (2026-08-04)
\item \texttt{https://docs.usdd.io/developers/core-contracts/jug.md} (2026-08-04)
\item \texttt{https://docs.usdd.io/system-architecture/regarding-the-dynamic-apy-pricing-mechanism.md} (2026-08-04)
\item \texttt{https://docs.usdd.io/system-architecture/susdd-mechanism.md} (2026-08-04)
\item \texttt{https://docs.usdd.io/system-architecture/smart-allocator.md} (2026-08-04)
\item \texttt{https://docs.usdd.io/developers/core-contracts/clip.md} (2026-08-04)
\item \texttt{https://docs.usdd.io/system-architecture/system-architecture.md} (2026-08-04)
\item \texttt{https://docs.usdd.io/user-guide/psm-peg-stability-module.md} (2026-08-04)
\item \texttt{https://github.com/decentralized-usd/psm} (2026-08-05)
\item \texttt{https://docs.usdd.io/developers/deployment-addresses.md} (2026-08-04)
\item \texttt{https://docs.usdd.io/developers/core-contracts.md} (2026-08-04)
\item \texttt{https://docs.usdd.io/governance/overview.md} (2026-08-04)
\end{itemize}

\subsection*{Babylon}
\begin{itemize}\setlength{\itemsep}{0pt}\small
\item \texttt{https://github.com/babylonlabs-io/babylon/blob/main/docs/staking-script.md} (2026-08-04)
\item \texttt{https://github.com/babylonlabs-io/babylon/blob/main/x/btcstaking/README.md} (2026-08-04)
\item \texttt{https://github.com/babylonlabs-io/babylon/blob/main/README.md} (2026-08-04)
\item \texttt{https://babylon-api.polkachu.com/babylon/btcstaking/v1/params} (2026-08-04)
\item \texttt{https://github.com/babylonlabs-io/networks/blob/main/bbn-1/covenant-committee/README.md} (2026-08-04)
\item \texttt{https://babylon-api.polkachu.com/babylon/finality/v1/params} (2026-08-04)
\item \texttt{https://github.com/babylonlabs-io/babylon/blob/main/x/finality/README.md} (2026-08-04)
\item \texttt{https://docs.babylonlabs.io/guides/overview/babylon\_genesis/} (2026-08-04)
\item \texttt{https://docs.babylonlabs.io/guides/overview/bitcoin\_staking/} (2026-08-04)
\item \texttt{https://github.com/babylonlabs-io/babylon/blob/main/x/incentive/README.md} (2026-08-04)
\item \texttt{https://github.com/babylonlabs-io/babylon/blob/main/x/costaking/README.md} (2026-08-04)
\item \texttt{https://api.github.com/repos/babylonlabs-io/babylon/contents/x} (2026-08-04)
\end{itemize}

\subsection*{Binance staked ETH}
\begin{itemize}\setlength{\itemsep}{0pt}\small
\item \texttt{https://eth.blockscout.com/api/v2/smart-contracts/0x9E021c9607bD3ADB7424D3b25a2D35763ff180BB} (2026-08-04)
\item \texttt{https://www.binance.com/en/support/faq/what-is-wbeth-e252366155174ba6887f6b32e3798273} (2026-08-04)
\item \texttt{https://eth.blockscout.com/api/v2/smart-contracts/0x81720695e43A39C52557Ce6386feB3FAAC215f06} (2026-08-04)
\item \texttt{https://eth.blockscout.com/api/v2/smart-contracts/0x542059d658624DF6452b22B10302A15a6AB59f10} (2026-08-04)
\item \texttt{https://eth.blockscout.com/api/v2/smart-contracts/0xa2E3356610840701BDf5611a53974510Ae27E2e1} (2026-08-04)
\item \texttt{https://eth.blockscout.com/api/v2/smart-contracts/0x79973d557CD9dd87eb61e250cc2572c990e20196} (2026-08-04)
\item \texttt{https://www.binance.com/en/proof-of-reserves} (2026-08-04)
\item \texttt{https://llamarisk.com/research/risk-collateral-risk-assessment-binance-wrapped-beacon-eth-wbeth} (2026-08-04)
\item \texttt{https://www.binance.com/en/research/projects/wrapped-beacon-eth-wbeth} (2026-08-04)
\item \texttt{https://www.binance.com/en/support/announcement/binance-introduces-wrapped-beacon-eth-wbeth-on-eth-staking-a1197f34d832445db41654ad01f56b4d} (2026-08-04)
\end{itemize}

\subsection*{EigenLayer}
\begin{itemize}\setlength{\itemsep}{0pt}\small
\item \texttt{https://github.com/Layr-Labs/eigenlayer-contracts/blob/main/docs/core/StrategyManager.md} (2026-08-04)
\item \texttt{https://github.com/Layr-Labs/eigenlayer-contracts/blob/main/docs/core/EigenPodManager.md} (2026-08-04)
\item \texttt{https://github.com/Layr-Labs/eigenlayer-contracts/blob/main/docs/core/DelegationManager.md} (2026-08-04)
\item \texttt{https://github.com/Layr-Labs/eigenlayer-contracts/blob/main/docs/core/AllocationManager.md} (2026-08-04)
\item \texttt{https://github.com/eigenfoundation/ELIPs/blob/main/ELIPs/ELIP-002.md} (2026-08-04)
\item \texttt{https://github.com/eigenfoundation/ELIPs/blob/main/ELIPs/ELIP-006.md} (2026-08-04)
\item \texttt{https://github.com/Layr-Labs/eigenlayer-contracts/blob/main/docs/README.md} (2026-08-04)
\item \texttt{https://api.github.com/repos/Layr-Labs/eigenlayer-contracts/releases/latest} (2026-08-04)
\item \texttt{https://github.com/Layr-Labs/eigenlayer-contracts/blob/main/docs/core/RewardsCoordinator.md} (2026-08-04)
\item \texttt{https://github.com/Layr-Labs/eigenlayer-contracts/blob/main/docs/multichain/README.md} (2026-08-04)
\item \texttt{https://docs.eigencloud.xyz/assets/files/EIGEN\_Token\_Whitepaper-0df8e17b7efa052fd2a22e1ade9c6f69.pdf} (2026-08-04)
\item \texttt{https://api.github.com/repos/Layr-Labs/eigenlayer-contracts/contents/src/contracts} (2026-08-04)
\item \texttt{https://api.github.com/repos/Layr-Labs/eigenlayer-contracts} (2026-08-04)
\end{itemize}

\subsection*{ether.fi}
\begin{itemize}\setlength{\itemsep}{0pt}\small
\item \texttt{https://github.com/etherfi-protocol/smart-contracts/blob/b4a0968087b178bc346cdf6bee6c0597bf4c42c7/src/core/EETH.sol} (2026-08-04)
\item \texttt{https://github.com/etherfi-protocol/smart-contracts/blob/b4a0968087b178bc346cdf6bee6c0597bf4c42c7/src/core/LiquidityPool.sol} (2026-08-04)
\item \texttt{https://github.com/etherfi-protocol/smart-contracts/blob/b4a0968087b178bc346cdf6bee6c0597bf4c42c7/src/core/WeETH.sol} (2026-08-04)
\item \texttt{https://github.com/etherfi-protocol/smart-contracts/blob/b4a0968087b178bc346cdf6bee6c0597bf4c42c7/src/oracle/EtherFiOracle.sol} (2026-08-04)
\item \texttt{https://github.com/etherfi-protocol/smart-contracts/blob/b4a0968087b178bc346cdf6bee6c0597bf4c42c7/src/staking/NodeOperatorManager.sol} (2026-08-04)
\item \texttt{https://github.com/etherfi-protocol/smart-contracts/blob/b4a0968087b178bc346cdf6bee6c0597bf4c42c7/src/staking/AuctionManager.sol} (2026-08-04)
\item \texttt{https://github.com/etherfi-protocol/smart-contracts/blob/b4a0968087b178bc346cdf6bee6c0597bf4c42c7/src/staking/StakingManager.sol} (2026-08-04)
\item \texttt{https://github.com/etherfi-protocol/smart-contracts/blob/b4a0968087b178bc346cdf6bee6c0597bf4c42c7/src/staking/EtherFiNodesManager.sol} (2026-08-04)
\item \texttt{https://etherfi.gitbook.io/etherfi/resources/ether.fi-whitepaper/ether.fi-staking} (2026-08-04)
\item \texttt{https://api.github.com/repos/etherfi-protocol/smart-contracts/contents/src/archive} (2026-08-04)
\item \texttt{https://github.com/etherfi-protocol/smart-contracts/blob/b4a0968087b178bc346cdf6bee6c0597bf4c42c7/README.md} (2026-08-04)
\item \texttt{https://github.com/etherfi-protocol/smart-contracts/blob/b4a0968087b178bc346cdf6bee6c0597bf4c42c7/src/withdrawals/WithdrawRequestNFT.sol} (2026-08-04)
\item \texttt{https://github.com/etherfi-protocol/smart-contracts/blob/b4a0968087b178bc346cdf6bee6c0597bf4c42c7/src/withdrawals/EtherFiRedemptionManager.sol} (2026-08-04)
\item \texttt{https://github.com/etherfi-protocol/smart-contracts/blob/b4a0968087b178bc346cdf6bee6c0597bf4c42c7/src/withdrawals/PriorityWithdrawalQueue.sol} (2026-08-04)
\item \texttt{https://etherfi.gitbook.io/etherfi/resources/ether.fi-whitepaper/ether.fi-re-staking} (2026-08-04)
\item \texttt{https://github.com/etherfi-protocol/smart-contracts/blob/b4a0968087b178bc346cdf6bee6c0597bf4c42c7/src/restaking/EtherFiRestaker.sol} (2026-08-04)
\item \texttt{https://github.com/eigenfoundation/ELIPs/blob/main/ELIPs/ELIP-002.md} (2026-08-04)
\item \texttt{https://etherfi.gitbook.io/etherfi/developers/contracts-and-integrations/deployed-contracts} (2026-08-04)
\item \texttt{https://etherfi.gitbook.io/etherfi/llms.txt} (2026-08-04)
\item \texttt{https://etherfi.gitbook.io/etherfi/solo-stakers/operation-solo-staker/obol-dvt-explained} (2026-08-04)
\end{itemize}

\subsection*{Lido}
\begin{itemize}\setlength{\itemsep}{0pt}\small
\item \texttt{https://docs.lido.fi/contracts/accounting-oracle} (2026-08-04)
\item \texttt{https://docs.lido.fi/deployed-contracts/} (2026-08-04)
\item \texttt{https://lido.fi/how-lido-works/protocol-fee} (2026-08-04)
\item \texttt{https://operatorportal.lido.fi/modules/simple-dvt-module} (2026-08-04)
\item \texttt{https://lido.fi/how-lido-works/lido-node-operators-set-overview} (2026-08-04)
\item \texttt{https://docs.lido.fi/staking-modules/csm/intro} (2026-08-04)
\item \texttt{https://operatorportal.lido.fi/modules/community-staking-module} (2026-08-04)
\item \texttt{https://docs.lido.fi/contracts/deposit-security-module/} (2026-08-04)
\item \texttt{https://docs.lido.fi/contracts/staking-router} (2026-08-04)
\item \texttt{https://research.lido.fi/t/mitigations-for-deposit-front-running-vulnerability/1239} (2026-08-04)
\item \texttt{https://hackmd.io/@lido/stVaults-design} (2026-08-04)
\item \texttt{https://docs.lido.fi/contracts/validators-exit-bus-oracle} (2026-08-04)
\item \texttt{https://github.com/lidofinance/lido-improvement-proposals/blob/develop/LIPS/lip-35.md} (2026-08-04)
\item \texttt{https://research.lido.fi/t/staking-router-v3-design-implementation-proposal-lip-35/11621} (2026-08-04)
\item \texttt{https://docs.lido.fi/contracts/withdrawal-queue-erc721} (2026-08-04)
\item \texttt{https://blog.lido.fi/lido-v3-is-live-modular-infrastructure-for-a-new-paradigm-of-ethereum-staking/} (2026-08-04)
\item \texttt{https://docs.lido.fi/guides/protocol-levers/} (2026-08-04)
\item \texttt{https://github.com/lidofinance/lido-improvement-proposals/blob/develop/LIPS/lip-28.md} (2026-08-04)
\end{itemize}

\subsection*{ApeX}
\begin{itemize}\setlength{\itemsep}{0pt}\small
\item \texttt{https://apex-pro.gitbook.io/apex-pro/apex-omni/perpetual-contracts-open-interest-and-leverage.md} (2026-08-04)
\item \texttt{https://www.apex.exchange/blog/detail/ApeX-Omni-Litepaper} (2026-08-04)
\item \texttt{https://apex-pro.gitbook.io/apex-pro/llms-full.txt} (2026-08-04)
\item \texttt{https://docs.zk.link/} (2026-08-04)
\item \texttt{https://www.apex.exchange/blog/detail/ApeX-Omni-zkLink-Elevating-Multi-chain-Liquidity-Trading-via-zkLink-X} (2026-08-04)
\item \texttt{https://apex-pro.gitbook.io/apex-pro/apex-omni/force-withdrawal.md} (2026-08-04)
\item \texttt{https://apex-pro.gitbook.io/apex-pro/what-is-apex-omni.md} (2026-08-04)
\item \texttt{https://apex-pro.gitbook.io/apex-pro/apex-omni/margin-risk-management-and-liquidations.md} (2026-08-04)
\item \texttt{https://apex-pro.gitbook.io/apex-pro/apex-omni/cross-collateral.md} (2026-08-04)
\item \texttt{https://apex-pro.gitbook.io/apex-pro/apex-omni/price-markers.md} (2026-08-04)
\item \texttt{https://apex-pro.gitbook.io/apex-pro/apex-omni/trading-and-funding-fees.md} (2026-08-04)
\item \texttt{https://apex-pro.gitbook.io/apex-pro/apex-omni/auto-deleveraging-system-adl.md} (2026-08-04)
\item \texttt{https://apex-pro.gitbook.io/apex-pro/apex-omni/protocol-vaults.md} (2026-08-04)
\item \texttt{https://apex-pro.gitbook.io/apex-pro/faqs/apex-protocol-faqs/is-apex-omni-safe-are-your-smart-contracts-audited.md} (2026-08-04)
\end{itemize}

\subsection*{Aster}
\begin{itemize}\setlength{\itemsep}{0pt}\small
\item \texttt{https://docs.asterdex.com/aster-chain/overview.md} (2026-08-04)
\item \texttt{https://docs.asterdex.com/overview/smart-contracts.md} (2026-08-04)
\item \texttt{https://docs.asterdex.com/trading/perpetuals/margin.md} (2026-08-04)
\item \texttt{https://docs.asterdex.com/trading/perpetuals/fees-and-specs/mark-price.md} (2026-08-04)
\item \texttt{https://docs.asterdex.com/trading/1001x/price-oracles.md} (2026-08-04)
\item \texttt{https://docs.asterdex.com/trading/perpetuals/fees-and-specs/funding-rate.md} (2026-08-04)
\item \texttt{https://docs.asterdex.com/trading/perpetuals/fees-and-specs/liquidations.md} (2026-08-04)
\item \texttt{https://docs.asterdex.com/trading/perpetuals/fees-and-specs/auto-deleveraging-adl.md} (2026-08-04)
\item \texttt{https://docs.asterdex.com/earn/overview/aster-alp.md} (2026-08-04)
\item \texttt{https://docs.asterdex.com/trading/shield-mode/liquidations.md} (2026-08-04)
\item \texttt{https://docs.asterdex.com/trading/perpetuals/order-types/hidden-order-and-hidden-position.md} (2026-08-04)
\item \texttt{https://docs.asterdex.com/aster-chain/staking/how-staking-works.md} (2026-08-04)
\item \texttt{https://docs.asterdex.com/aster-chain/governance.md} (2026-08-04)
\item \texttt{https://docs.asterdex.com/llms.txt} (2026-08-04)
\end{itemize}

\subsection*{edgeX}
\begin{itemize}\setlength{\itemsep}{0pt}\small
\item \texttt{https://edgex-1.gitbook.io/edgeX-documentation/edgex-v2/v2-technical-architecture.md} (2026-08-04)
\item \texttt{https://l2beat.com/scaling/projects/edgex} (2026-08-04)
\item \texttt{https://edgex-1.gitbook.io/edgeX-documentation/llms-full.txt} (2026-08-04)
\item \texttt{https://edgex-1.gitbook.io/edgeX-documentation/edgex-v1/self-custody.md} (2026-08-04)
\item \texttt{https://edgex-1.gitbook.io/edgeX-documentation/getting-started/deposits-and-withdrawals/deposits.md} (2026-08-04)
\item \texttt{https://edgex-1.gitbook.io/edgeX-documentation/edgex-v1.md} (2026-08-04)
\item \texttt{https://edgex-1.gitbook.io/edgeX-documentation/trading/trading-accounts-and-margin.md} (2026-08-04)
\item \texttt{https://edgex-1.gitbook.io/edgeX-documentation/api-v2.md} (2026-08-04)
\item \texttt{https://edgex-1.gitbook.io/edgeX-documentation/trading/last-price-index-price-and-oracle-price.md} (2026-08-04)
\item \texttt{https://edgex-1.gitbook.io/edgeX-documentation/trading/funding-fees.md} (2026-08-04)
\item \texttt{https://edgex-1.gitbook.io/edgeX-documentation/trading/liquidation-logic.md} (2026-08-04)
\end{itemize}

\subsection*{Hyperliquid}
\begin{itemize}\setlength{\itemsep}{0pt}\small
\item \texttt{https://hyperliquid.gitbook.io/hyperliquid-docs/hypercore/order-book.md} (2026-08-04)
\item \texttt{https://hyperliquid.gitbook.io/hyperliquid-docs/hypercore/overview.md} (2026-08-04)
\item \texttt{https://github.com/hyperliquid-dex/node} (2026-08-04)
\item \texttt{https://hyperliquid.gitbook.io/hyperliquid-docs/hypercore/bridge.md} (2026-08-04)
\item \texttt{https://hyperliquid.gitbook.io/hyperliquid-docs/hypercore/oracle.md} (2026-08-04)
\item \texttt{https://hyperliquid.gitbook.io/hyperliquid-docs/trading/robust-price-indices.md} (2026-08-04)
\item \texttt{https://hyperliquid.gitbook.io/hyperliquid-docs/trading/funding.md} (2026-08-04)
\item \texttt{https://hyperliquid.gitbook.io/hyperliquid-docs/trading/margin-tiers.md} (2026-08-04)
\item \texttt{https://hyperliquid.gitbook.io/hyperliquid-docs/hypercore/clearinghouse.md} (2026-08-04)
\item \texttt{https://hyperliquid.gitbook.io/hyperliquid-docs/trading/liquidations.md} (2026-08-04)
\item \texttt{https://hyperliquid.gitbook.io/hyperliquid-docs/trading/auto-deleveraging.md} (2026-08-04)
\item \texttt{https://hyperliquid.gitbook.io/hyperliquid-docs/hypercore/vaults/protocol-vaults.md} (2026-08-04)
\item \texttt{https://hyperliquid.gitbook.io/hyperliquid-docs/hypercore/staking.md} (2026-08-04)
\item \texttt{https://hyperliquid.gitbook.io/hyperliquid-docs/hyperliquid-improvement-proposals-hips/hip-3-builder-deployed-perpetuals.md} (2026-08-04)
\item \texttt{https://hyperliquid.gitbook.io/hyperliquid-docs/trading/builder-codes.md} (2026-08-04)
\item \texttt{https://hyperliquid.gitbook.io/hyperliquid-docs/hypercore/vaults/hypercore-vaults-legacy.md} (2026-08-04)
\item \texttt{https://hyperliquid.gitbook.io/hyperliquid-docs/risks.md} (2026-08-04)
\end{itemize}

\subsection*{Lighter}
\begin{itemize}\setlength{\itemsep}{0pt}\small
\item \texttt{https://docs.lighter.xyz/trading/order-types-and-matching.md} (2026-08-04)
\item \texttt{https://assets.lighter.xyz/whitepaper.pdf} (2026-08-04)
\item \texttt{https://l2beat.com/zk-catalog/lighterprover} (2026-08-04)
\item \texttt{https://l2beat.com/scaling/projects/lighter} (2026-08-04)
\item \texttt{https://docs.lighter.xyz/about-lighter/technical-architecture-lighter-core.md} (2026-08-04)
\item \texttt{https://docs.lighter.xyz/trading/liquidations-and-llp-insurance-fund.md} (2026-08-04)
\item \texttt{https://docs.lighter.xyz/trading/multi-asset-margin.md} (2026-08-04)
\item \texttt{https://docs.lighter.xyz/trading/fair-price-marking.md} (2026-08-04)
\item \texttt{https://docs.lighter.xyz/trading/funding.md} (2026-08-04)
\item \texttt{https://docs.lighter.xyz/trading/public-pools.md} (2026-08-04)
\item \texttt{https://docs.lighter.xyz/about-lighter/lit-utility.md} (2026-08-04)
\item \texttt{https://docs.lighter.xyz/trading/trading-fees.md} (2026-08-04)
\item \texttt{https://docs.lighter.xyz/trading/perprfq.md} (2026-08-04)
\end{itemize}

\subsection*{CIAN Yield Layer}
\begin{itemize}\setlength{\itemsep}{0pt}\small
\item \texttt{https://eth.blockscout.com/api?module=contract\&action=getsourcecode\&address=0xB13aa2d0345b0439b064f26B82D8dCf3f508775d} (2026-08-04)
\item \texttt{https://ethereum-rpc.publicnode.com} (2026-08-04)
\item \texttt{https://github.com/Ackee-Blockchain/public-audit-reports/blob/master/2025/ackee-blockchain-cian-yield-layer-report.pdf} (2026-08-04)
\item \texttt{https://docs.cian.app/llms-full.txt} (2026-08-04)
\item \texttt{https://eth.blockscout.com/api?module=contract\&action=getsourcecode\&address=0x49d9409111a6363d82c4371ffa43faea660c917b} (2026-08-04)
\item \texttt{https://eth.blockscout.com/api?module=contract\&action=getsourcecode\&address=0x6d425B3D302DD82cC611866eC8176d435307b616} (2026-08-04)
\end{itemize}

\subsection*{Convex Finance}
\begin{itemize}\setlength{\itemsep}{0pt}\small
\item \texttt{https://docs.convexfinance.com/convexfinance/guides/depositing/curve-lp-tokens.md} (2026-08-04)
\item \texttt{https://github.com/convex-eth/platform/blob/4f22038387ca4ce014dec3d7f5781a8e28051c13/contracts/contracts/Booster.sol} (2026-08-04)
\item \texttt{https://ethereum-rpc.publicnode.com} (2026-08-04)
\item \texttt{https://raw.githubusercontent.com/curvefi/curve-dao-contracts/master/contracts/gauges/LiquidityGaugeV3.vy} (2026-08-04)
\item \texttt{https://docs.convexfinance.com/convexfinance/general-information/why-convex/understanding-cvxcrv.md} (2026-08-04)
\item \texttt{https://github.com/convex-eth/platform/blob/4f22038387ca4ce014dec3d7f5781a8e28051c13/contracts/contracts/CrvDepositor.sol} (2026-08-04)
\item \texttt{https://raw.githubusercontent.com/curvefi/curve-dao-contracts/master/contracts/VotingEscrow.vy} (2026-08-04)
\item \texttt{https://github.com/convex-eth/voting} (2026-08-04)
\item \texttt{https://github.com/convex-eth/platform/blob/4f22038387ca4ce014dec3d7f5781a8e28051c13/contracts/contracts/Cvx.sol} (2026-08-04)
\item \texttt{https://github.com/convex-eth/platform/tree/4f22038387ca4ce014dec3d7f5781a8e28051c13/contracts/contracts} (2026-08-04)
\item \texttt{https://docs.convexfinance.com/convexfinance/faq/multisig-admin-rights.md} (2026-08-04)
\item \texttt{https://eth.blockscout.com/api/v2/smart-contracts/0x989AEb4d175e16225E39E87d0D97A3360524AD80} (2026-08-04)
\end{itemize}

\subsection*{Huma Finance V2}
\begin{itemize}\setlength{\itemsep}{0pt}\small
\item \texttt{https://docs.huma.finance/about-huma/what-is-huma.md} (2026-08-04)
\item \texttt{https://docs.huma.finance/products/huma-2.0/modes.md} (2026-08-04)
\item \texttt{https://docs.huma.finance/products/huma-2.0/lockup.md} (2026-08-04)
\item \texttt{https://docs.huma.finance/products/huma-2.0/redemption.md} (2026-08-04)
\item \texttt{https://docs.huma.finance/products/huma-2.0/overview.md} (2026-08-04)
\item \texttt{https://docs.huma.finance/ecosystem-resources/technical-docs/huma.md} (2026-08-04)
\item \texttt{https://docs.huma.finance/products/huma-institutional/introduction/admin-roles.md} (2026-08-04)
\item \texttt{https://docs.huma.finance/products/huma-institutional/tranches.md} (2026-08-04)
\item \texttt{https://github.com/00labs/huma-contracts-v2/blob/21aaad53023eb3f448ff09820fee1a3936b6175d/contracts/common/PoolConfig.sol\#L283} (2026-08-04)
\item \texttt{https://docs.huma.finance/products/huma-institutional/first-loss-covers.md} (2026-08-04)
\item \texttt{https://github.com/00labs/huma-contracts-v2/blob/21aaad53023eb3f448ff09820fee1a3936b6175d/contracts/credit/CreditLineManager.sol} (2026-08-04)
\item \texttt{https://github.com/00labs/huma-contracts-v2/blob/21aaad53023eb3f448ff09820fee1a3936b6175d/contracts/common/PoolConfig.sol\#L673-L675} (2026-08-04)
\item \texttt{https://github.com/00labs/huma-contracts-v2/blob/21aaad53023eb3f448ff09820fee1a3936b6175d/contracts/credit/ReceivableStorage.sol} (2026-08-04)
\item \texttt{https://docs.huma.finance/products/huma-institutional/introduction/pool-types.md} (2026-08-04)
\item \texttt{https://docs.huma.finance/products/huma-institutional/tranches/redemption.md} (2026-08-04)
\end{itemize}

\subsection*{Pendle}
\begin{itemize}\setlength{\itemsep}{0pt}\small
\item \texttt{https://docs.pendle.finance/pendle-v2-dev/Contracts/StandardizedYield} (2026-08-04)
\item \texttt{https://github.com/pendle-finance/pendle-core-v2-public/blob/e8c2cca4c9b329ba8a383a27d7318e5f8b35c843/contracts/core/YieldContracts/PendleYieldToken.sol\#L243-L249} (2026-08-04)
\item \texttt{https://docs.pendle.finance/pendle-v2/ProtocolMechanics/YieldTokenization/PT} (2026-08-04)
\item \texttt{https://github.com/pendle-finance/pendle-core-v2-public/blob/e8c2cca4c9b329ba8a383a27d7318e5f8b35c843/contracts/core/YieldContracts/PendleYieldToken.sol\#L328-L368} (2026-08-04)
\item \texttt{https://github.com/pendle-finance/pendle-core-v2-public/blob/e8c2cca4c9b329ba8a383a27d7318e5f8b35c843/contracts/core/YieldContracts/InterestManagerYT.sol\#L76} (2026-08-04)
\item \texttt{https://github.com/pendle-finance/pendle-core-v2-public/blob/e8c2cca4c9b329ba8a383a27d7318e5f8b35c843/contracts/core/YieldContracts/PendleYieldToken.sol\#L403-L413} (2026-08-04)
\item \texttt{https://docs.pendle.finance/pendle-v2/ProtocolMechanics/YieldTokenization/YT} (2026-08-04)
\item \texttt{https://github.com/pendle-finance/pendle-core-v2-public/blob/e8c2cca4c9b329ba8a383a27d7318e5f8b35c843/contracts/core/YieldContracts/PendleYieldToken.sol\#L435-L437} (2026-08-04)
\item \texttt{https://github.com/pendle-finance/pendle-core-v2-public/blob/e8c2cca4c9b329ba8a383a27d7318e5f8b35c843/contracts/core/Market/MarketMathCore.sol\#L335-L359} (2026-08-04)
\item \texttt{https://github.com/pendle-finance/pendle-core-v2-public/blob/e8c2cca4c9b329ba8a383a27d7318e5f8b35c843/contracts/core/Market/MarketMathCore.sol\#L276-L312} (2026-08-04)
\item \texttt{https://github.com/pendle-finance/pendle-core-v2-public/blob/e8c2cca4c9b329ba8a383a27d7318e5f8b35c843/contracts/core/Market/MarketMathCore.sol\#L41-L46} (2026-08-04)
\item \texttt{https://github.com/pendle-finance/pendle-core-v2-public/blob/e8c2cca4c9b329ba8a383a27d7318e5f8b35c843/contracts/core/Market/MarketMathCore.sol\#L223} (2026-08-04)
\item \texttt{https://ethereum-rpc.publicnode.com} (2026-08-04)
\item \texttt{https://github.com/pendle-finance/pendle-core-v2-public/blob/e8c2cca4c9b329ba8a383a27d7318e5f8b35c843/contracts/core/YieldContracts/PendleYieldContractFactoryUpg.sol\#L73-L74} (2026-08-04)
\item \texttt{https://github.com/pendle-finance/pendle-core-v2-public/blob/e8c2cca4c9b329ba8a383a27d7318e5f8b35c843/contracts/core/Market/PendleMarketFactoryV7Upg.sol\#L70-L88} (2026-08-04)
\item \texttt{https://raw.githubusercontent.com/pendle-finance/pendle-core-v2-public/e8c2cca4c9b329ba8a383a27d7318e5f8b35c843/deployments/1-core.json} (2026-08-04)
\item \texttt{https://github.com/pendle-finance/pendle-core-v2-public/blob/e8c2cca4c9b329ba8a383a27d7318e5f8b35c843/contracts/LiquidityMining/sPendle/StakedPendle.sol\#L41-L86} (2026-08-04)
\item \texttt{https://github.com/pendle-finance/pendle-core-v2-public/blob/e8c2cca4c9b329ba8a383a27d7318e5f8b35c843/contracts/LiquidityMining/sPendle/StakedPendle.sol\#L116-L127} (2026-08-04)
\item \texttt{https://docs.pendle.finance/pendle-v2-dev/Integration/PointsTracking} (2026-08-04)
\item \texttt{https://github.com/pendle-finance/pendle-core-v2-public/blob/e8c2cca4c9b329ba8a383a27d7318e5f8b35c843/contracts/core/Governance/PendleGovernanceProxy.sol\#L37-L40} (2026-08-04)
\item \texttt{https://github.com/pendle-finance/pendle-core-v2-public/blob/e8c2cca4c9b329ba8a383a27d7318e5f8b35c843/contracts/core/Governance/PendleGovernanceProxy.sol\#L33-L35} (2026-08-04)
\end{itemize}

\subsection*{Spark Savings}
\begin{itemize}\setlength{\itemsep}{0pt}\small
\item \texttt{https://docs.spark.fi/products/spark-savings} (2026-08-04)
\item \texttt{https://github.com/sky-ecosystem/sdai/blob/susds/src/SUsds.sol} (2026-08-04)
\item \texttt{https://github.com/sky-ecosystem/sp-beam/blob/master/src/SPBEAM.sol} (2026-08-04)
\item \texttt{https://ethereum-rpc.publicnode.com} (2026-08-04)
\item \texttt{https://github.com/sparkdotfi/xchain-ssr-oracle} (2026-08-04)
\item \texttt{https://github.com/sky-ecosystem/sdai/blob/susds/src/l2/SUsds.sol} (2026-08-04)
\item \texttt{https://github.com/sparkdotfi/spark-alm-controller/blob/dev/src/MainnetController.sol} (2026-08-04)
\item \texttt{https://docs.spark.fi/products/spark-liquidity-layer} (2026-08-04)
\item \texttt{https://github.com/sparkdotfi/spark-alm-controller/blob/dev/src/RateLimits.sol} (2026-08-04)
\item \texttt{https://github.com/sparkdotfi/spark-address-registry/blob/master/src/Ethereum.sol} (2026-08-04)
\item \texttt{https://eth.blockscout.com/api/v2/smart-contracts/0xa3931d71877C0E7a3148CB7Eb4463524FEc27fbD} (2026-08-04)
\end{itemize}

\subsection*{BTCB}
\begin{itemize}\setlength{\itemsep}{0pt}\small
\item \texttt{https://www.bnbchain.org/en/blog/btcb-on-binance-smart-chain-101} (2026-08-04)
\item \texttt{https://bsc-rpc.publicnode.com} (2026-08-04)
\item \texttt{https://www.binance.com/en/proof-of-reserves} (2026-08-04)
\item \texttt{https://bscscan.com/address/0x7130d2a12b9bcbfae4f2634d864a1ee1ce3ead9c\#code} (2026-08-04)
\item \texttt{https://bscscan.com/token/0x7130d2a12b9bcbfae4f2634d864a1ee1ce3ead9c} (2026-08-04)
\item \texttt{https://www.binance.com/en/proof-of-collateral} (2026-08-04)
\item \texttt{https://www.binance.com/en/collateral-btokens} (2026-08-04)
\end{itemize}

\subsection*{Coinbase Bridge}
\begin{itemize}\setlength{\itemsep}{0pt}\small
\item \texttt{https://help.coinbase.com/en/coinbase/trading-and-funding/sending-or-receiving-cryptocurrency/coinbase-wrapped-btc} (2026-08-04)
\item \texttt{https://www.coinbase.com/blog/coinbase-wrapped-btc-cbbtc-is-now-live} (2026-08-04)
\item \texttt{https://www.coinbase.com/cbbtc} (2026-08-04)
\item \texttt{https://www.coinbase.com/cbbtc/proof-of-reserves} (2026-08-04)
\item \texttt{https://ethereum-rpc.publicnode.com} (2026-08-04)
\item \texttt{https://raw.githubusercontent.com/DefiLlama/DefiLlama-Adapters/main/projects/coinbase-btc/index.js} (2026-08-04)
\item \texttt{https://api.llama.fi/protocol/coinbase-bridge} (2026-08-04)
\item \texttt{https://raw.githubusercontent.com/coinbase/wrapped-tokens-os/master/README.md} (2026-08-04)
\item \texttt{https://basescan.org/address/0x7458bfDC30034EB860B265E6068121D18Fa5Aa72\#code} (2026-08-04)
\item \texttt{https://basescan.org/address/0xcbb7c0000ab88b473b1f5afd9ef808440eed33bf} (2026-08-04)
\item \texttt{https://api.github.com/repos/coinbase/wrapped-tokens-os} (2026-08-04)
\end{itemize}

\subsection*{Hyperliquid Bridge}
\begin{itemize}\setlength{\itemsep}{0pt}\small
\item \texttt{https://hyperliquid.gitbook.io/hyperliquid-docs/for-developers/api/bridge2} (2026-08-04)
\item \texttt{https://arbitrum-one-rpc.publicnode.com} (2026-08-04)
\item \texttt{https://hyperliquid.gitbook.io/hyperliquid-docs/hypercore/bridge.md} (2026-08-04)
\item \texttt{https://raw.githubusercontent.com/hyperliquid-dex/contracts/master/Bridge2.sol} (2026-08-04)
\item \texttt{https://api.llama.fi/protocol/hyperliquid-bridge} (2026-08-04)
\item \texttt{https://www.circle.com/blog/native-usdc-cctp-v2-are-coming-to-hyperliquid-what-you-need-to-know} (2026-08-04)
\item \texttt{https://hyperliquid.gitbook.io/hyperliquid-docs/hypercore/staking.md} (2026-08-04)
\item \texttt{https://api.hyperliquid.xyz/info} (2026-08-04)
\end{itemize}

\subsection*{LayerZero V2}
\begin{itemize}\setlength{\itemsep}{0pt}\small
\item \texttt{https://docs.layerzero.network/v2/concepts/protocol/layerzero-endpoint} (2026-08-04)
\item \texttt{https://ethereum-rpc.publicnode.com} (2026-08-04)
\item \texttt{https://docs.layerzero.network/v2/concepts/modular-security/security-stack-dvns} (2026-08-04)
\item \texttt{https://docs.layerzero.network/v2/home/protocol/protocol-overview} (2026-08-04)
\item \texttt{https://docs.layerzero.network/v2/concepts/glossary} (2026-08-04)
\item \texttt{https://www.coindesk.com/tech/2026/04/20/kelp-dao-claims-layerzero-s-default-settings-are-what-actually-caused-the-usd290-million-disaster} (2026-08-04)
\item \texttt{https://www.chainalysis.com/blog/kelpdao-bridge-exploit-april-2026/} (2026-08-04)
\item \texttt{https://metadata.layerzero-api.com/v1/metadata/deployments} (2026-08-04)
\item \texttt{https://docs.layerzero.network/v2/home/token-standards/oft-standard} (2026-08-04)
\item \texttt{https://api.llama.fi/protocol/layerzero} (2026-08-04)
\end{itemize}

\subsection*{WBTC}
\begin{itemize}\setlength{\itemsep}{0pt}\small
\item \texttt{https://docs.wbtc.network/how-wbtc-works/custody-and-security} (2026-08-04)
\item \texttt{https://docs.wbtc.network/how-wbtc-works/mint-burn-mechanism} (2026-08-04)
\item \texttt{https://ethereum-rpc.publicnode.com} (2026-08-04)
\item \texttt{https://raw.githubusercontent.com/WrappedBTC/bitcoin-token-smart-contracts/master/ethereum/contracts/factory/Factory.sol} (2026-08-04)
\item \texttt{https://www.wbtc.network/transparency} (2026-08-04)
\item \texttt{https://raw.githubusercontent.com/WrappedBTC/DAO/master/MerchantGuide.md} (2026-08-04)
\item \texttt{https://raw.githubusercontent.com/DefiLlama/DefiLlama-Adapters/main/projects/wbtc.js} (2026-08-04)
\item \texttt{https://raw.githubusercontent.com/WrappedBTC/bitcoin-token-smart-contracts/master/ethereum/contracts/controller/Controller.sol} (2026-08-04)
\item \texttt{https://raw.githubusercontent.com/WrappedBTC/bitcoin-token-smart-contracts/master/ethereum/contracts/token/WBTC.sol} (2026-08-04)
\item \texttt{https://raw.githubusercontent.com/WrappedBTC/DAO/master/README.md} (2026-08-04)
\item \texttt{https://raw.githubusercontent.com/WrappedBTC/DAO/master/DeploymentVerification.md} (2026-08-04)
\item \texttt{https://docs.wbtc.network/resources/contract-addresses} (2026-08-04)
\item \texttt{https://www.coindesk.com/business/2026/08/04/bitgo-s-wbtc-move-pushes-layerzero-to-chainlink-tally-near-usd15-billion} (2026-08-04)
\end{itemize}

\subsection*{Binance Wallet}
\begin{itemize}\setlength{\itemsep}{0pt}\small
\item \texttt{https://www.binance.com/en-BH/support/faq/what-is-binance-web3-wallet-swap-7ebe19e8ff884cc09a9dbb064aff131f} (2026-08-05)
\item \texttt{https://www.binance.com/en/support/faq/what-is-binance-wallet-swap-7ebe19e8ff884cc09a9dbb064aff131f} (2026-08-05)
\item \texttt{https://www.binance.com/en/support/announcement/detail/4037408455264c02becd90b5873121c3} (2026-08-05)
\item \texttt{https://www.bnbchain.org/en/blog/protecting-users-from-sandwich-attacks-bnb-chain-introduces-mev-protection-with-several-wallets} (2026-08-05)
\item \texttt{https://bscscan.com/address/0xb300000b72deaeb607a12d5f54773d1c19c7028d} (2026-08-05)
\item \texttt{https://www.binance.com/en/support/faq/frequently-asked-questions-on-binance-web3-wallet-5a3fc86a702b43e1a4a398ebd8853b77} (2026-08-05)
\item \texttt{https://academy.binance.com/en/articles/what-is-binance-alpha} (2026-08-05)
\item \texttt{https://www.binance.com/en/support/announcement/detail/c58cbaf282ac45819f470176bdc17f3b} (2026-08-05)
\item \texttt{https://github.com/DefiLlama/dimension-adapters/blob/master/aggregators/binancewallet/index.ts} (2026-08-05)
\end{itemize}

\subsection*{Jupiter}
\begin{itemize}\setlength{\itemsep}{0pt}\small
\item \texttt{https://developers.jup.ag/docs/swap/} (2026-08-05)
\item \texttt{https://developers.jup.ag/docs/blog/ultra-v3} (2026-08-05)
\item \texttt{https://developers.jup.ag/docs/swap/routing} (2026-08-05)
\item \texttt{https://developers.jup.ag/docs/swap/fees} (2026-08-05)
\item \texttt{https://developers.jup.ag/docs/trigger/} (2026-08-05)
\item \texttt{https://developers.jup.ag/docs/llms.txt} (2026-08-05)
\item \texttt{https://github.com/jup-ag/metis-binary} (2026-08-05)
\end{itemize}

\subsection*{KyberSwap}
\begin{itemize}\setlength{\itemsep}{0pt}\small
\item \texttt{https://docs.kyberswap.com/kyberswap-solutions/kyberswap-aggregator} (2026-08-05)
\item \texttt{https://docs.kyberswap.com/developer-guide/aggregator-api/aggregator-api-specification/evm-swaps.md} (2026-08-05)
\item \texttt{https://docs.kyberswap.com/developer-guide/start-here/foundational-solutions/smart-settlement-better-swap-output-with-lower-slippage.md} (2026-08-05)
\item \texttt{https://github.com/KyberNetwork/kyber-exclusive-amm-sc} (2026-08-05)
\item \texttt{https://docs.kyberswap.com/user-guide/kyberswap-fairflow/fairflow-details.md} (2026-08-05)
\item \texttt{https://docs.kyberswap.com/user-guide/kyberswap-fairflow/solution-fairflow.md} (2026-08-05)
\item \texttt{https://docs.kyberswap.com/getting-started/fee-schedule.md} (2026-08-05)
\item \texttt{https://github.com/KyberNetwork/smart-intent-sc} (2026-08-05)
\item \texttt{https://docs.kyberswap.com/llms.txt} (2026-08-05)
\item \texttt{https://github.com/KyberNetwork/kyberswap-dex-lib} (2026-08-05)
\end{itemize}

\subsection*{LiquidMesh}
\begin{itemize}\setlength{\itemsep}{0pt}\small
\item \texttt{https://docs.liquidmesh.io/reference/quote-1.md} (2026-08-05)
\item \texttt{https://docs.liquidmesh.io/docs/getting-started.md} (2026-08-05)
\item \texttt{https://docs.liquidmesh.io/docs/quote-api.md} (2026-08-05)
\item \texttt{https://docs.liquidmesh.io/docs/boost-api.md} (2026-08-05)
\item \texttt{https://liquidmesh.io/} (2026-08-05)
\item \texttt{https://docs.liquidmesh.io/docs/access-eligibility-and-ip-screening.md} (2026-08-05)
\item \texttt{https://docs.liquidmesh.io/changelog/liquidmesh-security-update.md} (2026-08-05)
\item \texttt{https://docs.liquidmesh.io/docs/overview.md} (2026-08-05)
\item \texttt{https://docs.liquidmesh.io/docs/smart-contracts.md} (2026-08-05)
\item \texttt{https://bscscan.com/address/0x3d90f66B534Dd8482b181e24655A9e8265316BE9\#code} (2026-08-05)
\item \texttt{https://docs.liquidmesh.io/docs/supported-chains.md} (2026-08-05)
\end{itemize}

\subsection*{OKX DEX}
\begin{itemize}\setlength{\itemsep}{0pt}\small
\item \texttt{https://web3.okx.com/onchainos/dev-docs/trade/dex-swap-api-introduction} (2026-08-05)
\item \texttt{https://www.okx.com/en-eu/help/what-is-okx-dex-aggregator-plus} (2026-08-05)
\item \texttt{https://www.okx.com/en-us/help/dex-aggregator-and-intent-conflict-of-interest-policies} (2026-08-05)
\item \texttt{https://web3.okx.com/onchainos/dev-docs/trade/dex-api-addfee} (2026-08-05)
\item \texttt{https://developers.jup.ag/docs/swap/routing} (2026-08-05)
\item \texttt{https://web3.okx.com/onchainos/dev-docs/trade/dex-smart-contract} (2026-08-05)
\item \texttt{https://web3.okx.com/onchainos/dev-docs/home/change-log} (2026-08-05)
\item \texttt{https://github.com/okx/dex-solana-binary} (2026-08-05)
\item \texttt{https://github.com/Julian-dev28/WEB3-DEX-OPENSOURCE} (2026-08-05)
\end{itemize}

\subsection*{BlackRock BUIDL}
\begin{itemize}\setlength{\itemsep}{0pt}\small
\item \texttt{https://www.sec.gov/Archives/edgar/data/2013810/000201381026000002/primary\_doc.xml} (2026-08-04)
\item \texttt{https://securitize.io/blackrock/buidl} (2026-08-04)
\item \texttt{https://raw.githubusercontent.com/securitize-io/DSTokenInterfaces/master/contracts/dsprotocol/registry/DSRegistryServiceInterface.sol} (2026-08-04)
\item \texttt{https://raw.githubusercontent.com/securitize-io/DSTokenInterfaces/master/contracts/dsprotocol/token/DSTokenInterface.sol} (2026-08-04)
\item \texttt{https://www.sec.gov/Archives/edgar/data/1782266/000190359625000593/primary\_doc.xml} (2026-08-04)
\end{itemize}

\subsection*{Centrifuge}
\begin{itemize}\setlength{\itemsep}{0pt}\small
\item \texttt{https://docs.centrifuge.io/user/concepts/share-tokens.md} (2026-08-04)
\item \texttt{https://docs.centrifuge.io/user/concepts/pools.md} (2026-08-04)
\item \texttt{https://docs.centrifuge.io/developer/protocol/guides/publish-nav-and-oracles.md} (2026-08-04)
\item \texttt{https://raw.githubusercontent.com/centrifuge/protocol/v3.2.0/src/core/hub/Hub.sol} (2026-08-04)
\item \texttt{https://docs.centrifuge.io/developer/protocol/features/token-compliance.md} (2026-08-04)
\item \texttt{https://raw.githubusercontent.com/centrifuge/protocol/v3.2.0/src/core/spoke/BalanceSheet.sol} (2026-08-04)
\item \texttt{https://centrifuge.io/blog/centrifuge-derwa-tokens} (2026-08-04)
\item \texttt{https://docs.centrifuge.io/developer/security/pool-access-levels.md} (2026-08-04)
\item \texttt{https://docs.centrifuge.io/developer/security/guardian.md} (2026-08-04)
\item \texttt{https://www.anemoy.io/funds/jtrsy} (2026-08-04)
\end{itemize}

\subsection*{Circle USYC}
\begin{itemize}\setlength{\itemsep}{0pt}\small
\item \texttt{https://usyc.docs.hashnote.com/overview/product-structuring.md} (2026-08-04)
\item \texttt{https://usyc.docs.hashnote.com/overview/token-price.md} (2026-08-04)
\item \texttt{https://usyc.docs.hashnote.com/overview/subscription-and-redemption.md} (2026-08-04)
\item \texttt{https://usyc.docs.hashnote.com/overview/investor-onboarding.md} (2026-08-04)
\item \texttt{https://usyc.docs.hashnote.com/overview/smart-contracts.md} (2026-08-04)
\item \texttt{https://www.circle.com/pressroom/circles-usyc-now-supported-as-yield-bearing-off-exchange-collateral-for-binances-institutional-clients} (2026-08-04)
\item \texttt{https://usyc.docs.hashnote.com/overview/fees.md} (2026-08-04)
\item \texttt{https://www.sec.gov/Archives/edgar/data/1876042/000187604226000062/crcl-20251231.htm} (2026-08-04)
\item \texttt{https://usyc.docs.hashnote.com/overview/service-providers.md} (2026-08-04)
\item \texttt{https://developers.circle.com/tokenized/usyc/overview.md} (2026-08-04)
\end{itemize}

\subsection*{Maple Finance}
\begin{itemize}\setlength{\itemsep}{0pt}\small
\item \texttt{https://docs.maple.finance/technical-resources/protocol-overview/smart-contract-architecture.md} (2026-08-04)
\item \texttt{https://docs.maple.finance/integrate/technical-resources/collateral-and-yield-disclosure.md} (2026-08-04)
\item \texttt{https://docs.maple.finance/maple-institutional-for-lenders/defaults-and-impairments-1.md} (2026-08-04)
\item \texttt{https://docs.maple.finance/technical-resources/loans/defaults.md} (2026-08-04)
\item \texttt{https://docs.maple.finance/technical-resources/loans/impairments.md} (2026-08-04)
\item \texttt{https://docs.maple.finance/syrupusdc-usdt-usdg-for-lenders/risk.md} (2026-08-04)
\item \texttt{https://docs.maple.finance/technical-resources/singletons/pool-permission-manager.md} (2026-08-04)
\item \texttt{https://docs.maple.finance/maple-for-token-holders/syrup-tokenomics/staking.md} (2026-08-04)
\item \texttt{https://docs.maple.finance/technical-resources/admin-functions/timelocks.md} (2026-08-04)
\item \texttt{https://docs.maple.finance/technical-resources/crosschain.md} (2026-08-04)
\item \texttt{https://docs.maple.finance/maple-institutional-for-lenders/risk.md} (2026-08-04)
\item \texttt{https://docs.maple.finance/maple-institutional-for-lenders/lending.md} (2026-08-04)
\end{itemize}

\subsection*{Ondo Finance}
\begin{itemize}\setlength{\itemsep}{0pt}\small
\item \texttt{https://docs.ondo.finance/trust-and-security.md} (2026-08-04)
\item \texttt{https://docs.ondo.finance/developer-guides/usdy-instant-manager-integration.md} (2026-08-04)
\item \texttt{https://docs.ondo.finance/qualified-access-products/ousg/yield.md} (2026-08-04)
\item \texttt{https://docs.ondo.finance/general-access-products/usdy/rebasing.md} (2026-08-04)
\item \texttt{https://docs.ondo.finance/qualified-access-products/ousg/overview.md} (2026-08-04)
\item \texttt{https://docs.ondo.finance/ondo-stocks/token-and-quote-pricing.md} (2026-08-04)
\item \texttt{https://docs.ondo.finance/ondo-stocks/investing-and-redeeming.md} (2026-08-04)
\item \texttt{https://docs.ondo.finance/ondo-stocks/secondary-market-restrictions.md} (2026-08-04)
\item \texttt{https://docs.ondo.finance/qualified-access-products/ousg/instant-limits.md} (2026-08-04)
\item \texttt{https://docs.ondo.finance/addresses.md} (2026-08-04)
\item \texttt{https://eth.blockscout.com/api/v2/smart-contracts/0x5cd9e3a4c9933133b512da1b6ba4672160e0c665} (2026-08-04)
\item \texttt{https://github.com/ondoprotocol/usdy/blob/3912ca0698c2992e4db997d0855e62588c44e2c0/contracts/usdy/USDY.sol} (2026-08-04)
\item \texttt{https://docs.ondo.finance/tools/ondo-bridge.md} (2026-08-04)
\item \texttt{https://docs.ondo.finance/ondo-stocks/trust-and-transparency.md} (2026-08-04)
\item \texttt{https://docs.ondo.finance/general-access-products/usdy/important-notes.md} (2026-08-04)
\item \texttt{https://web.archive.org/web/20251013000000id\_/https://docs.ondo.finance/general-access-products/usdy/faq/investing-and-redeeming} (2026-08-04)
\item \texttt{https://web.archive.org/web/20251013000000id\_/https://docs.ondo.finance/general-access-products/usdy/faq/trust-and-transparency} (2026-08-04)
\item \texttt{https://ondo.finance/usdy} (2026-08-04)
\end{itemize}

\subsection*{Aevo}
\begin{itemize}\setlength{\itemsep}{0pt}\small
\item \texttt{https://docs.aevo.xyz/aevo-products/aevo-exchange/trading-on-aevo/options-specifications/eth-options.md} (2026-08-04)
\item \texttt{https://api.aevo.xyz/markets?asset=ETH\&instrument\_type=OPTION} (2026-08-04)
\item \texttt{https://docs.aevo.xyz/aevo-products/aevo-exchange/technical-architecture/margin-framework.md} (2026-08-04)
\item \texttt{https://docs.aevo.xyz/llms.txt} (2026-08-04)
\item \texttt{https://www.theblock.co/post/382461/aevos-legacy-ribbon-dov-vaults-exploited-for-2-7-million-following-oracle-upgrade} (2026-08-04)
\item \texttt{https://docs.aevo.xyz/aevo-products/aevo-exchange/technical-architecture/exchange-structure/off-chain-orderbook-and-risk-engine.md} (2026-08-04)
\item \texttt{https://docs.aevo.xyz/aevo-products/aevo-exchange/technical-architecture/margin-framework/portfolio-margin.md} (2026-08-04)
\item \texttt{https://docs.aevo.xyz/aevo-products/aevo-exchange/technical-architecture/liquidations.md} (2026-08-04)
\item \texttt{https://docs.aevo.xyz/aevo-products/aevo-exchange/technical-architecture/auto-deleveraging-adl.md} (2026-08-04)
\item \texttt{https://docs.aevo.xyz/aevo-products/aevo-exchange/technical-architecture/index-price.md} (2026-08-04)
\item \texttt{https://docs.aevo.xyz/aevo-products/aevo-exchange/technical-architecture/exchange-structure/layer-2-architecture.md} (2026-08-04)
\end{itemize}

\subsection*{Derive}
\begin{itemize}\setlength{\itemsep}{0pt}\small
\item \texttt{https://docs.derive.xyz/docs/supported-products-1.md} (2026-08-04)
\item \texttt{https://api.lyra.finance/public/get\_instruments} (2026-08-04)
\item \texttt{https://docs.derive.xyz/docs/settlements.md} (2026-08-04)
\item \texttt{https://docs.derive.xyz/docs/oracles-1.md} (2026-08-04)
\item \texttt{https://github.com/derivexyz/v2-core/tree/96796a61dcb1dc852e25518b00cc1a79fb3caeeb/src/feeds} (2026-08-04)
\item \texttt{https://docs.derive.xyz/docs/pm2.md} (2026-08-04)
\item \texttt{https://docs.derive.xyz/reference/overview} (2026-08-04)
\item \texttt{https://docs.derive.xyz/docs/self-custodial-withdrawals-escape-hatch.md} (2026-08-04)
\item \texttt{https://docs.derive.xyz/reference/fees-1.md} (2026-08-04)
\item \texttt{https://docs.derive.xyz/docs/portfolio-margin-1.md} (2026-08-04)
\item \texttt{https://github.com/derivexyz/v2-core/tree/96796a61dcb1dc852e25518b00cc1a79fb3caeeb/src/risk-managers} (2026-08-04)
\item \texttt{https://api.lyra.finance/public/get\_all\_currencies} (2026-08-04)
\item \texttt{https://docs.derive.xyz/docs/liquidations-1.md} (2026-08-04)
\item \texttt{https://github.com/derivexyz/v2-core/tree/96796a61dcb1dc852e25518b00cc1a79fb3caeeb} (2026-08-04)
\item \texttt{https://docs.derive.xyz/docs/staking-rewards-program.md} (2026-08-04)
\item \texttt{https://github.com/derivexyz/v2-core/tree/96796a61dcb1dc852e25518b00cc1a79fb3caeeb/src/assets} (2026-08-04)
\item \texttt{https://docs.derive.xyz/docs/how-do-i-know-my-funds-are-safe.md} (2026-08-04)
\item \texttt{https://docs.derive.xyz/docs/how-are-strikes-and-expiries-selected.md} (2026-08-04)
\end{itemize}

\subsection*{Hegic}
\begin{itemize}\setlength{\itemsep}{0pt}\small
\item \texttt{https://github.com/hegic/contracts/tree/54833086acc17d2a2bf90229040c456f2fda2624} (2026-08-04)
\end{itemize}

\subsection*{Panoptic}
\begin{itemize}\setlength{\itemsep}{0pt}\small
\item \texttt{https://github.com/panoptic-labs/panoptic-v2-core\#readme} (2026-08-04)
\item \texttt{https://panoptic.xyz/docs/contracts/smart-contracts-overview} (2026-08-04)
\item \texttt{https://panoptic.xyz/docs/faq/} (2026-08-04)
\item \texttt{https://panoptic.xyz/docs/product/streamia} (2026-08-04)
\item \texttt{https://panoptic.xyz/docs/panoptic-protocol/liquidations} (2026-08-04)
\item \texttt{https://panoptic.xyz/blog/panoptic-v2-is-coming} (2026-08-04)
\item \texttt{https://code4rena.com/audits/2025-12-panoptic-next-core} (2026-08-04)
\item \texttt{https://github.com/panoptic-labs/panoptic-v2-core/blob/d65310d6cfbaadb6910fa9446cc59c6541060749/contracts/RiskEngine.sol} (2026-08-04)
\item \texttt{https://github.com/panoptic-labs/panoptic-v2-core/tree/d65310d6cfbaadb6910fa9446cc59c6541060749/audits} (2026-08-04)
\end{itemize}

\subsection*{Rysk}
\begin{itemize}\setlength{\itemsep}{0pt}\small
\item \texttt{https://docs.rysk.finance/getting-started/solution-rysk-v12} (2026-08-04)
\item \texttt{https://docs.rysk.finance/getting-started/protocol-and-product/makers-how-to-integrate-in-the-rfq.md} (2026-08-04)
\item \texttt{https://docs.rysk.finance/getting-started/protocol-and-product/how-it-works.md} (2026-08-04)
\item \texttt{https://docs.rysk.finance/getting-started/protocol-and-product/faq.md} (2026-08-04)
\item \texttt{https://docs.rysk.finance/getting-started/protocol-and-product/products.md} (2026-08-04)
\item \texttt{https://docs.rysk.finance/rysk-premium/rysk-premium-explainer.md} (2026-08-04)
\item \texttt{https://sourcify.dev/server/v2/contract/1/0x684404F2AEBAD87a6803F13741B1d638Bfe2C671?fields=compilation} (2026-08-04)
\item \texttt{https://docs.rysk.finance/} (2026-08-04)
\item \texttt{https://docs.rysk.finance/resources/security.md} (2026-08-04)
\item \texttt{https://docs.rysk.finance/resources/oracle.md} (2026-08-04)
\item \texttt{https://api.llama.fi/protocols} (2026-08-04)
\item \texttt{https://docs.rysk.finance/resources/important-contracts.md} (2026-06-10)
\end{itemize}

\subsection*{Circle USDC}
\begin{itemize}\setlength{\itemsep}{0pt}\small
\item \texttt{https://www.sec.gov/Archives/edgar/data/1876042/000187604226000062/crcl-20251231.htm} (2026-08-04)
\item \texttt{https://6778953.fs1.hubspotusercontent-na1.net/hubfs/6778953/USDCAttestationReports/2026/2026\%20USDC\_Examination\%20Report\%20June\%2026.pdf} (2026-08-04)
\item \texttt{https://www.circle.com/legal/usdc-terms} (2026-08-04)
\item \texttt{https://www.circle.com/transparency} (2026-08-04)
\item \texttt{https://eth.blockscout.com/api/v2/smart-contracts/0x43506849D7C04F9138D1A2050bbF3A0c054402dd} (2026-08-04)
\item \texttt{https://eth.blockscout.com/api/v2/smart-contracts/0xA0b86991c6218b36c1d19D4a2E9Eb0cE3606eB48} (2026-08-04)
\item \texttt{https://ethereum-rpc.publicnode.com} (2026-08-04)
\item \texttt{https://github.com/circlefin/stablecoin-evm} (2026-08-04)
\end{itemize}

\subsection*{Global Dollar USDG}
\begin{itemize}\setlength{\itemsep}{0pt}\small
\item \texttt{https://www.paxos.com/terms-and-conditions/usdg-eu-whitepaper} (2026-08-04)
\item \texttt{https://globaldollar.com/about-usdg} (2026-08-04)
\item \texttt{https://eth.blockscout.com/api/v2/smart-contracts/0xfacd5ff359adf87822374275699dd518aaf9a65f} (2026-08-04)
\item \texttt{https://github.com/paxosglobal/paxos-token-contracts} (2026-08-04)
\item \texttt{https://www.paxos.com/usdg-transparency} (2026-08-04)
\item \texttt{https://globaldollar.com/earn-with-usdg} (2026-08-04)
\item \texttt{https://ethereum-rpc.publicnode.com} (2026-08-04)
\item \texttt{https://github.com/paxosglobal/usdg-contract} (2026-08-04)
\end{itemize}

\subsection*{PayPal USD}
\begin{itemize}\setlength{\itemsep}{0pt}\small
\item \texttt{https://www.sec.gov/Archives/edgar/data/0001633917/000163391726000082/pypl-20260630.htm} (2026-08-04)
\item \texttt{https://www.paxos.com/pyusd} (2026-08-04)
\item \texttt{https://www.paxos.com/pyusd-transparency} (2026-08-04)
\item \texttt{https://www.paxos.com/newsroom/occ-approves-paxos-application-to-convert-to-occ-trust-paxos-to-complete-conversion-imminently-to-become-a-federally-regulated-blockchain-infrastructure-provider} (2026-08-04)
\item \texttt{https://www.occ.gov/news-issuances/news-releases/2025/nr-occ-2025-125.html} (2026-08-04)
\item \texttt{https://eth.blockscout.com/api/v2/smart-contracts/0x8c35caa5fd5bdc64b6b11344ad57594a3676256a} (2026-08-04)
\item \texttt{https://ethereum-rpc.publicnode.com} (2026-08-04)
\item \texttt{https://github.com/paxosglobal/paxos-token-contracts} (2026-08-04)
\item \texttt{https://github.com/paxosglobal/pyusd-contract} (2026-08-04)
\item \texttt{https://eth.blockscout.com/api/v2/smart-contracts/0xfacd5ff359adf87822374275699dd518aaf9a65f} (2026-08-04)
\item \texttt{https://docs.paxos.com/guides/stablecoin/pyusd/mainnet} (2026-08-04)
\item \texttt{https://www.paypal.com/us/digital-wallet/manage-money/crypto/pyusd} (2026-08-04)
\end{itemize}

\subsection*{Tether USDT}
\begin{itemize}\setlength{\itemsep}{0pt}\small
\item \texttt{https://tether.to/public/Relevant\_Information\_Document\_-\_Tether\_International,\_S.A.\_de\_C.V..pdf} (2026-08-04)
\item \texttt{https://assets.ctfassets.net/vyse88cgwfbl/2kYf7r64h3tzwiu6F0CbUB/2997abd2f11ecea74a21528048b50707/Opinion\_\_\_Report\_-\_Tether\_International\_Financial\_Figure\_30-06-2026.pdf} (2026-08-04)
\item \texttt{https://tether.io/news/tether-posts-strong-q2-performance-generates-1-5b-net-operating-profit-maintains-4-11b-reserve-buffer-and-expands-gold-holdings-to-more-than-146-tons/} (2026-08-04)
\item \texttt{https://eth.blockscout.com/api/v2/smart-contracts/0xdAC17F958D2ee523a2206206994597C13D831ec7} (2026-08-04)
\item \texttt{https://ethereum-rpc.publicnode.com} (2026-08-04)
\item \texttt{https://tether.io/news/tether-supports-freeze-of-more-than-344-million-in-usdt-in-coordination-with-ofac-and-u-s-law-enforcement/} (2026-08-04)
\item \texttt{https://tether.io/news/tether-announces-the-launch-of-usat-the-federally-regulated-dollar-backed-stablecoin-made-in-america/} (2026-08-04)
\item \texttt{https://github.com/tethercoin/USDT} (2026-08-04)
\end{itemize}

\subsection*{USD1}
\begin{itemize}\setlength{\itemsep}{0pt}\small
\item \texttt{https://landing.bitgo.com/rs/552-OGK-141/images/USD1\%5FReserve\%5FAttestation\%5FReport\%5FJune\%5F2026.pdf?version=0} (2026-08-04)
\item \texttt{https://docs.worldlibertyfinancial.com/usd1-token/what-is-usd1} (2026-08-04)
\item \texttt{https://www.bitgo.com/legal/bitgo-coin-minting-services-terms/} (2026-08-04)
\item \texttt{https://fintel.io/so/us/31607a703} (2026-08-04)
\item \texttt{https://www.bitgo.com/usd1-attestations/} (2026-08-04)
\item \texttt{https://docs.worldlibertyfinancial.com/usd1-token/attestation-reports} (2026-08-04)
\item \texttt{https://eth.blockscout.com/api/v2/smart-contracts/0x694aa534bdef8ed63244eb902e7914e527891f08} (2026-08-04)
\item \texttt{https://eth.blockscout.com/api/v2/smart-contracts/0xee9b1a09aedaced9dcda74964ea447feb93861c2} (2026-08-04)
\item \texttt{https://ethereum-rpc.publicnode.com} (2026-08-04)
\item \texttt{https://docs.worldlibertyfinancial.com/usd1-token/contract-addresses} (2026-08-04)
\end{itemize}

\subsection*{Azuro}
\begin{itemize}\setlength{\itemsep}{0pt}\small
\item \texttt{https://gem.azuro.org/hub/blockchains/architecture} (2026-08-04)
\item \texttt{https://gem.azuro.org/knowledge-hub/how-azuro-works/components/pools} (2026-08-04)
\item \texttt{https://gem.azuro.org/knowledge-hub/how-azuro-works/components/vAMM} (2026-08-04)
\item \texttt{https://gem.azuro.org/knowledge-hub/how-azuro-works/components/odds} (2026-08-04)
\item \texttt{https://gem.azuro.org/knowledge-hub/how-azuro-works/components/reinforcement} (2026-08-04)
\item \texttt{https://gem.azuro.org/knowledge-hub/how-azuro-works/liquidity-tree} (2026-08-04)
\item \texttt{https://gem.azuro.org/knowledge-hub/how-azuro-works/protocol-actors/liquidity-providers} (2026-08-04)
\item \texttt{https://gem.azuro.org/knowledge-hub/how-azuro-works/protocol-actors/data-providers} (2026-08-04)
\item \texttt{https://gem.azuro.org/knowledge-hub/how-azuro-works/components/conditions} (2026-08-04)
\item \texttt{https://gem.azuro.org/hub/releases/29-cashouts} (2026-08-04)
\item \texttt{https://gem.azuro.org/knowledge-hub/how-azuro-works/components/betting-engines} (2026-08-04)
\item \texttt{https://gem.azuro.org/knowledge-hub/how-azuro-works/reward-distribution} (2026-08-04)
\end{itemize}

\subsection*{Grove Finance}
\begin{itemize}\setlength{\itemsep}{0pt}\small
\item \texttt{https://www.grove.finance/blog/inside-grove-onchain-allocation-framework} (2026-08-04)
\item \texttt{https://docs.grove.finance/grove-allocator} (2026-08-04)
\item \texttt{https://centrifuge.io/blog/grove-avalanche-centrifuge-jaaa} (2026-08-04)
\item \texttt{https://www.businesswire.com/news/home/20250624393365/en/Grove-Announces-Launch-of-Institutional-Grade-Credit-Infrastructure-DeFi-Protocol-with-\$1-Billion-Allocation-to-Tokenized-Janus-Henderson-Anemoy-AAA-CLO-Strategy} (2026-08-04)
\item \texttt{https://docs.grove.finance/grove-app} (2026-08-04)
\item \texttt{https://docs.grove.finance/faqs} (2026-08-04)
\item \texttt{https://docs.grove.finance/grove-token} (2026-08-04)
\item \texttt{https://docs.grove.finance/grove-basin} (2026-08-04)
\end{itemize}

\subsection*{Kalshi}
\begin{itemize}\setlength{\itemsep}{0pt}\small
\item \texttt{https://www.cftc.gov/filings/orgrules/rules07012525155.pdf} (2026-08-04)
\item \texttt{https://docs.kalshi.com/getting\_started/orderbook\_responses} (2026-08-04)
\item \texttt{https://kalshi-public-docs.s3.amazonaws.com/regulatory/rulebook/Kalshi\%20DCO\%20Rulebook\%201.4.pdf} (2026-08-04)
\item \texttt{https://www.cftc.gov/filings/ptc/ptc11142532444.pdf} (2026-08-04)
\item \texttt{https://kalshi.com/docs/kalshi-fee-schedule.pdf} (2026-08-04)
\item \texttt{https://kalshi.com/perpetuals/learn} (2026-08-04)
\end{itemize}

\subsection*{Polymarket}
\begin{itemize}\setlength{\itemsep}{0pt}\small
\item \texttt{https://raw.githubusercontent.com/gnosis/conditional-tokens-contracts/master/contracts/ConditionalTokens.sol} (2026-08-04)
\item \texttt{https://github.com/Polymarket/ctf-exchange-v2} (2026-08-04)
\item \texttt{https://help.polymarket.com/en/articles/13364518-how-are-prediction-markets-resolved} (2026-08-04)
\item \texttt{https://blog.uma.xyz/articles/managed-proposers} (2026-08-04)
\item \texttt{https://raw.githubusercontent.com/Polymarket/uma-ctf-adapter/main/src/UmaCtfAdapter.sol} (2026-08-04)
\item \texttt{https://docs.polymarket.com/developers/CTF/overview} (2026-08-04)
\item \texttt{https://github.com/polymarket/neg-risk-ctf-adapter} (2026-08-04)
\item \texttt{https://help.polymarket.com/en/articles/14762452-polymarket-exchange-upgrade-april-28-2026} (2026-08-04)
\item \texttt{https://www.polymarketexchange.com/files/legal/latest/rulebook} (2026-08-04)
\item \texttt{https://www.polymarketexchange.com/clearing/} (2026-08-04)
\end{itemize}

\subsection*{Steakhouse Financial}
\begin{itemize}\setlength{\itemsep}{0pt}\small
\item \texttt{https://github.com/morpho-org/metamorpho} (2026-08-04)
\item \texttt{https://docs.morpho.org/curate/concepts/roles/} (2026-08-04)
\item \texttt{https://steakhouse.financial/docs/risk-management/monitoring/allocation-process.md} (2026-08-04)
\item \texttt{https://steakhouse.financial/docs/products/stablecoin-products/vaults.md} (2026-08-04)
\item \texttt{https://steakhouse.financial/docs/risk-management/monitoring/vault-setup-and-controls.md} (2026-08-04)
\item \texttt{https://docs.morpho.org/curate/concepts/fee/} (2026-08-04)
\item \texttt{https://steakhouse.financial/docs/documents/disclaimers/vaults.md} (2026-08-04)
\item \texttt{https://steakhouse.financial/docs/llms.txt} (2026-08-04)
\end{itemize}

\end{document}